\documentclass[11pt,letterpaper]{article}
\usepackage[T1]{fontenc}
\usepackage[utf8]{inputenc}
\usepackage{lmodern}
\usepackage[margin=0.88in,headheight=14pt]{geometry}
\usepackage{mathtools,amssymb,amsthm}
\usepackage{microtype}
\usepackage{xcolor}
\usepackage{booktabs,longtable,tabularx,array}
\usepackage{enumitem}
\usepackage{listings}
\usepackage{fancyhdr}
\usepackage{xurl}
\usepackage[unicode,hidelinks]{hyperref}
\definecolor{ink}{HTML}{15364B}
\definecolor{accent}{HTML}{216780}
\definecolor{softgray}{HTML}{F3F5F7}
\hypersetup{pdftitle={Constructive Number Theory in iPRA},
 pdfauthor={Research report prepared for Vladimir Reshetnikov},
 pdfsubject={Finite proofs, primitive recursive bounds, and Lean formalization blueprint},
 colorlinks=true,linkcolor=ink,urlcolor=accent,citecolor=ink}
\newcommand{\N}{\mathbb N}
\newcommand{\Z}{\mathbb Z}
\newcommand{\Q}{\mathbb Q}
\newcommand{\ipra}{\ensuremath{\mathsf{iPRA}}}
\newtheorem{theorem}{Theorem}[section]
\newtheorem{lemma}[theorem]{Lemma}
\newtheorem{corollary}[theorem]{Corollary}
\newtheorem{proposition}[theorem]{Proposition}
\theoremstyle{definition}
\newtheorem{definition}[theorem]{Definition}
\theoremstyle{remark}
\newtheorem*{remark}{Remark}
\setlist{topsep=5pt,itemsep=3pt,parsep=0pt}
\setlength{\parskip}{4pt}
\setlength{\parindent}{0pt}
\setlength{\emergencystretch}{3em}
\allowdisplaybreaks[2]
\lstset{basicstyle=\small\ttfamily,columns=fullflexible,
 breaklines=true,frame=single,rulecolor=\color{accent!30},
 backgroundcolor=\color{softgray},aboveskip=8pt,belowskip=8pt,
 keepspaces=true,showstringspaces=false}
\pagestyle{fancy}
\fancyhf{}
\fancyhead[L]{\small\color{ink}Constructive Number Theory in \ipra}
\fancyhead[R]{\small Research report}
\fancyfoot[C]{\small\thepage}
\renewcommand{\headrulewidth}{0.3pt}
\newcommand{\ResultCount}{106}
\newcommand{\FoundationCount}{5}
\newcommand{\NumberTheoryCount}{101}


\begin{document}
\hypersetup{pageanchor=false}
\begin{titlepage}
\color{ink}
\vspace*{0.8cm}
{\large MATHEMATICAL RESEARCH REPORT\par}
\vspace{0.35cm}
\hrule height 1.2pt
\vspace{1.1cm}
{\Huge\bfseries Constructive\par Number Theory\par}
\vspace{0.3cm}
{\LARGE in intuitionistic primitive\par recursive arithmetic\par}
\vspace{0.9cm}
{\large Fully developed finite proofs, explicit witnesses,\par
and a Lean formalization blueprint\par}
\vspace{1.2cm}
{\Large\bfseries \NumberTheoryCount\ number-theoretic results\par}
\vspace{0.15cm}
{\large plus \FoundationCount\ foundational finite-proof lemmas\par}
\vfill
{\large Prepared for Vladimir Reshetnikov\par}
\vspace{0.35cm}
{\large 22 November 2026\par}
\vspace{0.6cm}
{\small Source scope: \emph{Number theory} and relevant immediately linked topics.\par
ProveIt integration snapshot: \texttt{bbb85238dac5}.\par
Proof status: complete human arguments; uncompiled Lean blueprint;\par
finite exact-arithmetic checks recorded separately.\par}
\end{titlepage}
\hypersetup{pageanchor=true}
\tableofcontents
\clearpage
\section*{Scope, results, and reading guide}
\addcontentsline{toc}{section}{Scope, results, and reading guide}

The task is to prove as much as possible of the mathematics represented by the Wikipedia article \emph{Number theory} and its relevant immediate outgoing links, while remaining in intuitionistic primitive recursive arithmetic. The article is used as a topic index. The proofs below are independently developed finite arguments, not an appeal to Wikipedia as a proof authority. The selected coverage spans divisibility, prime decomposition, congruences, finite fields, quadratic reciprocity, additive representations, Diophantine equations and approximation, finite analytic estimates, elementary algebraic number theory, $p$-adic lifting, and partitions~\cite{wiki-number,mit781}.

The report contains \ResultCount\ numbered mathematical results, of which \FoundationCount\ are foundational finite-proof lemmas and \NumberTheoryCount\ are number-theoretic theorems, lemmas, or corollaries. These counts refer to distinct numbered entries, not to individual consequences packed into a theorem. They are a substantial selection, not a claim to exhaust the indefinitely extensible subject or to maximize a formal theorem count. Each numbered result has a human-readable proof. The archive adds Lean definitions and proof obligations, exact-arithmetic reference algorithms, finite validation records, a theorem inventory, and source provenance.

\begin{center}
\begin{tabularx}{\linewidth}{p{0.27\linewidth}X}
\toprule
Area & Principal results developed in full\\
\midrule
Arithmetic and primes & Division, Euclidean algorithm and B\'ezout, linear equations, Euclid's lemma, unique factorization, valuations, lcm, bounded prime infinitude, and sieve correctness.\\
Congruences and functions & Generalized finite CRT, Euler and Fermat, Wilson and converse, RSA correctness, Lucas and Kummer, divisor and totient formulas, M\"obius inversion, Liouville's divisor sum, and even perfect numbers.\\
Quadratic theory & Polynomial root bounds, primitive roots, Euler's criterion, Gauss's lemma, reciprocity and supplements, complete two-square classification, and the four-square theorem.\\
Diophantine theory & Rational roots, Pythagorean triples, square-area and quartic descents, continued fractions, Farey separation, finite Dirichlet approximation, bounded Pell existence, and classification by fundamental powers.\\
Analytic and probabilistic & Exact inclusion--exclusion sieve, dyadic Chebyshev bounds, explicit divergence of prime reciprocals, zeta tails and Euler-product error, prime-divisor means, and finite CRT independence.\\
Algebraic and combinatorial & A concrete nonunique factorization ring, truncated valuations, unique Hensel lifting, compatible $p$-adic names, partition recursion, odd/distinct partitions, and conjugation.\\
Transcendence & An explicit Liouville rational name, factorial-tail bounds, and uniform positive separation from every nonzero integer polynomial.\\
\bottomrule
\end{tabularx}
\end{center}

\subsection*{Three representative explicit bounds}
The main theme is that the existence statements carry actual finite search bounds. For every $n$, a four-square witness lies in $[0,n]^4$, and the full proof establishes that the $(n+1)^4$-candidate search succeeds. For a nonsquare $D\ge2$, put
\[
 C=4D+2,\quad B=4(D+1),\quad K=2C^3,\quad U=(B+1)^{K+1},
 \quad H=4(D+1)^2(U+1)^2.
\]
Then a nontrivial solution of $x^2-Dy^2=1$ exists with $2\le x\le H$ and $1\le y\le H$. The proof uses rational approximation and a finite collision of equal norms and residue classes. It does not assume periodic continued fractions.

For the prime reciprocal sum a similarly explicit statement is
\[
 \sum_{p\le 2^{\,2\cdot4^{4k+1}}}\frac1p>k\qquad(k\in\N).
\]
The proof uses only finite rational products and the dyadic harmonic lower bound. Such bounds are deliberately generous. Primitive recursive arithmetic constrains the form of the construction and its proof; it does not require a practical algorithm or an optimal asymptotic estimate.

\subsection*{Correspondence with the source scope}
The mathematical map below records the portion of the linked topic actually developed. Links are scope references; the later proofs supply the mathematics.
\begin{longtable}{p{0.31\linewidth}p{0.61\linewidth}}
\toprule
Root article or immediate topic & Treatment and exact scope\\
\midrule\endhead
\emph{Number theory}, elementary subsection & Integer arithmetic, congruences, primes, factorial and divisor identities, right triangles, and partitions. The \emph{Elementary number theory} URL currently redirects to this article.\\
\emph{Prime number} & Factorization, infinitude, two residue classes, primality by finite tests, dyadic distribution bounds, prime reciprocal divergence, and Euler products.\\
\emph{Finite field} & Prime fields $\mathbf F_p$, polynomial roots, the cyclic multiplicative group, and power equations. The classification of every finite field is outside this report.\\
\emph{Diophantine approximation} & Rational finite Dirichlet approximation, simultaneous approximation, continued fractions, explicit transfer with slack to fixed rational names, and Liouville's explicit transcendental example.\\
\emph{Algebraic number theory} & The pair-ring $\Z[\sqrt{-5}]$, norm, algebraic-integrality certificates, and an explicit failure of unique factorization. General class groups and class field theory are outside the proved scope.\\
\emph{$p$-adic number} & Finite precision valuations, Hensel's simple-root lifting, and a fixed primitive recursive compatible residue sequence. No unrestricted completion or quantification over all $p$-adic numbers is assumed.\\
\emph{Probabilistic number theory} & Exact finite means and divisibility independence over a complete CRT period. No central limit theorem is inferred.\\
\emph{Analytic number theory} and zeta discussion & Finite sieve identities, Chebyshev-scale estimates, and Euler products for integer exponents $s\ge2$ with a rational error bound. No complex analytic continuation is used.\\
\emph{Diophantine geometry} & Rational-root obstructions, the right-triangle parametrization, Pell curves, and explicit finite descent examples. No general finiteness theorem for higher-genus curves is claimed.\\
\bottomrule
\end{longtable}

\subsection*{Limits of the claims}
For discrete statements the report proves their usual mathematical content. The four-square theorem is proved for every natural number, and quadratic reciprocity is proved for all distinct odd primes. For continuous statements the proved content is the specified arithmetical version: a fixed rational name and a primitive recursive accuracy schedule, or a finite certificate with a stated error. These formulations are stated explicitly at the point of use.

The report proves the obstruction to three squares but not the converse characterization. It proves primes in the classes $1$ and $3$ modulo $4$, not the full arithmetic-progression theorem. It supplies the exponent-four case of Fermat's theorem, not a proof for every exponent. It does not claim proofs of the Riemann hypothesis, Goldbach's conjecture, general class field theory, or unrestricted Diophantine finiteness theorems. Failure to include a theorem here is not a proof that it cannot be formulated or proved in \ipra. In particular the repository has a stronger arithmetic prime-number endpoint, discussed separately in Section~\ref{sec:lean}; this report's own proof portfolio stops at the explicit estimates it develops in full.

\subsection*{How to read the package}
Read Section~\ref{sec:foundations} first for the translation discipline and the distinction between human proofs, reference computations, and object-level formal proofs. The elementary chapter supplies the common arithmetic. The quadratic and Diophantine chapters can then largely be read independently. The final formalization section gives exact repository interfaces, concrete Lean goal types, and a staged completion route. The machine-readable inventory and validation summaries make the coverage and evidence auditable without relying on the present synopsis.

\clearpage
\section{The arithmetic setting and the proof discipline}\label{sec:foundations}

\subsection{What is being proved}
This report gives first-order arithmetic versions of a broad collection of number-theoretic results. For an integer theorem such as the four-square theorem the version is the ordinary theorem, with a specified primitive recursive witness. For an analytic theorem the version is an explicitly quantified finite inequality or a statement about a specified primitive recursive sequence of rational approximations. These two situations must be distinguished. A first-order theory of numbers does not quantify over an unrestricted set of real numbers, all subsets of the integers, or all functions on the integers.

Our base theory \ipra\ has one sort, natural numbers including zero, intuitionistic first-order logic with equality, and a symbol for each primitive recursive definition. It includes the defining equations, successor injectivity, $Sx\ne0$, and induction for quantifier-free formulas with parameters. The defining equations for a primitive recursion are
\[
 f(0,\vec x)=g(\vec x),\qquad
 f(Sn,\vec x)=h(n,f(n,\vec x),\vec x).
\]
Induction is the axiom schema
\[
 A(0,\vec x)\ \wedge\ \forall n\bigl(A(n,\vec x)\to A(Sn,\vec x)\bigr)
 \ \longrightarrow\ \forall n A(n,\vec x),
\]
where $A$ is quantifier-free. Induction on a formula with an unbounded existential quantifier is not an available rule. The notation $\ipra\vdash\varphi$ means derivability in this object theory, not merely truth of a corresponding theorem in Lean's natural numbers.

The logical rules are the usual assumption and substitution rules; introduction and elimination for $\wedge,\vee,\to,\forall,\exists$; elimination of falsehood; equality reflexivity and replacement of equals. Universal introduction requires that the generalized variable is absent from the open assumptions. Existential elimination requires a fresh witness variable absent from the conclusion and other assumptions. Negation abbreviates implication to falsehood. There is no rule of unrestricted excluded middle, double-negation elimination, choice, or reflection from semantic truth to derivability. Arithmetic equality and bounded finite questions will nevertheless be decidable by explicit computations.

\subsection{Numerical encodings}
All encodings below have primitive recursive constructors, accessors, and validity tests. An integer is represented by its sign and absolute value, canonically by the natural code $2a$ for $a\ge0$ and $2a-1$ for $-a<0$. A rational is a signed numerator and a positive denominator; equality and order are tested by cross multiplication, so normalization is optional. A rational identity in a proof can always be cleared of its positive denominators and becomes an integer identity. An integer identity becomes an equality between natural-number expressions by moving negative summands to the other side.

For definiteness, pairs may be encoded by
\[
 \langle a,b\rangle=2^a(2b+1)-1.
\]
Decode $a$ as the largest $j\le c+1$ such that $2^j\mid c+1$, then decode $b$ by exact division. This is a bounded search. A finite list of length $l$ with entries at most $M$ can be represented by the triple $(l,M,A)$, where
\[
 A=\sum_{i<l} a_i(M+1)^i,\qquad 0\le A<(M+1)^l
\]
when $M\ge1$; the all-zero case uses base $2$. The $i$th digit is obtained by quotient and remainder. Pair the three numbers to obtain a single code. Lists whose sizes change can be re-encoded with a new maximum at each step. Products, sums, concatenation, bounded permutations, coefficient vectors, and finite tables are therefore primitive recursive. We never require a bound on running time that is polynomial or practically small.

The division functions used in the encodings can first be defined by an $a+1$-step scan of candidate quotients at input $(a,b)$, with a default output for $b=0$. Their defining equations do not presuppose Euclid's theorem; correctness for $b>0$ is proved in the elementary chapter. Similarly, an integer square root is defined by searching $j\le n$ for the largest $j$ with $j^2\le n$. The elementary order lemmas verify $r^2\le n<(r+1)^2$. Thus the encodings do not assume their later number-theoretic applications.

\begin{lemma}[Bounded Boolean reflection]\label{fd:reflection}
For each fixed formula $A(\vec x)$ made from primitive recursive atomic equalities by propositional connectives and bounded quantifiers, there is a primitive recursive Boolean $[A](\vec x)$ for which \ipra\ proves
\[
 [A]=0\ \vee\ [A]=1,
 \qquad A\ \longleftrightarrow\ [A]=1.
\]
Consequently bounded-formula induction is derivable, by applying quantifier-free induction to $[A]=1$.
\end{lemma}
\begin{proof}
Equality is decided by simultaneous predecessor comparison, or by testing whether both truncated differences vanish. The equations for this test and induction on its input prove equality or its negation. On Boolean values use multiplication for conjunction, $1-b$ for negation, and $\max$ for disjunction. Implication between decidable formulas is represented by $\max(1-a,b)$. These equivalences are intuitionistically valid because the induction hypotheses provide the relevant decisions.

For $\forall i<t\,B(i,\vec x)$ use the product $P(t)=\prod_{i<t}[B(i,\vec x)]$. Its equations are $P(0)=1$ and $P(t+1)=P(t)[B(t,\vec x)]$. Quantifier-free induction proves that it is a Boolean. For each fixed index $i$, induction proves $i<j\wedge P(j)=1\to[B(i,\vec x)]=1$. In the other direction, retain $\forall i<t\,B(i,\vec x)$ as an open assumption and induct on the quantifier-free assertion $j\le t\to P(j)=1$. The assumption supplies the next factor by universal elimination. Thus this argument does not assume bounded-formula induction in order to establish it.

For an existential quantifier, put $b_i=[B(i,\vec x)]$ and $F(j)=\prod_{i<j}(1-b_i)$. With the final bound $t$ fixed, initialize $w(0)=t$ and update $w(j)$ to $j$ exactly when $w(j)=t$ and $b_j=1$. Direct quantifier-free induction, guarded by $j\le t$, proves
\[
 F(j)=1\ \longrightarrow\ w(j)=t,\qquad
 F(j)=0\ \longrightarrow\ w(j)<j\ \wedge\ b_{w(j)}=1.
\]
For each fixed $i$, another quantifier-free induction proves $i<j\wedge b_i=1\to F(j)=0$. Hence $1-F(t)=1$ is equivalent to the bounded existential, with witness $w(t)$. These direct Boolean invariants establish the scan before using any bounded-formula induction. Nest the construction in the syntactic order of the bounded formula. This is an effective metatheoretic construction for each fixed formula, not a universal evaluator for arbitrary primitive recursive programs.
\end{proof}

\begin{lemma}[Bounded search, finite sums, and finite products]\label{fd:search}
Let $b(i,\vec x)$ be a fixed primitive recursive Boolean and let $T(\vec x)$ be primitive recursive. The first $i<T$ with $b(i,\vec x)=1$, with the failure value $T$, is primitive recursive. Its successful answer satisfies the predicate and is least; failure is equivalent to all $T$ tests being zero. Finite sums and products of fixed primitive recursive functions, including signed or rational ones, are primitive recursive.
\end{lemma}
\begin{proof}
At stage $j$ store the least successful index below $j$, using the sentinel $T$ until one is found. If the stored value is $T$ and $b(j)=1$, replace it by $j$; otherwise retain it. The invariant is: either the stored value is a successful least index below $j$, or it is $T$ and every previous test is zero. Every quantifier in this invariant is bounded by $j$ or $T$; Lemma~\ref{fd:reflection} replaces it with a quantifier-free Boolean equation. Start at $j=0$ and iterate exactly $T$ steps. The update equations and the invariant prove the claim. For sums store the previous sum and add the next term; for products multiply. Rational updates use pairs of integers and positive denominators. A common denominator such as the product of all input denominators provides an alternative integer-only implementation.
\end{proof}

\begin{lemma}[Clocked loops and simultaneous invariants]\label{fd:loops}
Suppose a state update $F(s,\vec x)$, initial state $s_0(\vec x)$, and iteration count $T(\vec x)$ are primitive recursive. If a bounded arithmetic invariant $I(j,s,\vec x)$ holds initially and is preserved by the update, \ipra\ proves it at every $j\le T$. The resulting state and every selected output component are primitive recursive.
\end{lemma}
\begin{proof}
Define $s(0)=s_0$ and $s(j+1)=F(s(j),\vec x)$. Substitute this term in the Boolean characteristic of $I$. Quantifier-free induction proves $[I(j,s(j),\vec x)]=1$. If the desired invariant only concerns the first $T$ stages, freeze the state thereafter, or use the implication $j\le T\to I$. Simultaneous recursions are one recursion on a tuple code. Finite strong induction is handled by storing the previously constructed witnesses and the product of their correctness bits. At stage $j$ every preceding case is read from that table; the induction formula is still a Boolean equation. This justifies the finite strong-induction presentations used below.
\end{proof}

\begin{lemma}[Finite descent]\label{fd:descent}
Let $\operatorname{Bad}(s)$ and a height $h(s)$ be primitive recursive. Suppose \ipra\ proves that a bad state has positive height and that a primitive recursive update $F$ maps every bad state to a bad state of strictly smaller height. Then \ipra\ proves that no bad state exists.
\end{lemma}
\begin{proof}
Starting from a proposed bad state, run $F$ for $h(s)$ steps. At stage $j\le h(s)$ maintain
\[
 \operatorname{Bad}(s_j),\qquad j+h(s_j)\le h(s).
\]
Strict decrease makes the second inequality survive the next step. At $j=h(s)$ it forces $h(s_j)=0$, contradicting positivity. Both assertions are bounded arithmetic, and the iteration has a supplied finite clock. This argument does not use an unbounded least-counterexample principle. In a witness-producing descent, freeze the state once the target is reached and apply the same argument to show that failure to reach it within the clock is impossible.
\end{proof}

\begin{lemma}[Finite pigeonhole principle with witnesses]\label{fd:php}
For a coded list $a_0,\ldots,a_M$ with $a_i<M$, one can find $i<j\le M$ with $a_i=a_j$ by a primitive recursive search, and \ipra\ proves its success.
\end{lemma}
\begin{proof}
Test all at most $(M+1)^2$ pairs. If none succeeds, for each $v<M$ the count $c_v(t)=\#\{i<t:a_i=v\}$ is at most one. Induction on $t$ proves $\sum_{v<M}c_v(t)=t$, because inserting $a_t$ increases exactly its one count. At $t=M+1$, the left side is at most $M$ and the right side is $M+1$, a contradiction. The counts and both nested sums are bounded loops. The finite decision of whether a collision exists, together with the contradiction in its failure branch, produces the successful pair.
\end{proof}

\subsection{Witnesses, schemas, and analytic names}
Whenever a theorem says ``choose a prime divisor'', ``take an inverse'', ``take a square root'', or ``select a collision'', the indicated finite bounds determine a search using Lemma~\ref{fd:search}. Witnesses returned by separate branches can be combined by a primitive recursive conditional. The implication from hypotheses to witness correctness is proved by the stated invariants; a default value is used when the hypotheses fail. Accordingly the functions are total even when a mathematical operation is described only on its usual domain.

Statements about an arbitrary arithmetic function $f$ are schemas: fix a primitive recursive definition of $f$, then make the displayed construction and derivation for that definition. A uniform theorem can instead take a finite table as numerical input. It is incorrect to use a single primitive recursive function $\operatorname{eval}(e,x)$ evaluating every primitive recursive program with code $e$: diagonalizing $\operatorname{eval}(e,e)+1$ rules out such an evaluator. Uniform evaluation of a finite polynomial or a fixed finite table is harmless because those are restricted, explicitly bounded data formats.

A real name, when needed, is a fixed primitive recursive rational sequence $r(k)$ accompanied by a derived modulus, for example
\[
 \forall l\ge k\quad |r(l)-r(k)|\le 2^{-k}.
\]
The displayed formula concerns natural-number codes and is interpreted by clearing denominators. The external real number represented by the name is explanatory notation. Within \ipra\ the theorem is the family of rational estimates about $r$. No step asks whether an arbitrary represented real is exactly zero, an integer, positive, or rational. A strict comparison needs a finite separation margin; the approximation chapter makes that margin explicit.

\subsection{Proof status and the ProveIt integration target}
The human-readable proofs in this report are complete mathematical arguments with an explicit reduction to the preceding finite proof rules. The supplied Lean files are definitions and a formalization blueprint; they have not been compiled in this environment, which has no Lean executable. They are not presented as new kernel-checked derivations. Finite Python checks exercise identities and witness implementations; they do not establish the universally quantified theorems.

The repository snapshot used for integration is
\begin{center}\small\ttfamily bbb85238dac533fc4b7b84a750cca864fe35fed5.\end{center}
Its top-level arithmetic overview distinguishes the checked libraries from delivered research packets, which it does not certify merely by filing them~\cite{proveit-readme}. We inspected the definitions of \texttt{PR}, \texttt{PTm}, \texttt{PRAAx}, and \texttt{PRA.qfInd}, and selected source files for search, division, B\'ezout, bounded sums, pigeonhole, descent, reciprocity, and the arithmetic prime number theorem. These are useful existing targets, not an automatic certification of this report. The formalization section records exact source paths and the obligations needed to turn this package into derivations of type
\[
 \texttt{([] : List (Fm PTm 0))}\ \vdash[\texttt{PRAAx}]\ \varphi.
\]

\clearpage
\section{Divisibility, primes, and finite arithmetic}
\label{sec:elementary}

This chapter develops the integer arithmetic needed by the later chapters.
Every search is bounded.  A statement about an arbitrary primitive recursive
function is a schema with that function fixed, or, equivalently for the
finite assertions used here, a statement about a coded finite table.  It
never means that arithmetic has a total interpreter for all primitive
recursive programs.  The conventional topics correspond to the elementary
portion of the source article and to the university course outline
\cite{mit781}; the arguments and arithmetic bounds below are supplied in full.

We use natural numbers including zero.  A signed integer is represented by
its sign and absolute value; addition, multiplication, comparison, and
absolute value have their usual primitive recursive definitions on these
codes.  For natural numbers, $d\mid n$ means that $n=dq$ for some $q\le n$,
with the separate convention $0\mid0$.  For integers one uses a signed
quotient of absolute value at most $|n|$ when $d\ne0$.  Thus the divisibility
predicate is decidable by a bounded test.  A bounded maximum, minimum,
conjunction, sum, or product is implemented by a fold over its displayed
finite range.  For example, primality is the bounded predicate
\[
 \operatorname{Prime}(p)\ \Longleftrightarrow\
 p\ge2\ \text{ and }\ (\forall d\le p)
       (d\mid p\Longrightarrow d=1\text{ or }d=p).
\]
A bounded universal quantifier in a loop invariant abbreviates a Boolean
fold.  Induction on that Boolean assertion is quantifier-free induction.

\subsection{Division and the Euclidean algorithm}

\begin{theorem}[Division with a bounded algorithm]\label{el:division}
For $a\in\N$ and $b>0$ there are unique natural numbers $q,r$ such that
$a=bq+r$ and $r<b$.  They satisfy $q\le a$ and $r\le a$, and are primitive
recursive functions of $a,b$.
\end{theorem}
\begin{proof}
Start with $(q_0,r_0)=(0,0)$.  To pass from $i$ to $i+1$, replace
$(q_i,r_i)$ by $(q_i,r_i+1)$ if $r_i+1<b$, and by $(q_i+1,0)$ otherwise.
The invariant is
\[
 i=bq_i+r_i,\qquad r_i<b,\qquad q_i\le i,\quad r_i\le i.
\]
In the second branch $r_i<b$ and the failure of $r_i+1<b$ give
$r_i+1=b$.  This proves preservation by direct arithmetic.  After exactly
$a$ steps the required pair has been obtained.  For uniqueness, if
$bq+r=bq'+r'$ and $q<q'$, then
$bq+r<b(q+1)\le bq'\le bq'+r'$, a contradiction.  Decidable trichotomy
therefore gives $q=q'$ and then $r=r'$.  No least-number principle is used.
\end{proof}

We write these outputs as $\lfloor a/b\rfloor$ and $a\bmod b$.
For a signed integer $z$ the same construction on $|z|$, followed when
necessary by $(-q-1,b-r)$, gives a unique residue in $[0,b)$.

\begin{theorem}[Euclidean algorithm and B\'ezout identity]
\label{el:bezout}
There are primitive recursive functions $g(a,b)\in\N$ and
$u(a,b),v(a,b)\in\Z$ satisfying
\[
 g\mid a,\quad g\mid b,\quad ua+vb=g,
 \qquad d\mid a\text{ and }d\mid b\Longrightarrow d\mid g.
\]
Here $g(0,0)=0$; otherwise $g>0$.  Consequently $g$ is the greatest
common divisor.  A computation with at most $a+b+1$ Euclidean steps
suffices, and the displayed coefficients may be chosen with
\[
 |u|,|v|\le (a+b+2)^{a+b+1}.
\]
\end{theorem}
\begin{proof}
Maintain two remainders $x,y$ and two coefficient rows $U,V$, initially
$x=a$, $y=b$, $U=(1,0)$, $V=(0,1)$.  If $y=0$, keep the state fixed.
Otherwise put $q=\lfloor x/y\rfloor$ and replace
\[
 (x,y,U,V)\quad\text{by}\quad(y,x-qy,V,U-qV).
\]
The two remainders remain the linear combinations described by the rows.
For any fixed $d$, the conditions $d\mid x,d\mid y$ are equivalent to
$d\mid y,d\mid x-qy$: one direction is subtraction and the other is
$x=(x-qy)+qy$.  In particular the initial common divisors and the current
common divisors coincide.

Whenever a nonzero second remainder is replaced, its successor is a
smaller natural number.  After the first step both remainders are at most
$\max(a,b)$; the second remainder starts at $b$ and decreases by at least
one on every active step.  A loop with $a+b+1$ steps has therefore reached
$y=0$.  A quantifier-free termination invariant is: either $y_i=0$, or
$y_i+i\le b$ .
At the terminal state let $g=x$ and let $(u,v)=U$.  The common-divisor
invariant gives all the claimed divisibilities, including $g\mid a,b$.
If the initial pair is nonzero then an active step never makes both
remainders zero, so $g>0$.  A positive common divisor of $a,b$ divides
the positive number $g$ and is therefore at most $g$.

Every quotient is at most $a+b$.  If all four coefficient magnitudes are
at most $C$, a step makes them at most $(a+b+1)C$.  Starting with $C=1$
and using the larger base $a+b+2$ gives the asserted bound.  All invariants
are equalities, inequalities, and bounded divisibility tests.  One may
preserve the common-divisor assertion for each fixed $d$ separately;
there is no need to induct on an unbounded universal formula.
\end{proof}

\begin{theorem}[Linear Diophantine equations]\label{el:linear}
Let $a,b\in\N$, not both zero, and $c\in\Z$, and put $g=g(a,b)$.
The equation $ax+by=c$ has integer solutions exactly when $g\mid c$.
When it does, one solution is
\[
 x_0=u(a,b)(c/g),\qquad y_0=v(a,b)(c/g).
\]
All solutions are
\[
 x=x_0+(b/g)t,\qquad y=y_0-(a/g)t\qquad(t\in\Z).
\]
For $a=b=0$ solvability is exactly $c=0$.
\end{theorem}
\begin{proof}
A common divisor divides every linear combination, proving necessity.
Multiplying the B\'ezout identity by $c/g$ proves sufficiency.  Divide
$a,b$ by $g$ and write $a',b'$; their gcd is one, since a common divisor
$h>1$ would make $gh$ a common divisor of $a,b$ greater than $g$.
If $(x,y)$ is a solution, then
$a'(x-x_0)=-b'(y-y_0)$.  For coprime integers $A,B$, a B\'ezout
identity $rA+sB=1$ shows that $B\mid Az$ implies $B\mid z$:
multiply the identity by $z$.  If $b'>0$, this proves
$x-x_0=b't$, and cancellation gives $y-y_0=-a't$.
If $b'=0$, then $a'=1$, so $x=x_0$ and $y$ is arbitrary; the same
formula still describes every solution.  The constructed solution has
coordinate magnitudes at most
$(a+b+2)^{a+b+1}|c|$.  The assertion about an arbitrary proposed solution
uses its own coordinates as parameters, not an unbounded search for $t$.
\end{proof}

\subsection{Prime decomposition and its consequences}

\begin{lemma}[A least divisor is prime]\label{el:leastprime}
Every $n\ge2$ has a prime divisor at most $n$, obtained by taking the
least $d\in[2,n]$ dividing $n$.  If $n$ is composite, it has a prime
divisor $p$ with $p^2\le n$.
\end{lemma}
\begin{proof}
The search succeeds because $n\mid n$.  If its output $d$ were
composite, it would have a divisor $e$ with $1<e<d$; then $e\mid n$,
contrary to the checked minimality.  If $n$ is composite and $p$ is its
least prime divisor, write $n=pk$ with $k>1$.  A prime divisor of $k$
is at least $p$, so $k\ge p$ and $p^2\le n$.
All alternatives are decidable within $[0,n]$.
\end{proof}

\begin{theorem}[Euclid's lemma]\label{el:euclidlemma}
For a prime $p$, if $p\mid ab$, then $p\mid a$ or $p\mid b$.
Consequently, if $\gcd(a,m)=1$ and $m\mid ab$, then $m\mid b$.
\end{theorem}
\begin{proof}
Test $p\mid a$.  If it fails, the positive divisor $g(p,a)$ of $p$
cannot equal $p$, so it equals one.  A B\'ezout identity $up+va=1$,
multiplied by $b$, expresses $b$ as a sum of two multiples of $p$.
The coprime version follows by precisely the same multiplication.
This proof uses excluded middle only for the bounded predicate
$p\mid a$, whose characteristic function is already primitive recursive.
\end{proof}

\begin{theorem}[Fundamental theorem of arithmetic]
\label{el:fta}
Every $n\ge1$ has a primitive recursively computed nondecreasing list
of primes $p_0,\ldots,p_{\ell-1}$, with $\ell\le n$, whose product is $n$.
Any other nondecreasing prime list with product $n$ is identical to it.
Equivalently, the multiset of prime factors is unique.
\end{theorem}
\begin{proof}
Keep a remainder $r$, initially $n$, and an initially empty factor list.
If $r>1$, append its least prime divisor and divide $r$ by that divisor;
if $r=1$, do nothing.  An active step strictly decreases $r$, and the
invariant is
\[
 n=r\prod_{j<\ell}p_j,\qquad r\ge1,\qquad
 \ell\le i,\qquad (r>1\Longrightarrow r+i\le n).
\]
A loop of $n$ steps therefore ends with $r=1$.  A prime factor of the
new remainder was also a divisor of the old remainder, so it is at least
the prime just removed; the list is nondecreasing.  Every list entry is
at most $n$, and a length-$n$ array in base $n+1$, together with its
length, is a primitive recursive code.

For uniqueness, the first prime of one nonempty list divides the
product of the other.  Repeated use of Euclid's lemma, over at most the
length of that list, gives an entry divisible by this prime.  Since that
entry is prime, the entries are equal.  Interchanging the lists shows
that their first entries are equal: each is at most an entry equal to
the first entry of the other list.  Cancel that positive common entry
and repeat.  A prime product is at least $2^{\ell}$, so a list with
product one must be empty.  Thus neither list can end earlier than the
other.  This is an induction on a bounded comparison loop, with the
invariant that the canceled prefixes agree and the remaining products
agree.  It requires no induction on an existential factorization claim.
\end{proof}

For $p\ge2$ and $n\ge1$ define
\[
 v_p(n)=\max\{e\le n:p^e\mid n\}.
\]
The bound is valid because $p^e\ge e+1$ by induction.  This definition is
used only with $p$ prime when the valuation laws are invoked.

\begin{theorem}[Valuations, gcd, and lcm]\label{el:valuation}
For primes $p$ and positive $a,b$,
\[
 v_p(ab)=v_p(a)+v_p(b),\qquad
 a\mid b\ \Longleftrightarrow\
       (\forall p\le\max(a,b))
       (\operatorname{Prime}(p)\Longrightarrow v_p(a)\le v_p(b)).
\]
The valuations of $\gcd(a,b)$ and $\operatorname{lcm}(a,b)$ are the
respective minima and maxima.  In particular
\[
 \gcd(a,b)\operatorname{lcm}(a,b)=ab.
\]
\end{theorem}
\begin{proof}
Collect the occurrences of $p$ in the unique prime lists.  Concatenating
and sorting the lists for $a$ and $b$ gives a prime list for $ab$, so
uniqueness says that the occurrences add.  If $a\mid b$, write $b=ac$
and use additivity.  Conversely, subtract the exponent of each prime
in $a$ from that in $b$ and multiply the resulting nonnegative prime
powers.  The product $c$ satisfies $ac=b$, proving divisibility.
Take the products with minimum and maximum exponents to construct $G,L$.
The exponent criterion proves that $G$ divides both inputs and every
common divisor divides $G$; similarly $L$ is a common multiple dividing
every common multiple.  Thus $G$ is the Euclidean gcd and $L$ is the
least positive common multiple.  Since
$\min(r,s)+\max(r,s)=r+s$, their product is $ab$.
All prime scans can be limited to $\max(a,b)$; any proposed common
multiple is handled with its own finite bound.  The output $L\le ab$.
\end{proof}

Integer fractions use the primitive recursive canonical reduction proved in
Theorem~\ref{dio:normal}.  The normalization map sends $(a,b)$, $b>0$,
to $(a/g,b/g)$ with $g=\gcd(|a|,b)$; all fraction calculations in this
chapter may use these canonical finite codes.

\begin{theorem}[Euclid's theorem with a bound]\label{el:euclid}
For every $N\in\N$ a prime $p$ with
\[
             N<p\le N!+1
\]
can be computed primitive recursively.  Consequently there is no finite
list containing all primes.
\end{theorem}
\begin{proof}
Take the least prime divisor of $N!+1\ge2$.  If this prime were at most
$N$, it would divide both $N!$ and $N!+1$, and hence divide one, a
contradiction.  For a supplied finite prime list, use $N$ equal to its
maximum (zero for the empty list).  The factorial is defined by primitive
recursion, so this is an explicit uniform witness rather than an
unbounded search for the next prime.
\end{proof}

\begin{theorem}[Correctness of the sieve of Eratosthenes]
\label{el:sieve}
There is a primitive recursive Boolean array marking exactly the primes
at most $N$.  The usual sieve, which crosses off multiples of $d$ from
$d^2$ onward whenever $d$ is still unmarked, is correct.
\end{theorem}
\begin{proof}
Initially mark all entries $2,\ldots,N$ as candidates.  Scan
$d=2,\ldots,\lfloor\sqrt N\rfloor$, where the square root is the bounded
maximum of $r\le N$ with $r^2\le N$.  If $d$ is a candidate, cross off
$d^2+kd\le N$ for $k\le N$.
Inductively a candidate $d$ must be prime: if it were composite, its
least prime divisor $q<d$ would have crossed it off, because $d\ge q^2$.
No prime is ever crossed off, since every crossed-off integer is a
product of two integers at least two.  Finally, every composite $n\le N$
has a least prime divisor $p$ with $p^2\le n$; the step $d=p$ occurs,
$p$ is still a candidate, and $n$ is crossed off.  The array can be coded
by a number below $2^{N+1}$.  At most $N^2$ crossing-off checks and $N$
outer steps are enough.  The invariant ``a crossed-off entry has a known
proper factor, and every processed prime has crossed off its prescribed
multiples'' is a bounded Boolean assertion on that array.
\end{proof}

\subsection{Congruences, finite residue systems, and powers}

For $m>0$ put $a\equiv b\pmod m$ when $a\bmod m=b\bmod m$.
Division identifies this with $m\mid(a-b)$.  Substitution in differences
proves reflexivity, symmetry, transitivity, and compatibility with
addition and multiplication.  B\'ezout's identity gives the useful rule
\[
 \gcd(a,m)=1,\quad ax\equiv ay\pmod m
       \quad\Longrightarrow\quad x\equiv y\pmod m.
\]
All residue representatives below lie in a stated finite interval.

\begin{theorem}[Solving and counting linear congruences]
\label{el:linearcong}
Let $m>0$, $a\in\N$, $b\in\Z$, and $g=\gcd(a,m)$.
The congruence $ax\equiv b\pmod m$ is solvable exactly when $g\mid b$.
When solvable, it has exactly $g$ representatives in $[0,m)$, namely
\[
 x_0+k(m/g)\qquad(0\le k<g),
\]
where $0\le x_0<m/g$ is computed by a B\'ezout inverse modulo $m/g$.
\end{theorem}
\begin{proof}
A solution gives $ax-my=b$, proving necessity.  If $g\mid b$, divide
by $g$ to obtain $a'x\equiv b'\pmod {m'}$ with
$\gcd(a',m')=1$.  Compute $u$ with $ua'+vm'=1$ and take
$x_0=(ub')\bmod m'$.  Multiplication by $u$ shows that every solution
is congruent to $x_0$ modulo $m'$.  The integers in $[0,m)$ with that
residue are exactly $x_0+km'$ for $k<g$, by division of a proposed
representative by $m'$.  They are distinct because their differences
have magnitude less than $m$ and a nonzero difference of their indices
cannot be a multiple of $g$.  This includes $m'=1$ and $a=0$.
\end{proof}

\begin{theorem}[Generalized Chinese remainder theorem for two moduli]
\label{el:crt2}
For $m,n>0$, the system
$x\equiv a\pmod m$, $x\equiv b\pmod n$ has a solution exactly when
$\gcd(m,n)\mid(b-a)$.  There is then a unique solution in
$[0,L)$, where $L=\operatorname{lcm}(m,n)\le mn$.
\end{theorem}
\begin{proof}
Necessity follows by subtracting the congruences.  Put $g=\gcd(m,n)$,
$m'=m/g$, $n'=n/g$.  If $g\mid b-a$, solve
$m't\equiv(b-a)/g\pmod {n'}$ using the preceding theorem.  The integer
$a+mt$ satisfies both congruences.  Reduce it modulo $L=mn'$ to obtain
the required representative; both moduli divide $L$, so reduction
preserves the congruences.  Two solutions differ by a multiple of both
$m$ and $n$, and hence by a multiple of $L$ by the valuation criterion.
Two representatives in $[0,L)$ with that property are equal.
\end{proof}

\begin{theorem}[Chinese remainder theorem for a finite system]
\label{el:crt}
For positive moduli $m_0,\ldots,m_{k-1}$ and integer residues $a_i$,
a simultaneous solution exists exactly when
\[
 \gcd(m_i,m_j)\mid(a_i-a_j)\qquad(i,j<k).
\]
When it exists, there is a unique representative below
$L=\operatorname{lcm}(m_0,\ldots,m_{k-1})$, and it is primitive
recursively computable.  In particular pairwise coprime moduli admit
arbitrary residues, with $L=\prod_i m_i$.
\end{theorem}
\begin{proof}
First prove the pairwise coprime case by successively applying
Theorem~\ref{el:crt2}.  The product of the previously used moduli is
coprime to the next: otherwise a prime dividing their gcd would divide
one previous modulus, contradicting the hypothesis.  Thus each step
is possible and gives one representative below the enlarged product.

For general moduli take $M=\prod_i m_i$ (the empty product is one).
For each prime $p\le M$ occurring in some modulus, set
$e_p=\max_{i<k}v_p(m_i)$ and select the least index $j_p$ attaining it.
The compatibility hypotheses imply
\[
 a_{j_p}\equiv a_i\pmod {p^{v_p(m_i)}}
\]
for every $i$: the indicated prime power divides
$\gcd(m_i,m_{j_p})$.  Use the coprime case to solve
$x\equiv a_{j_p}\pmod {p^{e_p}}$ for the finitely many occurring primes.
Its modulus is $L=\prod_p p^{e_p}\le M$.  For each $i$, the difference
$x-a_i$ is divisible by every prime power in $m_i$, hence by $m_i$.
Necessity follows by subtraction, and uniqueness follows because a
difference divisible by all moduli is divisible by their lcm.
All maxima, index choices, and prime scans are bounded by $k$ and $M$.
For $k=0$ the assertion has the unique representative zero modulo one.
\end{proof}

Define
\[
 \varphi(n)=\#\{1\le r\le n:\gcd(r,n)=1\}\qquad(n\ge1).
\]
For $n>1$ the same count uses the residues $1,\ldots,n-1$.
The inverse of a coprime residue is computed by B\'ezout or by a scan
of $[0,n)$, so inverse and multiplication are primitive recursive maps
on an explicitly finite list.

\begin{theorem}[Euler's theorem]\label{el:euler}
If $n\ge1$ and $\gcd(|a|,n)=1$, then
\[
                    a^{\varphi(n)}\equiv1\pmod n.
\]
\end{theorem}
\begin{proof}
For $n=1$ every pair of integers is congruent.  For $n>1$, enumerate
the coprime residues as $r_0<\cdots<r_{h-1}$, where $h=\varphi(n)$.
Multiplication by $a$, followed by reduction, sends this list to itself:
a common divisor of $ar_i$ and $n$ would, using an inverse for $a$,
divide $r_i$.  An inverse for $a$ also gives the inverse map, so this
is a permutation of the finite list.  Products of a finite list are
unchanged by such a permutation.  Concretely, scan the image list,
find each next entry in the remaining source list, and swap it into
place; every search has length at most $h$ and each swap preserves
the product.  Hence
\[
 a^h\prod_{i<h}r_i\equiv\prod_{i<h}r_i\pmod n.
\]
The product is coprime to $n$: multiply the inverses of the entries.
Coprime cancellation proves the assertion.  The finite permutation
argument needs only bounded arrays and quantifier-free induction.
\end{proof}

\begin{corollary}[Fermat's little theorem]\label{el:fermat}
For a prime $p$, $p\nmid a$ implies $a^{p-1}\equiv1\pmod p$.
For every integer $a$ one has $a^p\equiv a\pmod p$.
\end{corollary}
\begin{proof}
The nonzero residues modulo a prime are exactly its coprime residues,
so $\varphi(p)=p-1$.  Apply Euler's theorem and multiply by $a$.
If $p\mid a$, the second congruence holds directly.  The case split
is the bounded divisibility test.
\end{proof}

\begin{theorem}[Order of a residue]\label{el:order}
For $n>1$ and $\gcd(|a|,n)=1$, the least positive exponent
$d\le\varphi(n)$ with $a^d\equiv1\pmod n$ exists and is primitive
recursive.  For every $t\in\N$,
\[
 a^t\equiv1\pmod n\ \Longleftrightarrow\ d\mid t,
 \qquad\text{and consequently }d\mid\varphi(n).
\]
\end{theorem}
\begin{proof}
Euler's theorem supplies a witness within the search range.  Write
$t=qd+r$, $0\le r<d$.  Since $a^d\equiv1$, one has
$a^t\equiv a^r$.  If this is one and $r>0$, the search defining $d$
would have found $r$ earlier.  Thus $r=0$.  Conversely, a multiple
of $d$ gives a power of $a^d$ and hence residue one.  Substitute
$t=\varphi(n)$ for the last assertion.
\end{proof}

\begin{theorem}[Wilson's theorem and converse]\label{el:wilson}
For every $n\ge2$,
\[
 \operatorname{Prime}(n)\quad\Longleftrightarrow\quad
 (n-1)!\equiv-1\pmod n.
\]
\end{theorem}
\begin{proof}
Let $p$ be prime.  Every residue $1,\ldots,p-1$ has a unique inverse.
An inverse equals the original residue exactly when
$p\mid(x-1)(x+1)$, which, by Euclid's lemma, gives $x=1$ or $x=p-1$
within the specified interval.  For odd $p$ these are distinct and
all other entries form pairs with product one.  To make the pairing
literal, maintain a Boolean ``used'' array; take the least unused
non-fixed entry, pair it with its computed inverse, and mark both.
The inverse is unused, since if it had already appeared then its
partner would have been the present entry.  After at most $p$ steps
all non-fixed entries have been paired.  The factorial residue is
therefore $1\cdot(p-1)=-1$.  For $p=2$, directly $1!=-1\pmod2$.

Conversely, if $n$ is composite choose a proper divisor $d$ with
$1<d<n$.  It divides $(n-1)!$.  If $n$ divided $(n-1)!+1$, then $d$
would also divide that number and hence divide one, a contradiction.
Because primality is decidable, this exclusion of the composite branch
proves primality constructively.
\end{proof}

\begin{theorem}[RSA correctness, including nonunits]\label{el:rsa}
Let $p,q$ be distinct primes, $N=pq$, and let positive $e,d$ satisfy
$ed\equiv1\pmod{(p-1)(q-1)}$.  For every integer message $m$,
\[
                    (m^e)^d\equiv m\pmod N.
\]
All key arithmetic and message transformations here are primitive
recursive.  This assertion makes no claim about cryptographic security.
\end{theorem}
\begin{proof}
Since $ed\ge1$, write $ed=1+t(p-1)(q-1)$ with $t\ge0$.
Modulo $p$, test whether $p\mid m$.  If so, both $m^{ed}$ and $m$
are zero.  Otherwise Fermat's theorem gives
$m^{ed}=m(m^{p-1})^{t(q-1)}\equiv m$.  The same proof works modulo
$q$.  These distinct prime moduli are coprime, so the difference is
divisible by their product.  The exponentiation identity
$(m^e)^d=m^{ed}$ is a direct induction.  To compute encryption or
decryption one can reduce after every multiplication; a loop of
length equal to the exponent keeps all residues below $N$.
\end{proof}

\begin{theorem}[Base-expansion divisibility rules]\label{el:digits}
If $b\ge2$ and $n=\sum_{i<k}d_i b^i$, $0\le d_i<b$, then
\[
 n\equiv\sum_{i<k}d_i\pmod{b-1},\qquad
 n\equiv\sum_{i<k}(-1)^i d_i\pmod{b+1}.
\]
If $m\mid b^j$, divisibility by $m$ depends only on the last $j$
base-$b$ digits.  Thus the usual decimal rules for $3,9,11$, and
powers dividing $10^j$ follow.
\end{theorem}
\begin{proof}
Modulo $b-1$ every power of $b$ is one; modulo $b+1$ its $i$th power
is $(-1)^i$.  These claims follow by induction on $i$ and can be
substituted into the finite sum.  For the last assertion, split the
sum at $j$; every term of the remaining part contains $b^j$.
The digit list itself is obtained by repeated division by $b$ for at
most $n+1$ steps, padding by zeros after the quotient reaches zero.
The invariant says that the original number is the extracted low
part plus $b^i$ times the current quotient.
\end{proof}

\subsection{Factorials, binomial coefficients, and their valuations}

Define $\binom nk$ by Pascal recursion, setting it to zero when $k>n$.
Finite summation of Pascal's rule proves
$(x+y)^n=\sum_{k=0}^n\binom nkx^ky^{n-k}$ by induction on $n$:
multiply the previous sum by $x+y$, shift the second index, and add
adjacent coefficients.  The same recursion proves
$k!(n-k)!\binom nk=n!$ for $k\le n$.  Indeed, at an interior
index $1\le k\le n$, multiplying Pascal's identity for
$\binom{n+1}{k}$ by $k!(n+1-k)!$ gives the sum
$(n+1-k)n!+kn!=(n+1)!$ by the two induction identities.
The boundary indices are immediate.

\begin{theorem}[Prime binomial divisibility and the Frobenius identity]
\label{el:binomial}
For a prime $p$ and $0<k<p$ one has $p\mid\binom pk$.
Consequently $(x+y)^p\equiv x^p+y^p\pmod p$ for integer $x,y$,
and the coefficientwise identity
$(A+B)^p=A^p+B^p$ holds for finite polynomials modulo $p$.
\end{theorem}
\begin{proof}
The factorial product $k!(p-k)!$ contains no factor divisible by $p$,
by repeated Euclid's lemma; it is coprime to $p$.  The factorial
identity shows that $p$ divides this product times $\binom pk$,
because $p\mid p!$.  Cancel the coprime product.  All the inner
coefficients of the binomial expansion therefore vanish modulo $p$.
The same calculation works in the finite polynomial ring: addition
and multiplication are bounded coefficient sums, and the binomial
expansion uses only commutativity and distributivity.  There is no
appeal to an infinite polynomial object.
\end{proof}

\begin{theorem}[Legendre's factorial valuation formula]
\label{el:legendre}
For a prime $p$ and $n\ge0$,
\[
        v_p(n!)=\sum_{j=1}^{n}\left\lfloor\frac{n}{p^j}\right\rfloor.
\]
If $n=\sum_{i<k}d_i p^i$ is its base-$p$ expansion and
$s_p(n)=\sum_{i<k}d_i$, this is equivalently
\[
                 (p-1)v_p(n!)=n-s_p(n).
\]
\end{theorem}
\begin{proof}
Valuation additivity gives $v_p(n!)=\sum_{a=1}^n v_p(a)$.
For each $a\le n$, the exponent $v_p(a)$ counts exactly the indices
$j\in[1,n]$ with $p^j\mid a$: the valid powers form the initial segment
through $v_p(a)$, and $p^n>n$ for $n>0$.  Interchange the two finite
sums.  The number of multiples of $p^j$ in $[1,n]$ is
$\lfloor n/p^j\rfloor$, by division, proving the first formula.
For the second, division of the base-$p$ expansion gives
$\lfloor n/p^j\rfloor=\sum_{i\ge j}d_i p^{i-j}$.
Summing and using
$(p-1)\sum_{h=0}^{i-1}p^h=p^i-1$ gives
$(p-1)v_p(n!)=\sum_i d_i(p^i-1)=n-s_p(n)$.
For $n=0$ all sums are empty and $0!=1$ has valuation zero.
\end{proof}

\begin{theorem}[Lucas's binomial congruence]\label{el:lucas}
Let $p$ be prime and write $n=\sum_i n_i p^i$ and
$k=\sum_i k_i p^i$ with digits in $[0,p)$, padding with zeros to the
same finite length.  Then
\[
                    \binom nk\equiv\prod_i\binom{n_i}{k_i}\pmod p,
\]
where a binomial coefficient with lower index exceeding the upper
is zero.
\end{theorem}
\begin{proof}
Iterate the finite polynomial identity of
Theorem~\ref{el:binomial} to obtain
$(1+X)^{p^i}\equiv1+X^{p^i}$ coefficientwise modulo $p$.
Exponent addition therefore gives
\[
 (1+X)^n\equiv\prod_i(1+X^{p^i})^{n_i}\pmod p.
\]
Expanding the right side chooses indices $j_i$ with
$0\le j_i\le n_i<p$, contributing the exponent
$\sum_i j_i p^i$ and coefficient $\prod_i\binom{n_i}{j_i}$.
Uniqueness of base-$p$ division says that this exponent is $k$ exactly
when $j_i=k_i$ for every $i$.  If some digit $k_i>n_i$ no term exists
and the coefficient is zero; otherwise precisely one digit tuple
exists, yielding the formula.  This coefficient comparison is finite:
arrays of length $n+1$ suffice, with integer coefficients bounded by
$2^n$, and digit lists have length at most $n+k+1$.  For $k>n$ the
same comparison uses the zero coefficient outside degree $n$.
\end{proof}

\begin{theorem}[Kummer's carry formula]\label{el:kummer}
For a prime $p$ and $0\le k\le n$, the exponent of $p$ in
$\binom nk$ equals the number of carries when adding $k$ and $n-k$
in base $p$.  Equivalently,
\[
 (p-1)v_p\!\left(\binom nk\right)
     =s_p(k)+s_p(n-k)-s_p(n).
\]
\end{theorem}
\begin{proof}
The binomial coefficient is positive for $k\le n$, as follows from
Pascal recursion.  Taking valuations in
$k!(n-k)!\binom nk=n!$ and applying Legendre's formula gives the
displayed digit-sum identity.  To interpret it, let $a_i,b_i,n_i$
be the base-$p$ digits of $k,n-k,n$, padded with zeros, and let
$c_i$ be the carry into position $i$, with $c_0=0$.
The grade-school addition recursion is the division of
$a_i+b_i+c_i$ by $p$.  Since $a_i,b_i<p$ and $c_i\le1$, each carry
$c_{i+1}$ is zero or one, and
\[
                 a_i+b_i+c_i=n_i+p c_{i+1}.
\]
There are at most $n+2$ relevant positions, and the last carry is
zero after padding because the sum is $n$.  Sum these equations
without positional weights.  The sums of incoming and outgoing
carries agree, leaving
$s_p(k)+s_p(n-k)-s_p(n)=(p-1)\sum_i c_{i+1}$.
Cancel $p-1>0$.  The carry array is computed by a fixed bounded loop,
so both the equality and the count use only quantifier-free induction.
\end{proof}

\subsection{Arithmetic functions and finite inversion}

For $n\ge1$ and $r\in\N$ define
\[
 \tau(n)=\sum_{d=1}^n[d\mid n],\qquad
 \sigma_r(n)=\sum_{d=1}^n[d\mid n]d^r,\qquad \sigma=\sigma_1.
\]
Here $[P]$ is one or zero according as the decidable predicate $P$
holds or fails.  Define $\mu(n)=0$ when a square of a prime divides
$n$; otherwise define $\mu(n)=(-1)^t$, where $t$ is the number of
its distinct prime factors.  In particular $\mu(1)=1$.
An arithmetic function is called multiplicative if its value at one
is one and it satisfies $f(ab)=f(a)f(b)$ whenever $a,b$ are positive
and coprime.  Statements involving an unspecified $f$ below are
fixed-function schemas, or statements about supplied finite tables.

\begin{theorem}[Divisor formulas and multiplicativity]
\label{el:divisorfunc}
If $n=\prod_{i<t}p_i^{e_i}$ with distinct primes and $e_i>0$, then
\[
 \tau(n)=\prod_{i<t}(e_i+1),\qquad
 \sigma_r(n)=\prod_{i<t}\left(\sum_{j=0}^{e_i}p_i^{rj}\right).
\]
The functions $\tau$ and $\sigma_r$ are multiplicative.  For $r=1$,
\[
             (p-1)\sigma(p^e)=p^{e+1}-1.
\]
\end{theorem}
\begin{proof}
The valuation criterion gives a bijection between positive divisors
$d$ of $n$ and tuples $(j_0,\ldots,j_{t-1})$ with $0\le j_i\le e_i$:
the divisor is $\prod_i p_i^{j_i}$.  Counting tuples gives the formula
for $\tau$.  Summing $d^r=\prod_i p_i^{rj_i}$ and distributing the
finite product gives the second formula.  For coprime $a,b$ their
prime supports are disjoint, so both products concatenate; this proves
multiplicativity.  The last equation follows by subtracting the finite
geometric sum from $p$ times that sum.  All tuple codes are bounded by
$\prod_i(e_i+1)\le(n+1)^n$; direct divisor scans use only $n$ steps.
\end{proof}

\begin{theorem}[Squarefree decomposition and parity of divisor functions]
\label{el:squareparity}
Every $n\ge1$ has unique positive integers $s,t$ such that
$n=st^2$ and $s$ is squarefree (no prime square divides it).
Both are primitive recursively computable and at most $n$.
Furthermore,
\[
 \tau(n)\text{ is odd}\ \Longleftrightarrow\ n\text{ is a square},
\]
\[
 \sigma(n)\text{ is odd}\ \Longleftrightarrow\
 n\text{ is a square or twice a square}.
\]
\end{theorem}
\begin{proof}
For every prime exponent $e$ in the factorization of $n$, divide it
as $e=2q+r$ with $r\in\{0,1\}$.  Put the factor $p^r$ into $s$ and
$p^q$ into $t$.  Unique factorization proves $n=st^2$, squarefreeness,
and uniqueness of these choices.  The product formula
$\tau(n)=\prod_p(e_p+1)$ is odd exactly when every $e_p$ is even,
which is exactly the condition for $n$ to be a square.
For an odd prime $p$, the sum $1+p+\cdots+p^e$ has parity $e+1$;
for $p=2$ it is always odd.  Thus $\sigma(n)$ is odd exactly when
every odd-prime exponent is even.  In that case the exponent of
two is either even, giving a square, or odd, giving twice a square.
The converse follows from the same parity calculation.
All root witnesses are the displayed products of half-exponents,
so the existential conclusions carry explicit bounded constructions.
\end{proof}

\begin{theorem}[Totient formula and multiplicativity]
\label{el:totient}
For a prime $p$ and $e\ge1$,
$\varphi(p^e)=p^{e-1}(p-1)$.  The totient is multiplicative, and hence
\[
             \varphi(n)=\prod_{p^e\parallel n}p^{e-1}(p-1)
             \qquad(n\ge1).
\]
\end{theorem}
\begin{proof}
A residue modulo $p^e$ is noncoprime exactly when it is divisible by
$p$, using prime factorization.  Among $1,\ldots,p^e$ precisely
$p^{e-1}$ numbers are multiples of $p$.  Subtracting their count proves
the prime-power formula.  If $a,b$ are coprime, the Chinese remainder
map from residues modulo $ab$ to pairs of residues modulo $a,b$ is
bijective.  A residue is coprime to $ab$ exactly when each component
is coprime to its modulus: a prime common divisor would divide $a$
or $b$.  Restricting the finite bijection to these decidable sublists
gives $\varphi(ab)=\varphi(a)\varphi(b)$.  Factorization then proves
the product formula; the empty product gives $\varphi(1)=1$.
\end{proof}

\begin{theorem}[The divisor sum of the totient]\label{el:phisum}
For every $n\ge1$,
\[
                         \sum_{d\mid n}\varphi(d)=n.
\]
\end{theorem}
\begin{proof}
Partition $a=1,\ldots,n$ by $g=\gcd(a,n)$.  For a fixed divisor $g$
of $n$, the assignment $a=gu$ identifies this part with the integers
$1\le u\le n/g$ coprime to $n/g$: dividing a B\'ezout identity for
$a,n$ by $g$ proves one direction, and multiplying a coprimality
identity for $u,n/g$ by $g$ proves the other.  This part has
$\varphi(n/g)$ members.  The divisor map $g\mapsto n/g$ is an
involution, so counting all $n$ integers gives the stated sum.
This is a disjoint finite classification, implementable by a double
scan bounded by $n^2$.
\end{proof}

\begin{theorem}[M\"obius cancellation]\label{el:mobiuscancel}
The function $\mu$ is multiplicative and
\[
 \sum_{d\mid n}\mu(d)=\varepsilon(n),\qquad
 \varepsilon(n)=\begin{cases}1&n=1,\\0&n>1.\end{cases}
\]
\end{theorem}
\begin{proof}
For coprime inputs, their prime supports are disjoint.  If either has
a repeated prime factor, the product does too and both sides of
multiplicativity are zero.  Otherwise the numbers of prime factors
add and their parity signs multiply.  This proves multiplicativity.
For $n=1$ the divisor sum is $\mu(1)=1$.  For $n>1$, fix the least
prime $p\mid n$.  Divisors with a squared prime factor contribute
zero.  Pair each squarefree divisor not containing $p$ with its
product by $p$; both divide $n$, the pairing is reversible, and the
values of $\mu$ are opposite.  The finite sum is therefore zero.
Equivalently, distributing over the distinct prime support gives a
product of factors $1-1$.  The pairing is explicit and bounded by $n$.
\end{proof}

\begin{theorem}[Liouville's function and square detection]
\label{el:liouville}
Let $\Omega(n)$ be the length of the prime factor list, counting
multiplicity, and put $\lambda(n)=(-1)^{\Omega(n)}$ for $n\ge1$.
Then $\lambda(ab)=\lambda(a)\lambda(b)$ for all positive $a,b$, and
\[
 \sum_{d\mid n}\lambda(d)=
 \begin{cases}1&n\text{ is a square},\\0&n\text{ is not a square}.
 \end{cases}
\]
\end{theorem}
\begin{proof}
Concatenating prime factor lists and using their uniqueness gives
$\Omega(ab)=\Omega(a)+\Omega(b)$, without any coprimality assumption.
This proves complete multiplicativity of $\lambda$.  A divisor of
$n=\prod_p p^{e_p}$ chooses each exponent independently, so finite
distributivity gives
\[
 \sum_{d\mid n}\lambda(d)=\prod_{p\mid n}
             (1-1+1-\cdots+(-1)^{e_p}).
\]
Each parenthesis is one when $e_p$ is even and zero when it is odd,
by pairing consecutive terms.  The product is therefore one exactly
when every exponent is even.  This is the square criterion just proved.
For $n=1$ the empty product is one, as required.
\end{proof}

For fixed integer-valued primitive recursive functions $f,g$ define
Dirichlet convolution by
\[
                  (f*g)(n)=\sum_{d\mid n} f(d)g(n/d),\qquad n\ge1.
\]
This is a primitive recursive operation for each fixed pair of function
definitions.  One can instead take arrays containing $f(1),\ldots,f(N)$
and $g(1),\ldots,g(N)$ to obtain a single first-order assertion for
all indices at most $N$.

\begin{theorem}[Finite Dirichlet convolution algebra]
\label{el:convolution}
Dirichlet convolution is commutative and associative, has identity
$\varepsilon$, and distributes over pointwise addition.  If $f,g$
are multiplicative, then $f*g$ is multiplicative.
\end{theorem}
\begin{proof}
For commutativity replace a divisor $d$ by $n/d$; this substitution
is a bijection of the finite divisor list.  Both
$((f*g)*h)(n)$ and $(f*(g*h))(n)$ expand as
\[
                   \sum_{abc=n}f(a)g(b)h(c),
\]
where $1\le a,b,c\le n$.  The equality follows by a bijective
reindexing of this bounded triple sum, proving associativity.
Only the divisor one contributes to $(\varepsilon*f)(n)$, and
finite distributivity proves the additive rule.

If $\gcd(a,b)=1$, each divisor of $ab$ has a unique form $uv$ with
$u\mid a$ and $v\mid b$, by the prime exponent description.  The
complementary divisor is $(a/u)(b/v)$.  The two factors in each of
these products are coprime.  Multiplicativity of $f$ and $g$ therefore
gives
\[
 (f*g)(ab)=\sum_{u\mid a}\sum_{v\mid b}
 f(u)f(v)g(a/u)g(b/v)=(f*g)(a)(f*g)(b).
\]
The value at one is one.  For a formal theorem about finite tables,
only the multiplicativity identities at the bounded arguments used
in this display are hypotheses; no quantification over functions is
part of the arithmetic language.
\end{proof}

\begin{theorem}[M\"obius inversion]\label{el:mobius}
For fixed primitive recursive integer-valued $f,F$, the identities
\[
 F(n)=\sum_{d\mid n}f(d)\quad(n\ge1)
\]
are equivalent to
\[
 f(n)=\sum_{d\mid n}\mu(d)F(n/d)\quad(n\ge1).
\]
For each $N$, the same equivalence holds using only the corresponding
identities at indices $1,\ldots,N$.  In particular
\[
                   \varphi(n)=\sum_{d\mid n}\mu(d)(n/d).
\]
\end{theorem}
\begin{proof}
Let $\mathbf1(n)=1$.  M\"obius cancellation says
$\mu*\mathbf1=\varepsilon$.  If $F=\mathbf1*f$, associativity gives
$\mu*F=(\mu*\mathbf1)*f=f$.  Conversely, if $f=\mu*F$, then
$\mathbf1*f=(\mathbf1*\mu)*F=F$.
Every convolution evaluated at $n$ uses only divisors of $n$, all at
most $n$, so the proof applies verbatim to the bounded version with
$n\le N$.  Apply it to Theorem~\ref{el:phisum} to obtain the totient
formula.  This proof is therefore a quantifier-free induction on
finite sums after the bounded hypotheses have been folded into one
Boolean antecedent, rather than an induction on a functional equation.
\end{proof}

\begin{theorem}[Divisor summatory and hyperbola identities]
\label{el:summatory}
For $N\in\N$ and $r\in\N$,
\[
 \sum_{n=1}^N\sigma_r(n)=
      \sum_{d=1}^N d^r\left\lfloor\frac Nd\right\rfloor.
\]
Writing $h=\lfloor\sqrt N\rfloor$, the divisor-counting case also has
\[
 \sum_{n=1}^N\tau(n)=
       2\sum_{d=1}^{h}\left\lfloor\frac Nd\right\rfloor-h^2.
\]
\end{theorem}
\begin{proof}
Both sides of the first identity count the finite pairs $(d,k)$ with
$d,k\ge1$ and $dk\le N$, giving pair weight $d^r$.  On the left
classify by $n=dk$; on the right fix $d$ and count
$k\le\lfloor N/d\rfloor$.  This is a reindexing in the square
$[1,N]^2$.  In the unweighted case every admissible pair has $d\le h$
or $k\le h$, because otherwise $dk\ge(h+1)^2>N$.  The two regions
each have $\sum_{d=1}^{h}\lfloor N/d\rfloor$ points.  Their intersection
is the full square $[1,h]^2$, since $h^2\le N$, and has $h^2$ points.
Adding the two region counts and subtracting their intersection gives
the second formula.  For $N=0$ all regions and sums are empty.
\end{proof}

\begin{theorem}[Euclid--Euler characterization of even perfect numbers]
\label{el:perfect}
An even positive integer $N$ satisfies $\sigma(N)=2N$ exactly when
\[
                  N=2^{e-1}(2^e-1)
\]
for some $e\ge2$ such that $2^e-1$ is prime.  The exponent $e$ is then
prime.  For a supplied $N$ all witnesses and all tests are bounded
primitive recursive computations.
\end{theorem}
\begin{proof}
If $q=2^e-1$ is prime and $e\ge2$, then $q$ is odd, and
multiplicativity of $\sigma$ gives
\[
 \sigma(2^{e-1}q)=(1+2+\cdots+2^{e-1})(1+q)
                  =(2^e-1)2^e=2N.
\]
Conversely, write an even $N$ as $2^s m$ with $s\ge1$ and $m$ odd,
using $s=v_2(N)$.  The equation for perfection becomes
\[
             (2^{s+1}-1)\sigma(m)=2^{s+1}m.
\]
Put $q=2^{s+1}-1$.  Since $q$ and $2^{s+1}$ are consecutive, they are
coprime, so coprime cancellation gives $q\mid m$.  Write $m=qt$;
then the displayed equation simplifies to
$\sigma(m)=(q+1)t=m+t$.  If $t>1$, the three distinct numbers
$1,t,m$ are divisors of $m$ (here $q\ge3$), so
$\sigma(m)\ge1+t+m>t+m$, a contradiction.  Thus $t=1$ and $m=q$.
The identity $\sigma(q)=q+1$ forces $q$ to be prime, because a proper
divisor $1<d<q$ would add an additional positive summand besides
$1$ and $q$.  Taking $e=s+1$ gives the required form.
Finally, if $e=uv$ with $u,v>1$, then
\[
 2^e-1=(2^u-1)(1+2^u+\cdots+2^{u(v-1)})
\]
has two factors greater than one, contradicting primality.
The exponent $s$ is obtained by a scan bounded by $N$, and the proof
uses only factorization, finite sums, and explicit inequalities.  It
says nothing about the existence or nonexistence of odd perfect numbers.
\end{proof}

\subsection{Quantifier-free proof obligations for this chapter}

The constructions above provide witness terms, not merely assertions that
an algorithm eventually terminates.  Their common verification pattern is
as follows.  Code a loop state by a finite tuple of natural numbers; for
signed coordinates, code a sign and magnitude.  Set a total clock, extend
the transition by the identity after its stopping condition, and define
the state after $i$ steps by primitive recursion.  Express the invariant
as equality of natural-valued characteristic functions.  Its preservation
is a quantifier-free statement with the original inputs as parameters.
Apply quantifier-free induction to $i$, then decode the state at the clock.
Finite permutations and reindexings are verified by explicit bounded
maps and their inverses.  No infinite set, quotient type, arbitrary
choice function, or unbounded minimization is needed.

For clarity, some useful clocks and magnitude bounds are collected here.
They are deliberately generous; optimal bit complexity is a separate
question.
\begin{center}
\begin{tabular}{p{0.27\linewidth}p{0.62\linewidth}}
\hline
Construction & Sufficient explicit arithmetic bound\\
\hline
Division of $a$ by $b>0$ & $a$ counter steps; outputs at most $a$.\\
Extended Euclid & $a+b+1$ steps; coefficient magnitudes at most
$(a+b+2)^{a+b+1}$.\\
Prime factorization & $n$ division stages, each divisor scan at most $n$;
factor array of length at most $n$ with entries at most $n$.\\
Sieve to $N$ & At most $N^2$ crossing-off checks; bit-array code
below $2^{N+1}$.\\
CRT for $k$ moduli & Representative below $\prod_{i<k}m_i$;
prime and exponent scans below that product.\\
Orders and inverses & Residue scan below $n$ and exponent scan
at most $\varphi(n)\le n$.\\
Arithmetic functions & Divisor and prime scans bounded by $n$;
$\sigma_r(n)\le n^{r+1}$ for $n\ge1$.\\
Summatory identities & Reindexings inside $[1,N]^2$.\\
\hline
\end{tabular}
\end{center}

Each theorem concerning a positive input is encoded with a decidable
antecedent, so total definitions may return a harmless default on
nonpositive or malformed inputs.  This handles zero denominators and
invalid finite-array codes without adding a partial function to \ipra.

\clearpage
\section{Finite fields, quadratic reciprocity, and sums of squares}
\label{sec:quadratic}

This section develops three connected parts of number theory: the elementary
structure of the nonzero residues modulo a prime, quadratic reciprocity, and
representations by two and four squares.  The conventional topics and proof
traditions may be compared with \cite{lynn-qr,kumar-qr,conrad-descent}.
The proofs below include their own constructions and bounds.  In particular,
``descent'' always means an explicitly bounded iteration, not an appeal to an
unbounded least-number principle.

\subsection{Residue arithmetic and polynomial roots}

For a prime $p$, write $\mathbf F_p$ as notation for the integers
$0,\ldots,p-1$ with addition and multiplication followed by remainder modulo
$p$. This is a finite representation, not a new sort of the arithmetic.
An inverse of $a\ne0$ is the first $b<p$ with $ab\equiv1\pmod p$;
B\'ezout's identity proves that this bounded search succeeds.  In particular,
$ab=0$ in $\mathbf F_p$ implies $a=0$ or $b=0$.  Polynomials are finite
coefficient arrays and their values are computed by Horner's recurrence.
Trailing zero coefficients are removed by a bounded scan.  All assertions
about roots mean finite tests over $0,\ldots,p-1$.

\begin{theorem}[Polynomial root bound]\label{qr:roots}
A nonzero polynomial of degree $d$ over $\mathbf F_p$ has at most $d$
distinct roots.
\end{theorem}
\begin{proof}
For a constant nonzero polynomial there are no roots.  If a polynomial $f$
of positive degree has no root, the conclusion also holds.  Otherwise choose
the first root $a<p$.  Synthetic division computes an array $g$ of degree
$d-1$ such that
\[
 f(X)=(X-a)g(X)+f(a)=(X-a)g(X).
\]
For completeness, if $f=\sum_{i=0}^d c_iX^i$, set $b_{d-1}=c_d$ and,
working downwards, $b_{i-1}=c_i+ab_i$ for $1\le i<d$.  The coefficient
of $X^i$ in $(X-a)\sum_{j=0}^{d-1}b_jX^j$ is $b_{i-1}-ab_i=c_i$;
the remaining constant term agrees because evaluation at $a$ gives zero.
If $b\ne a$ is another root of $f$, cancellation in $\mathbf F_p$ gives
$g(b)=0$.  Thus the roots of $f$ consist of $a$ and at most $d-1$ further
roots, by induction.

Here the induction can be made quantifier-free without a hidden polynomial
sort. Fix $p$ and use the characteristic function which checks the assertion
for every coefficient array of length at most $d+1$ with entries below $p$.
Such arrays, padded with zeros, have codes below $p^{d+1}$.  The induction
step just given reduces the relevant array to one of length at most $d$;
all root counts are sums of $p$ zero-one values.  This is bounded induction
compiled into the quantifier-free induction scheme of \ipra.
\end{proof}

For $1\le a<p$, define $\operatorname{ord}_p(a)$ to be the first positive
$e\le p-1$ with $a^e\equiv1\pmod p$. Fermat's little theorem proves that
the search succeeds; exponentiation is always the ordinary primitive
recursive operation, or its remainder-reduced version.

\begin{lemma}[Orders and coprime products]\label{qr:orders}
If $d=\operatorname{ord}_p(a)$, then
\[
 a^k\equiv1\pmod p\quad\Longleftrightarrow\quad d\mid k
 \qquad(k\in\N),
\]
and $d\mid p-1$.  If $u,v$ have coprime orders $s,t$, respectively,
then $uv$ has order $st$.
\end{lemma}
\begin{proof}
Divide $k$ by $d$, writing $k=qd+r$ with $0\le r<d$.  If $a^k=1$,
then $a^r=1$, so the minimality in the bounded definition of $d$ forces
$r=0$.  The reverse implication follows by exponentiation. Apply this to
$k=p-1$ for the second assertion.  For the last assertion, $(uv)^{st}=1$.
If $(uv)^k=1$, raising to $t$ gives $u^{kt}=1$, hence $s\mid kt$.
Coprimality gives $s\mid k$ by B\'ezout. Similarly $t\mid k$, and therefore
$st\mid k$. The least positive such exponent is $st$.
\end{proof}

\begin{theorem}[Primitive roots modulo a prime]\label{qr:primitive}
For every prime $p$ there is $g<p$ whose powers run through all the
nonzero residues modulo $p$. A primitive recursive function can choose
such a $g$ by searching $1\le g<p$.
\end{theorem}
\begin{proof}
For $p=2$, take $g=1$. Put $N=p-1$ and factor $N$ as
$\prod_{i<r}q_i^{e_i}$ with distinct primes $q_i$. For each $i$, the
polynomial $X^{N/q_i}-1$ has degree $N/q_i<N$. By the root bound it
cannot vanish at all $N$ nonzero residues. Choose the first $a_i<p$
for which $a_i^{N/q_i}\ne1$, and set
\[
 b_i=a_i^{N/q_i^{e_i}}\pmod p.
\]
The order of $b_i$ divides $q_i^{e_i}$. If it were a proper divisor,
it would divide $q_i^{e_i-1}$, contradicting
$b_i^{q_i^{e_i-1}}=a_i^{N/q_i}\ne1$. Thus its order is exactly
$q_i^{e_i}$. The preceding lemma shows that $g=\prod_{i<r}b_i$
has order $N$. If $g^i=g^j$ for $0\le i<j<N$, cancellation gives
$g^{j-i}=1$, contradicting the order. Consequently these $N$ powers
are distinct, and finite counting shows that they exhaust the $N$
nonzero residues.

The factor list has length at most $N$, every residue search has $p$
candidates, and each exponent is at most $N$. Alternatively the final
function may simply test every $g<p$ for the $N$ required inequalities;
the argument proves that this single bounded search succeeds.
\end{proof}

\begin{corollary}[Orders of powers and the number of primitive roots]
\label{qr:primitive-count}
For a primitive root $g$ and $N=p-1$,
\[
 \operatorname{ord}_p(g^k)=\frac{N}{\gcd(N,k)}.
\]
There are exactly $\varphi(N)$ primitive roots modulo $p$. More generally,
if $b\ne0$, $k\ge1$, $b=g^t$ with $0\le t<N$, and $d=\gcd(k,N)$,
then $x^k=b$ has no nonzero solution when $d\nmid t$, and exactly $d$
nonzero solutions when $d\mid t$.
\end{corollary}
\begin{proof}
The equation $(g^k)^j=1$ is equivalent to $N\mid kj$. Dividing by
$d=\gcd(N,k)$ and using coprimality shows that this is equivalent to
$N/d\mid j$. This proves the order formula. Since $g^0,\ldots,g^{N-1}$
list the units once, exactly those exponents coprime to $N$ give order
$N$, proving the count. Finally, write $x=g^j$; then $x^k=b$ is
$kj\equiv t\pmod N$. Necessarily $d\mid t$. When $d\mid t$, division
by $d$ gives a congruence with invertible coefficient modulo $N/d$,
so it has one class $j_0$ there. Its $d$ distinct representatives modulo
$N$ are $j_0+\ell N/d$, $0\le\ell<d$.
\end{proof}

\subsection{Quadratic characters and reciprocity}

For an odd prime $p$ and an integer $a$, define the Legendre symbol by
the finite decision
\[
 \chi_p(a)=\left(\frac ap\right)=
 \begin{cases}
 0,&a\equiv0\pmod p,\\
 1,&a\not\equiv0\pmod p\ \text{and some }1\le x<p\text{ satisfies }x^2\equiv a,\\
 -1,&\text{otherwise}.
 \end{cases}
\]
The word ``otherwise'' is justified by a complete finite search, not by
excluded middle for an arbitrary proposition. The output is encoded as a
signed integer. Put $h=(p-1)/2$.

\begin{lemma}[The square residues]\label{qr:squares-half}
The $h$ residues $1^2,2^2,\ldots,h^2$ are distinct, and they are
exactly the nonzero square residues. Each has precisely two square roots
modulo $p$.
\end{lemma}
\begin{proof}
If $i^2=j^2$, then $(i-j)(i+j)=0$. For $1\le i,j\le h$, the second
factor is a positive integer less than $p$, so it cannot vanish modulo
$p$; hence $i=j$. Every nonzero residue $x$ has either $x\le h$ or
$p-x\le h$, and $x^2=(p-x)^2$. This proves exhaustion. The same
factorization shows that the only roots of $X^2-i^2$ are $i$ and $-i$;
they differ because $p$ is odd and $i\ne0$.
\end{proof}

\begin{theorem}[Euler's criterion]\label{qr:euler-criterion}
For an odd prime $p$ and $p\nmid a$,
\[
 a^{(p-1)/2}\equiv\chi_p(a)\pmod p.
\]
\end{theorem}
\begin{proof}
If $a=b^2\ne0$, Fermat's little theorem gives $a^h=b^{p-1}=1$.
Thus all $h$ nonzero square residues are roots of $X^h-1$. The root
bound permits no further roots, so a nonsquare $a$ satisfies $a^h\ne1$.
On the other hand, $(a^h)^2=a^{p-1}=1$. The identity
$(a^h-1)(a^h+1)=0$ forces $a^h=-1$. These are exactly the two
nonzero values specified by $\chi_p$.
\end{proof}

\begin{corollary}[Multiplicativity]\label{qr:character}
For all integers $a,b$,
\[
 \chi_p(ab)=\chi_p(a)\chi_p(b).
\]
\end{corollary}
\begin{proof}
If either residue is zero, both sides are zero. Otherwise Euler's criterion
makes the two sides congruent modulo $p$. Both are in $\{-1,1\}$.
Since an odd prime does not divide $2$, these two values are distinct
modulo $p$, and congruence implies equality as signed integers.
\end{proof}

\begin{theorem}[Gauss's lemma]\label{qr:gauss}
Let $p\nmid a$. For $1\le i\le h$, let $r_i$ be the residue of
$ai$ in $\{1,\ldots,p-1\}$, and let
\[
 m=\#\{i:1\le i\le h,\ 2r_i>p\}.
\]
Then $\chi_p(a)=(-1)^m$.
\end{theorem}
\begin{proof}
Define $s_i=r_i$ when $r_i\le h$ and $s_i=p-r_i$ otherwise.
These numbers lie in $\{1,\ldots,h\}$. If $s_i=s_j$, then
$ai\equiv aj$ or $ai\equiv-aj$. Cancellation gives $i\equiv j$
or $i\equiv-j$. The first gives $i=j$; the second is impossible
because $2\le i+j\le p-1$. Thus $s_1,\ldots,s_h$ is a permutation
of $1,\ldots,h$. Multiplying all congruences gives
\[
 a^h h!\equiv\prod_{i=1}^h r_i
       \equiv(-1)^m\prod_{i=1}^h s_i=(-1)^m h!\pmod p.
\]
No factor of $h!$ is divisible by $p$, so cancel it and use Euler's
criterion. The comparison of the two signs is justified as in the previous
corollary. Every permutation assertion here is a bounded statement about
$h$ entries; the product equality follows from counting each possible
entry once.
\end{proof}

\begin{lemma}[The floor-sum form]\label{qr:eisenstein}
If $a$ is a positive odd integer coprime to the odd prime $p$, then
\[
 \chi_p(a)=(-1)^{\displaystyle\sum_{i=1}^{h}\lfloor ai/p\rfloor}.
\]
\end{lemma}
\begin{proof}
Use the residues $r_i$ and folded residues $s_i$ of Gauss's lemma and
put $K=\sum_i\lfloor ai/p\rfloor$. Summing $ai=p\lfloor ai/p\rfloor+r_i$
and replacing each negative-side residue by $p-s_i$ yields the exact
integer identity
\[
 (a-1)\sum_{i=1}^{h}i
   =p(K+m)-2\sum_{\{i:r_i>h\}}s_i.
\]
Here we used $\sum_i s_i=\sum_i i$, already proved by the finite
permutation. Since $a-1$ is even and $p$ is odd, $K+m$ is even.
Therefore $(-1)^K=(-1)^m$, and Gauss's lemma proves the assertion.
The identity, rather than any geometric parity intuition, is the
quantifier-free invariant needed for formalization.
\end{proof}

\begin{theorem}[Quadratic reciprocity]\label{qr:reciprocity}
For distinct odd primes $p,q$,
\[
 \left(\frac qp\right)\left(\frac pq\right)
       =(-1)^{((p-1)/2)((q-1)/2)}.
\]
\end{theorem}
\begin{proof}
Put $h=(p-1)/2$, $k=(q-1)/2$. Consider all $hk$ pairs
$(i,j)$ with $1\le i\le h$, $1\le j\le k$. None satisfies
$qi=pj$: otherwise $p\mid i$, impossible because $0<i<p$.
For fixed $i$, the number with $pj<qi$ is $\lfloor qi/p\rfloor$.
Indeed $qi/p<q/2=k+1/2$, and $qi/p$ is not integral, so this floor
lies between $0$ and $k$. For fixed $j$, the number with $qi<pj$
is similarly $\lfloor pj/q\rfloor$. The decidable alternatives
$pj<qi$ and $qi<pj$ partition the rectangle. Consequently
\[
 \sum_{i=1}^{h}\left\lfloor\frac{qi}{p}\right\rfloor+
 \sum_{j=1}^{k}\left\lfloor\frac{pj}{q}\right\rfloor=hk.
\]
Apply the floor-sum formula twice and multiply the signs. The counting
identity is proved arithmetically by two nested bounded sums, so no
geometric theorem or principle concerning an infinite set has entered.
\end{proof}

\begin{theorem}[The two supplementary laws]\label{qr:supplements}
For every odd prime $p$,
\[
 \left(\frac{-1}{p}\right)=(-1)^{(p-1)/2},
 \qquad
 \left(\frac2p\right)=(-1)^{(p^2-1)/8}.
\]
In particular, $-1$ is a square precisely when $p\equiv1\pmod4$,
and $2$ is a square precisely when $p\equiv1$ or $7\pmod8$.
\end{theorem}
\begin{proof}
The first identity follows immediately by applying Euler's criterion to
$-1$. For the second, the residues of $2,4,\ldots,2h$ are those
integers themselves because $2h=p-1$. In Gauss's lemma, a residue
$2i$ is on the negative side exactly when $4i>p$. Hence
\[
 m=h-\lfloor p/4\rfloor.
\]
Write the odd integer $p$ in one of the four forms $8t+1,8t+3,8t+5,8t+7$.
The corresponding values of $m$ are $2t,2t+1,2t+1,2t+2$.
The integers $(p^2-1)/8$ in those four cases have the same parities:
$8t^2+2t$, $8t^2+6t+1$, $8t^2+10t+3$, and $8t^2+14t+6$.
Thus Gauss's sign is exactly the displayed sign. These four branches
also give the stated residue-class criteria.
\end{proof}

\subsection{Two squares: prime representation and exact classification}

Write $N_2(x,y)=x^2+y^2$ for signed integers $x,y$.  A pair of
nonnegative square roots representing $n$ can always be sought in the
finite box $0\le x,y\le n$.  We use signed coordinates during the
constructions because subtraction makes the identities shorter.

\begin{lemma}[The two-square product identity]\label{qr:two-product}
For all integers $x,y,r,s$,
\[
 (x^2+y^2)(r^2+s^2)=(xr+ys)^2+(xs-yr)^2.
\]
Consequently products of sums of two squares are sums of two squares.
\end{lemma}
\begin{proof}
Expanding the right side gives $x^2r^2+2xyrs+y^2s^2+x^2s^2-2xyrs+y^2r^2$.
The mixed terms cancel and the remaining four terms factor as claimed.
This is a polynomial identity over the signed-integer coding, hence an
identity provable from the elementary ring equations of \ipra.
\end{proof}

\begin{lemma}[A strict, bounded two-square descent]\label{qr:two-descent}
Let $p$ be prime and $1<m<p$. From $x^2+y^2=mp$ one can compute
integers $u,v$ and $1\le m'<m$ with $u^2+v^2=m'p$.
\end{lemma}
\begin{proof}
Choose balanced residues $r\equiv x$ and $s\equiv y\pmod m$ with
$|r|,|s|\le m/2$. They are computed by taking ordinary remainders
and subtracting $m$ exactly when twice the remainder exceeds $m$.
Since $r^2+s^2\equiv x^2+y^2\equiv0\pmod m$, write
$r^2+s^2=mm'$. The estimate $r^2+s^2\le m^2/2$ gives
$m'\le m/2<m$. If $m'=0$, both $r$ and $s$ vanish, so $m$ divides
$x$ and $y$, hence $m^2\mid mp$. Cancellation yields $m\mid p$,
contrary to $1<m<p$. Thus $m'>0$.

The numerators $xr+ys$ and $xs-yr$ are divisible by $m$, by reducing
them modulo $m$. Define
\[
 u=(xr+ys)/m,\qquad v=(xs-yr)/m.
\]
The product identity gives $u^2+v^2=m'p$. All divisions displayed here
are exact integer divisions established before they are performed.
A state has $m<p$ and $|x|,|y|<p$; the new state has the same bounds.
The intermediate products have absolute value at most $p^2$, and a
sum of two has absolute value at most $2p^2$.
\end{proof}

\begin{theorem}[Fermat's two-square theorem for primes]\label{qr:two-prime}
A prime $p$ is a sum of two integer squares if and only if $p=2$ or
$p\equiv1\pmod4$. In either positive case a primitive recursive
function returns a pair $(x,y)$ with $0\le x,y\le p$ and
$x^2+y^2=p$.
\end{theorem}
\begin{proof}
The prime $2$ is $1^2+1^2$. Suppose $p\equiv1\pmod4$.
The first supplementary law, together with the defining bounded search
for the Legendre symbol, supplies $a$ with $a^2\equiv-1\pmod p$.
Replacing $a$ by $p-a$ if needed gives $1\le a\le(p-1)/2$.
Thus $a^2+1=mp$ for some $1\le m<p$; the upper bound follows
from $a^2+1<p^2$. If $m>1$, use the preceding descent step.
Repeat this operation exactly $p$ times, freezing the state when $m=1$.
At every active step the positive multiplier drops by at least one;
therefore it must equal $1$ by the end. Absolute values of the final
coordinates give the requested pair.

Conversely, squares are $0$ or $1$ modulo $4$. If an odd prime is
a sum of two squares, exactly one square is odd, and the sum is
$1$ modulo $4$. This covers all primes.

To make the fuel argument explicit, if $m_0<p$ is the initial
multiplier and $m_t$ the multiplier after $t$ frozen-or-descending
steps, the bounded invariant is
$m_t\le\max(1,m_0-t)$, where subtraction is truncated in the
arithmetic coding. Quantifier-free induction proves it, and at $t=p$
it gives $m_p=1$. Thus the proof constructs the solution; it does not
merely negate the possibility of infinite descent.
\end{proof}

\begin{lemma}[Primes congruent to three modulo four]\label{qr:two-obstruction}
If $q$ is prime, $q\equiv3\pmod4$, and $q\mid x^2+y^2$, then
$q\mid x$ and $q\mid y$. If $n=x^2+y^2>0$, its $q$-adic
valuation $v_q(n)$ is even.
\end{lemma}
\begin{proof}
If $q\nmid y$, divide the congruence $x^2+y^2=0$ by $y^2$
in $\mathbf F_q$. It gives a square root of $-1$, contradicting the
first supplementary law. Thus $q\mid y$, and then $q\mid x^2$
forces $q\mid x$ by primality.

For the valuation statement, choose $t$ to be the smaller of the two
valuations of the nonzero coordinates, ignoring a zero coordinate.
At least one coordinate is nonzero because $n>0$. Equivalently, choose
the greatest $t\le n$ for which $q^t$ divides both coordinates;
there cannot be a larger relevant exponent, since a nonzero coordinate
has absolute value at most $n$ and $2^{n+1}>n$. Write
$x=q^t u$, $y=q^t v$. At least one of $u,v$ is not divisible by $q$.
The first part implies $q\nmid u^2+v^2$. Since
$n=q^{2t}(u^2+v^2)$, unique factorization gives $v_q(n)=2t$.
All the valuation selections used here are bounded by $n$.
\end{proof}

\begin{theorem}[The sum-of-two-squares classification]\label{qr:two-classification}
For $n>0$, the following conditions are equivalent:
\begin{enumerate}
\item $n=x^2+y^2$ for integers $x,y$;
\item every prime $q\equiv3\pmod4$ has even exponent in the
prime factorization of $n$.
\end{enumerate}
There is a primitive recursive decision procedure for this condition,
and a primitive recursive witness function returning nonnegative
$x,y\le n$ whenever it holds. Also $0=0^2+0^2$.
\end{theorem}
\begin{proof}
Necessity is the preceding lemma, applied to every prime factor of $n$.
For sufficiency, take the finite prime factorization. Each factor $2$
has representation $(1,1)$, and each prime $p\equiv1\pmod4$ has
a representation from the prime theorem. A prime factor
$q\equiv3\pmod4$ occurs to an even exponent $2e$, so the whole
factor $q^{2e}$ has representation $(q^e,0)$. Multiply these finitely
many representations using the two-square identity, starting with
$(1,0)$. The norm invariant states after each multiplication that the
current pair represents the product of the factors already processed.
The final product is $n$. Taking absolute values gives nonnegative
coordinates, each at most $n$.

The decision is a scan over primes $q\le n$ and their valuations.
For a particularly simple total witness function, scan the entire box
$\{0,\ldots,n\}^2$ and return its first representation, with $(0,0)$
as a failure value when none exists. The proof establishes success
under condition (2); the search itself has exactly $(n+1)^2$ candidates.
Thus the internal statement may use only the equality
$N_2(\operatorname{twoSq}(n))=n$ and a zero-one condition checking (2).
\end{proof}

\subsection{Four squares by finite descent}

Write $N_4(x)=x_0^2+x_1^2+x_2^2+x_3^2$. For two signed integer
quadruples define the bilinear quadruple $B(x,r)$ by
\begin{align*}
 B_0&=x_0r_0+x_1r_1+x_2r_2+x_3r_3,\\
 B_1&=x_0r_1-x_1r_0-x_2r_3+x_3r_2,\\
 B_2&=x_0r_2+x_1r_3-x_2r_0-x_3r_1,\\
 B_3&=x_0r_3-x_1r_2+x_2r_1-x_3r_0.
\end{align*}

\begin{lemma}[Euler's four-square identity]\label{qr:four-product}
For the displayed bilinear map,
\[
 N_4(B(x,r))=N_4(x)N_4(r).
\]
Moreover, if $r_i\equiv x_i\pmod m$ for all $i$ and
$m\mid N_4(x)$, then $m\mid B_i(x,r)$ for every $i$.
\end{lemma}
\begin{proof}
In the sum of the four squared coordinates, every term $x_i^2r_j^2$
occurs once. Every mixed term has a second occurrence with the opposite
sign; cancellation leaves $\sum_{i,j}x_i^2r_j^2=N_4(x)N_4(r)$.
For the divisibility assertion substitute $r_i=x_i$ modulo $m$.
The first coordinate becomes $N_4(x)$, and each remaining coordinate
becomes zero by commutativity of multiplication. This also explains the
particular sign convention in $B$: the divisibility needed for descent
is built directly into the identity.
\end{proof}

\begin{lemma}[A small represented multiple of an odd prime]\label{qr:four-seed}
For every odd prime $p$ one can find $a,b\le(p-1)/2$ and
$1\le m<p$ such that
\[
 a^2+b^2+1^2+0^2=mp.
\]
\end{lemma}
\begin{proof}
Put $h=(p-1)/2$. The residues $a^2$ for $0\le a\le h$ are
pairwise distinct: the usual factorization of $a^2-b^2$ gives this,
including the zero case. The residues $-1-b^2$ for $0\le b\le h$
are also pairwise distinct. Each list has $h+1$ elements, whereas
there are only $p=2h+1$ residues. If they had no common element,
the concatenated list would contain $p+1$ distinct residues in a box
of size $p$, contrary to finite counting. A scan of the
$(h+1)^2$ pairs therefore returns an intersection, so
$p\mid a^2+b^2+1$. This positive number is at most
$2h^2+1=(p^2-2p+3)/2<p^2$ for $p\ge3$. Division by $p$
gives the required multiplier.

The counting contradiction eliminates a decidable proposition: complete
search either finds a pair or certifies all pairs unequal. Hence it
constructs a pair intuitionistically. No infinitary pigeonhole principle
is used.
\end{proof}

\begin{lemma}[A strict four-square descent, including the even case]
\label{qr:four-descent}
If $p$ is an odd prime, $1<m<p$, and $N_4(x)=mp$, one can
compute $1\le m'<m$ and $y$ with $N_4(y)=m'p$.
\end{lemma}
\begin{proof}
First suppose $m$ is even. Since $mp$ is even, an even number of
the four coordinates is odd. There are therefore pairs of equal parity:
by a permutation, arrange $x_0\equiv x_1\pmod2$ and
$x_2\equiv x_3\pmod2$. This can be found by testing the $24$
permutations. Set
\[
 y=\left(\frac{x_0+x_1}{2},\frac{x_0-x_1}{2},
          \frac{x_2+x_3}{2},\frac{x_2-x_3}{2}\right).
\]
The coordinates are integers and
$N_4(y)=N_4(x)/2=(m/2)p$. Thus take $m'=m/2$, which is
positive and smaller than $m$.

Now suppose $m$ is odd. Choose balanced residues
$r_i\equiv x_i\pmod m$ with $|r_i|\le(m-1)/2$.
Their norm is divisible by $m$, so write $N_4(r)=mm'$.
The strict estimate
\[
 N_4(r)\le4\left(\frac{m-1}{2}\right)^2=(m-1)^2<m^2
\]
gives $m'<m$. If $m'=0$, all $r_i$ vanish, so every $x_i$
is divisible by $m$. Then $m^2\mid mp$, hence $m\mid p$,
contrary to $1<m<p$. Thus $m'\ge1$. The previous lemma makes
all coordinates of $B(x,r)$ divisible by $m$. Define
$y_i=B_i(x,r)/m$. Its norm is
\[
 N_4(y)=\frac{N_4(x)N_4(r)}{m^2}
        =\frac{(mp)(mm')}{m^2}=m'p.
\]
This completes both branches.

The even branch is essential for this proof: balanced residues for an
even modulus only give $N_4(r)\le m^2$, which need not be strict.
The parity operation handles that case directly. For every valid state,
$|x_i|<p$ because $N_4(x)<p^2$. The same holds for the output,
and every bilinear numerator has absolute value at most $4p^2$.
These bounds cover the whole finite iteration and all signed arithmetic.
\end{proof}

\begin{theorem}[Lagrange's four-square theorem]\label{qr:four-squares}
Every $n\in\N$ is a sum of four squares of nonnegative integers.
There is a primitive recursive function $\operatorname{fourSq}$ such
that, provably in \ipra,
\[
 \sum_{i=0}^3\operatorname{fourSq}_i(n)^2=n,
 \qquad \operatorname{fourSq}_i(n)\le n.
\]
\end{theorem}
\begin{proof}
For an odd prime $p$, the seed lemma supplies a represented multiple
$mp$ with $1\le m<p$. Apply the descent step when $m>1$ and
freeze the state when $m=1$. Exactly $p$ iterations suffice, by the
same invariant $m_t\le\max(1,m_0-t)$ used in the two-square proof.
This yields a representation of $p$. The prime $2$ is
$1^2+1^2+0^2+0^2$.

The numbers $0$ and $1$ have the evident representations. For $n>1$,
compute its prime factor list and multiply the prime representations
using $B$, starting from $(1,0,0,0)$. The four-square identity proves,
by induction over this list, that the current norm is the product of
the already processed factors. At the end the norm is $n$; take
absolute values of the coordinates. A nonnegative coordinate with
square at most $n$ is at most $n$.

All fuel bounds are explicit: factorization uses at most $n$ extracted
factors, each prime routine uses at most $p\le n$ descent steps,
and the final product uses at most $n$ bilinear operations. A simpler
specification of the witness function is the first quadruple in
$\{0,\ldots,n\}^4$ of norm $n$, with a default zero quadruple.
This is a bounded scan of $(n+1)^4$ candidates; the constructive proof
above establishes its success for every input. The universal theorem
therefore has a quantifier-free matrix in the language with this
primitive recursive witness function.
\end{proof}

\begin{corollary}[Four squares over the nonnegative rationals]
\label{qr:rational-four}
For integers $a\ge0$ and $b\ge1$, the rational $a/b$ is a sum of
four rational squares, with a primitive recursive construction from
$(a,b)$.
\end{corollary}
\begin{proof}
Represent $ab=x_0^2+x_1^2+x_2^2+x_3^2$ by the preceding theorem.
Then $a/b=\sum_i(x_i/b)^2$. This statement is an exact equality
after multiplication by $b^2$, so it needs no analytic theory of
rational or real numbers. The numerators satisfy $x_i\le ab$ and
the common positive denominator is the input $b$.
\end{proof}

\begin{proposition}[The three-square obstruction and optimality]
\label{qr:three-obstruction}
No integer of the form $4^a(8b+7)$, with $a,b\in\N$, is a sum
of three integer squares. In particular four squares are sometimes
necessary: $7$ cannot be expressed using three squares.
\end{proposition}
\begin{proof}
The square residues modulo $8$ are $0,1,4$. A sum of three such
residues is never $7$: with no $4$ it is at most $3$; with one $4$
it is $4,5$, or $6$; with two $4$ it is $0$ or $1$ modulo $8$;
with three it is $4$. Thus $8b+7$ is not a sum of three squares.
If a sum of three squares is divisible by $4$, all three coordinates
are even, since a square is $0$ or $1$ modulo $4$ and a sum of at
most three ones is zero modulo $4$ only when there are none. Divide
all three coordinates by $2$ and the equation by $4$, and repeat
this operation exactly $a$ times. A representation of
$4^a(8b+7)$ would thereby yield the already impossible representation
of $8b+7$. This is a bounded descent on the supplied exponent $a$.
No converse to this obstruction, and thus no proof here of the full
three-square theorem, is asserted.
\end{proof}

\subsection{Internal arithmetic statements and Lean obligations}
\label{qr:formalization}

All functions in this section are given by composition, primitive recursion,
or explicitly bounded search. For clarity, a common implementation can use
these numerical interfaces:
\begin{itemize}
\item $\operatorname{eval}(p,c,d,x)$: Horner evaluation of a base-$p$
coefficient code $c<p^{d+1}$; $\operatorname{rootCount}(p,c,d)$ is the
sum of its zero tests over $x<p$.
\item $\operatorname{order}(p,a)$: the first $1\le e\le p-1$ with
$a^e\bmod p=1$, defaulting to $0$; $\operatorname{generator}(p)$:
the first $1\le g<p$ with order $p-1$, defaulting to $0$.
\item $\operatorname{legendre}(p,a)$: the signed three-valued finite
search above; $\operatorname{fold}(p,r)=\min(r,p-r)$ for $r<p$;
$\operatorname{negCount}(p,a)$: the sum of the comparisons
$2(ai\bmod p)>p$ for $1\le i\le(p-1)/2$.
\item $\operatorname{balance}(m,x)$: the signed balanced remainder;
$\operatorname{twoStep}$ and $\operatorname{fourStep}$: the exact
formulas in the descent lemmas, with invalid states sent to a fixed
zero state. Valid states have $1\le m<p$ and norm $mp$.
\item $\operatorname{twoSq}(n)$ and $\operatorname{fourSq}(n)$:
the bounded lexicographic searches specified in the main theorems.
A tuple is coded in base $n+1$, so its search bound is respectively
$(n+1)^2$ and $(n+1)^4$.
\end{itemize}
The state records and coefficient codes are numerical encodings.  In
particular, a statement about a list means a statement about the bounded
digit tests of its code. The arithmetic predicates ``prime'', ``coprime'',
``has order'', and ``all $3\bmod4$ exponents are even'' are zero-one
primitive recursive functions. Signed equations can equally be converted
to equations between natural numbers by moving negative summands to the
opposite side.

A Lean formalization should have two distinct layers.  The first is a
semantic development on \texttt{Nat}, coded signed integers, and finite
arrays; the second constructs or checks the corresponding derivations
in the explicit syntax of \ipra. Proving a result in Lean's ambient
logic, even without invoking \texttt{Classical}, is not by itself the
second layer. The specific proof obligations are as follows.
\begin{enumerate}
\item Prove the Horner synthetic-division identity, root-set inclusion,
and counting inequality. Instantiate bounded induction on degree with
all coefficient codes below $p^{d+1}$.
\item Prove the remainder/exponent identities for orders, the finite
root-count witness for each factor of $p-1$, and the coprime-order product
lemma. No general finite-field classification theorem is needed.
\item Prove that folded residues give a bounded permutation. Derive
both the product identity for Gauss's lemma and the exact sum identity
for floor parity. Prove the reciprocity rectangle identity as equality
of two finite sums of zero-one comparisons.
\item Normalize the polynomial identities $N_2$ and $N_4$ in a
commutative-ring semantics, then supply equation-level certificates
for their use in \ipra. Prove exact divisibility before constructing
any quotient in a descent step.
\item For four-square descent, formalize the equal-parity pairing of
four integers separately from the odd-modulus balanced-residue estimate.
This prevents a weak inequality from being mistaken for strict descent.
\item Compile the frozen-loop invariant
$m_t\le\max(1,m_0-t)$ and the norm invariant to quantifier-free
induction. Finish with bounded-search success and its primitive
recursive definition certificate.
\end{enumerate}
This list is a completion route, not a claim that the Lean implementation
of all these obligations has already been completed. It is deliberately
compatible with reusing the repository's folded-residue, Gauss-product,
reciprocity, and supplementary-law derivations once their exact interfaces
and assumptions have been checked.

\clearpage
\section{Diophantine equations, continued fractions, and bounded approximation}
\label{sec:diophantine}

This chapter uses signed integers, positive denominators, finite lists, and
bounded search only. A rational is represented by $(a,b)$ with $a\in\Z$
and $b\geq1$; equality and order are cross-multiplication tests.
Rational inequalities therefore abbreviate integer inequalities.
An algorithmic square root means the integer function
\[
 \operatorname{isqrt}(n)=\max\{j\leq n:j^2\leq n\}.
\]
Its specification $s^2\leq n<(s+1)^2$ follows by a bounded scan,
including $n=0$. No real-number decision procedure is assumed.

\subsection{Rational arithmetic and power obstructions}

\begin{theorem}[Canonical rational representation]\label{dio:normal}
Every pair $(a,b)$ with $b>0$ has a unique equivalent pair $(u,v)$ with
$v>0$ and $\gcd(|u|,v)=1$. Normalization, addition, multiplication,
comparison, and reciprocal of a nonzero rational are primitive recursive.
\end{theorem}
\begin{proof}
Put $g=\gcd(|a|,b)$, $u=a/g$, and $v=b/g$. Since $b>0$, $g>0$;
the quotients are exact. A common divisor larger than $1$, multiplied
by $g$, would contradict the greatest-common-divisor property.
If $u/v=u'/v'$ and both pairs are reduced, then $uv'=u'v$.
The coprime version of Euclid's lemma gives $v\mid v'$ and $v'\mid v$.
Positivity gives $v=v'$, and cancellation gives $u=u'$.
Addition and multiplication form
$(uv'+u'v,vv')$ and $(uu',vv')$ and normalize.
Comparison tests $uv'\leq u'v$. The reciprocal switches numerator
and denominator and moves a negative sign to the numerator.
Every operation uses integer arithmetic and gcd. Zero becomes $(0,1)$.
\end{proof}

\begin{theorem}[Rational root test]\label{dio:rational-root}
Let $f(X)=c_0+c_1X+\cdots+c_dX^d$, with $d\geq1$, $c_d\ne0$,
and integer coefficients. If a reduced rational $u/v$, $v>0$, is a
root, then $v\mid c_d$; if $u\ne0$, then $u\mid c_0$.
A rational root of a monic integer polynomial is therefore an integer.
If $c_0\ne0$, all rational roots can be found by a finite search with
$|u|\leq|c_0|$ and $1\leq v\leq|c_d|$.
\end{theorem}
\begin{proof}
Multiply the equation by $v^d$:
\[
 c_0v^d+c_1uv^{d-1}+\cdots+c_du^d=0.
\]
Reduction modulo $v$ shows $v\mid c_du^d$. Coprimality gives
$\gcd(u^d,v)=1$, hence $v\mid c_d$. Reduction modulo $u$ similarly
gives $u\mid c_0v^d$, hence $u\mid c_0$. For $c_d=1$ we obtain $v=1$.
If $c_0\ne0$, zero is not a root, and the divisibilities bound every
candidate. Horner evaluation after clearing denominators checks each
candidate by an integer equality. The degree and coefficient list are
finite data, so the complete test is primitive recursive.
\end{proof}

\begin{corollary}[Rational powers and irrational roots]\label{dio:powers}
For $N\in\N$ and $k\geq2$, a nonnegative rational whose $k$th power is
$N$ is an integer. Such a rational exists precisely when
$N=t^k$ for some $0\leq t\leq N+1$.
For $N>0$, the equivalent valuation criterion is
$k\mid v_p(N)$ for every prime $p\leq N$.
In particular, if $D$ is nonsquare and $q>0$, then
$p^2-Dq^2\ne0$ for every integer $p$.
\end{corollary}
\begin{proof}
Apply the rational root test to $X^k-N$. An integer root is bounded
by $N+1$, including $N=0$. For $N>0$, unique factorization makes
$t^k=N$ equivalent to every prime exponent being divisible by $k$;
the reverse implication uses the finite product
$\prod_{p\leq N}p^{v_p(N)/k}$.
If $p^2=Dq^2$, normalize $|p|/q$ and apply the square case.
Thus the irrationality of $\sqrt D$ has the purely arithmetical content
$q>0\Longrightarrow p^2\ne Dq^2$.
\end{proof}

\subsection{Pythagorean triples and two finite descents}

\begin{theorem}[Primitive Pythagorean triples]\label{dio:pythagoras}
Suppose $a,b,c>0$, $a^2+b^2=c^2$, and $\gcd(a,b)=1$.
Exactly one leg is even. Calling it $b$, there are unique integers
$m>n>0$ such that
\[
 a=m^2-n^2,\quad b=2mn,\quad c=m^2+n^2,\quad
 \gcd(m,n)=1,\quad m\not\equiv n\pmod2.
\]
Conversely these conditions produce a primitive positive triple.
The parameters can be recovered by bounded square-root tests.
\end{theorem}
\begin{proof}
The legs cannot both be even by coprimality. They cannot both be odd,
because their squared sum would be $2$ modulo $4$. Thus $b$ is even
and $a,c$ odd. A prime dividing $a,c$ divides $b^2$, hence $b$;
therefore $\gcd(a,c)=1$. The positive integers
\[
 U=(c+a)/2,\qquad V=(c-a)/2
\]
satisfy $U>V$, $UV=(b/2)^2$, and $\gcd(U,V)=1$: a common divisor
divides $c$ and $a$. In a coprime product that is a square, each prime
exponent in each factor is even. Hence $U=m^2$ and $V=n^2$ for unique
positive $m,n$. The displayed formulas follow. Coprimality follows
from that of $U,V$, and opposite parity follows because $c$ is odd.

Conversely expand
$(m^2-n^2)^2+(2mn)^2=(m^2+n^2)^2$.
The first leg is odd. An odd prime dividing both legs divides $m$ or
$n$, and then the other, contradicting coprimality.
The positive square roots of $U,V$ prove uniqueness; they are bounded
by $c$, so recovery is primitive recursive.
\end{proof}

\begin{corollary}[All integer right triangles]\label{dio:all-triples}
Every positive integer right triangle is a positive integer multiple of
exactly one primitive triple, up to exchanging its legs.
\end{corollary}
\begin{proof}
Put $g=\gcd(a,b)$. Since $g^2\mid c^2$, comparison of prime
exponents gives $g\mid c$. Divide by $g$ and apply the preceding
theorem. Conversely scaling preserves the equation.
The scale is forced to equal the gcd of the two legs.
\end{proof}

\begin{theorem}[No integer right triangle has positive square area]
\label{dio:square-area}
There are no positive integers $a,b,c,t$ satisfying
$a^2+b^2=c^2$ and $ab=2t^2$.
\end{theorem}
\begin{proof}
We give a size-decreasing transformation with explicitly bounded
iteration. First divide a proposed triangle by $g=\gcd(a,b)$.
The primitive triangle has integral area $A_0$, since one leg is even.
Then $t^2=g^2A_0$, so $g\mid t$ by prime exponents; its area is the
square $(t/g)^2$. This normalization does not increase the hypotenuse.

For a primitive counterexample the area, in the parameters of
Theorem~\ref{dio:pythagoras}, is
\[
 mn(m+n)(m-n).
\]
The four factors are pairwise coprime. Coprimality involving $m$ or
$n$ follows from $\gcd(m,n)=1$; a common divisor of $m+n,m-n$
divides $2m,2n$, and these two factors are odd.
Thus each factor is a square, say
\[
 m=u^2,\quad n=v^2,\quad m+n=w^2,\quad m-n=z^2.
\]
Both $w,z$ are odd. Since $w^2-z^2=2v^2$ and odd squares are $1$
modulo $8$, $v$ is even. Put $r=(w+z)/2$, $s=(w-z)/2$.
These are positive integers, and
\[
 r^2+s^2=(w^2+z^2)/2=m=u^2,\qquad
 rs/2=(w^2-z^2)/8=(v/2)^2.
\]
So $(r,s,u)$ is another integer right triangle of square area, with
smaller hypotenuse because $u^2=m<m^2+n^2=c$.

The step uses gcd, exact division, and bounded square-root searches.
Make it a total primitive recursive function by returning a fixed
default tuple on invalid input. Starting at hypotenuse $c_0$, iterate
it $c_0$ times. The quantifier-free invariant is that the current tuple
is a counterexample with $1\leq c_i\leq c_0-i$. At $i=c_0$ this is
impossible. This uses quantifier-free induction, not an unrestricted
least-counterexample principle.
\end{proof}

\begin{theorem}[Fermat's quartic descent]\label{dio:fermat-four}
There are no positive integers $x,y,z$ with $x^4+y^4=z^2$.
Consequently $x^{4k}+y^{4k}=z^{4k}$ has no positive integer solution
for any $k\geq1$.
\end{theorem}
\begin{proof}
Normalize by $g=\gcd(x,y)$. The equation gives $g^4\mid z^2$,
hence $g^2\mid z$, so division yields a solution with coprime $x,y$
and no larger $z$. They cannot both be odd, by reduction modulo $16$.
Exchange them so $x$ is odd and $y$ even.
The primitive right triangle $(x^2,y^2,z)$ has parameters
\[
 x^2=m^2-n^2,\qquad y^2=2mn,\qquad z=m^2+n^2.
\]
If $m$ were even and $n$ odd, the first equality would make a square
$3$ modulo $4$. Hence $m$ is odd and $n$ even.
The coprime factors $m,2n$ have square product, so $m=u^2$ and
$2n=w^2$. The latter implies $w=2v$ and $n=2v^2$.

Now $(x,n,m)$ is a primitive right triangle: a prime dividing $x,n$
would also divide $m$. Its even leg is $2v^2$ and its hypotenuse is
$u^2$. Another parametrization gives coprime positive $r>s$ with
\[
 x=r^2-s^2,\qquad v^2=rs,\qquad u^2=r^2+s^2.
\]
Therefore $r=A^2,s=B^2$, and $A^4+B^4=u^2$ is a new positive
solution. It is smaller since $u^2=m<z=m^2+n^2$.
Normalize and apply this transformation with fuel equal to the original
$z_0$. The quantifier-free invariant records the quartic equation and
$1\leq z_i\leq z_0-i$, giving a contradiction at the last stage.
All root searches are bounded by the current hypotenuse.
For the last assertion substitute $X=x^k,Y=y^k,Z=z^{2k}$.
\end{proof}

\subsection{Finite continued fractions and Farey arithmetic}

For $a_0\geq0$ and $a_i\geq1$ for $i>0$, write
$[a_0;a_1,\ldots,a_n]$ for the usual right-nested rational fraction.
Define continuants by
\[
 \begin{gathered}
 p_{-2}=0,\quad p_{-1}=1,\quad
 q_{-2}=1,\quad q_{-1}=0,\\
 p_i=a_ip_{i-1}+p_{i-2},\qquad
 q_i=a_iq_{i-1}+q_{i-2}\quad(i\geq0).
 \end{gathered}
\]
Negative subscripts name four initial constants; an implementation
stores a four-entry state.

\begin{theorem}[Finite continued fractions for rationals]\label{dio:cf-rational}
Every nonnegative rational has a unique finite simple continued fraction
whose last coefficient is at least $2$, unless there is just one
coefficient. For $a/b$, $a\geq0,b>0$, it can be computed in at most
$b+1$ Euclidean stages, with every coefficient at most $a+b$.
\end{theorem}
\begin{proof}
Divide $a=a_0b+r$. If $r=0$, stop; otherwise replace $(a,b)$ by
$(b,r)$ and repeat. Every nonterminal positive second entry decreases,
so $b+1$ stages suffice. Substitution backwards gives the continued
fraction. Later quotients are positive. At a final stage after the
first, the first entry is larger than, and a multiple of, the second,
so the quotient is at least $2$. All entries are at most $a+b$.
For uniqueness, the tail of a nonintegral canonical fraction is greater
than $1$, so its first coefficient is the floor of its value.
Subtract the coefficient, invert the positive remainder, and repeat.
This recovers every coefficient. A final $1$ accounts exactly for the
excluded alternate expression $[\ldots,c,1]=[\ldots,c+1]$.
All floors here are rational floors, computed by integer division.
\end{proof}

\begin{theorem}[Continuant identity, coprimality, and growth]
\label{dio:continuants}
For every valid finite coefficient list,
\[
 [a_0;\ldots,a_n]=p_n/q_n,\qquad
 p_nq_{n-1}-p_{n-1}q_n=(-1)^{n+1}.
\]
In particular $\gcd(p_n,q_n)=1$ and $q_n>0$.
The denominators are nondecreasing from $q_0=1$ and strictly increase
after $q_1$. Whenever the indicated terms exist,
\[
 q_{n+2}\geq2q_n,\qquad q_{2k}\geq2^k.
\]
If every coefficient is at most $A\geq1$, then
$p_n,q_n\leq(A+1)^{n+1}$.
\end{theorem}
\begin{proof}
Induction by matrix multiplication gives
\[
 \begin{pmatrix}a_0&1\\1&0\end{pmatrix}\cdots
 \begin{pmatrix}a_n&1\\1&0\end{pmatrix}
 =
 \begin{pmatrix}p_n&p_{n-1}\\q_n&q_{n-1}\end{pmatrix}.
\]
Substitution in the nested fraction gives its value $p_n/q_n$.
Each factor has determinant $-1$, proving the determinant identity.
Any common divisor of $p_n,q_n$ divides $1$.
The recurrence gives $q_0=1,q_1=a_1\geq1$, and
$q_{n+2}\geq q_{n+1}+q_n>q_{n+1}$.
Thus $q_{n+2}\geq2q_n$, and induction gives the even-index bound.
For the upper bound start with $p_0\leq A,q_0=1$ and
$p_1\leq A^2+1\leq(A+1)^2,q_1\leq A$. Thereafter use
\[
 A(A+1)^n+(A+1)^{n-1}\leq(A+1)^{n+1}.
\]
All these are invariants of a finite matrix-product loop.
\end{proof}

\begin{theorem}[Exact tail and error identities]\label{dio:cf-error}
For a positive rational $t$ and $n\geq0$,
\[
 [a_0;\ldots,a_n,t]
 =\frac{tp_n+p_{n-1}}{tq_n+q_{n-1}},\qquad
 \left|[a_0;\ldots,a_n,t]-\frac{p_n}{q_n}\right|
 =\frac1{q_n(tq_n+q_{n-1})}.
\]
For $t\geq1$ the distance from $p_n/q_n$ is at most $1/q_n^2$.
For $n\geq1$ the value lies between adjacent convergents.
If a next coefficient $a_{n+1}$ is specified and $t\geq a_{n+1}$,
its distance is at most $1/(q_nq_{n+1})$.
\end{theorem}
\begin{proof}
Insert $t$ in the matrix identity and subtract $p_n/q_n$.
The numerator of the difference is
$p_{n-1}q_n-p_nq_{n-1}$, of absolute value $1$.
All denominators are positive. For $n\geq1$ the value is the weighted
mean of $p_n/q_n,p_{n-1}/q_{n-1}$ with positive weights
$tq_n,q_{n-1}$. Finally $tq_n+q_{n-1}\geq q_n$ if $t\geq1$,
and is at least $q_{n+1}$ under the stronger hypothesis.
For a finite extension $t$ is a positive rational, so all claims are
finite integer identities after clearing denominators.
\end{proof}

\begin{corollary}[A primitive recursive continued-fraction real]
\label{dio:cf-real}
Fix a primitive recursive coefficient function $a(n)$ with $a(0)\geq0$
and $a(n)\geq1$ for $n>0$. Its convergents satisfy
\[
 m>n\quad\Longrightarrow\quad
 |p_m/q_m-p_n/q_n|\leq q_n^{-2}.
\]
Consequently $r_k=p_{2k+2}/q_{2k+2}$ is a primitive recursive rational
name with $|r_\ell-r_k|\leq2^{-2k-2}$ whenever $\ell>k$.
\end{corollary}
\begin{proof}
For specified finite $m>n$, use the preceding theorem with
$t=[a_{n+1};\ldots,a_m]\geq1$. Then
$q_{2k+2}\geq2^{k+1}$ gives the stated estimate.
Each prefix is produced by primitive recursion.
This is a schema for each fixed primitive recursive coefficient
function and its positivity proofs; it does not say that arbitrary
reals have primitive recursive continued-fraction expansions.
The estimates are themselves the arithmetical Cauchy data, without
completion or choice.
\end{proof}

\begin{theorem}[Farey separation and the mediant]\label{dio:farey}
Suppose $a/b<c/d$, $b,d>0$, and $bc-ad=1$.
Every rational $p/q$ strictly between them, $q>0$, has $q\geq b+d$.
The mediant $(a+c)/(b+d)$ is reduced, lies strictly between them,
and attains that denominator bound.
\end{theorem}
\begin{proof}
The strict inequalities give $bp-aq\geq1$ and $cq-dp\geq1$.
Hence
\[
 q=d(bp-aq)+b(cq-dp)\geq d+b.
\]
Cross multiplication places the mediant strictly between the endpoints.
A common divisor of its numerator and denominator divides
$b(a+c)-a(b+d)=bc-ad=1$.
In particular this applies to appropriately ordered adjacent convergents.
\end{proof}

\subsection{Finite Dirichlet approximation}

\paragraph{The bounded collision principle.}
We use Theorem~\ref{fd:php}: a list of $M+1$ integers in
$\{0,\ldots,M-1\}$ contains an equal pair, which can be found by
checking at most $(M+1)^2$ pairs.
For the present applications, one can maintain a Boolean array of
length $M$ marking encountered values.
Processing a new distinct value changes exactly one zero to one.
After $j$ distinct entries the array has exactly $j$ ones, by induction.
It cannot contain $M+1$ ones. Thus the scan finds a collision.
Array counts and tests for repeated entries in a finite prefix are
bounded Boolean operations; their induction invariants are quantifier-free.

\begin{theorem}[Rational Dirichlet approximation]\label{dio:dirichlet}
For a rational $\alpha$ and $Q\geq1$, there are integers $p,q$ with
\[
 1\leq q\leq Q,\qquad |q\alpha-p|<1/Q.
\]
They can be found by a bounded scan of $Q+1$ rational fractional parts.
\end{theorem}
\begin{proof}
For $j=0,\ldots,Q$, let
$f_j=j\alpha-\lfloor j\alpha\rfloor\in[0,1)$ and assign it box
$\lfloor Qf_j\rfloor<Q$.
Two indices $i<j$ have equal box numbers. Put
$q=j-i$ and $p=\lfloor j\alpha\rfloor-\lfloor i\alpha\rfloor$.
Then $q\alpha-p=f_j-f_i$, of absolute value less than $1/Q$.
For $\alpha=a/b$, $b>0$, compute the fractional parts from
remainders of $ja$ modulo $b$. There is no real floor.
The witness satisfies $|p|\leq Q|\alpha|+1$.
\end{proof}

\begin{theorem}[Simultaneous rational Dirichlet approximation]
\label{dio:simultaneous}
For a finite rational list $\alpha_1,\ldots,\alpha_d$, $d\geq1$,
and $N\geq1$, there are $q,p_1,\ldots,p_d$ satisfying
\[
 1\leq q\leq N^d,\qquad |q\alpha_i-p_i|<1/N
 \quad(1\leq i\leq d).
\]
The construction and witness bounds are primitive recursive in the
finite input list and $N$.
\end{theorem}
\begin{proof}
For $j=0,\ldots,N^d$, form the vector of box numbers
$\lfloor N\{j\alpha_i\}\rfloor<N$. Encode it in base $N$;
for $N=1$ every vector has the single code zero.
There are $N^d$ codes and $N^d+1$ vectors. A collision at $j>k$,
with $q=j-k$, gives the result by subtracting the integer parts in
every coordinate. All errors are less than the box width $1/N$.
One may use $|p_i|\leq N^d|\alpha_i|+1$.
Exponentiation, finite-list operations, and the collision scan are
primitive recursive.
\end{proof}

\begin{theorem}[Certified transfer to represented reals]\label{dio:real-transfer}
Fix rational approximation names for a finite list of quantities
$\xi_1,\ldots,\xi_d$. Suppose a chosen precision supplies rationals
$r_i$ with certificates $|\xi_i-r_i|\leq\varepsilon$.
The rational construction produces $1\leq q\leq N^d$ and integers
$p_i$ with
\[
 |q\xi_i-p_i|<1/N+N^d\varepsilon.
\]
In particular $\varepsilon\leq1/(2N^{d+1})$ gives errors
less than $3/(2N)$.
\end{theorem}
\begin{proof}
Take witnesses for the rational list $r_i$. Then
\[
 |q\xi_i-p_i|\leq q|\xi_i-r_i|+|qr_i-p_i|
 <N^d\varepsilon+1/N.
\]
Arithmetically, replace $\xi_i$ by all later rational approximants and
the supplied error inequalities; the same calculation proves each
finite inequality, uniformly in the later precision index.
For names with error $2^{-k}$, the generous choice $k=2N^{d+1}$
suffices, since $2^k\geq k$ by quantifier-free induction.
Thus the required precision is primitive recursive.
This is a schema for fixed names and their approximation proofs.
It retains explicit slack; it neither assumes a decidable real floor
nor claims the sharp real version with the identical denominator bound.
\end{proof}

\subsection{Pell's equation with an explicit primitive recursive bound}

The following proof uses rational approximation and a finite collision.
It needs neither periodic continued fractions nor a limiting argument.

\begin{lemma}[Bounded norm with a prescribed lower denominator]
\label{dio:pell-approx}
Let $D\geq2$ be nonsquare and $Y\geq0$. Set
\[
 B=4(D+1),\qquad Q=B(Y+1),\qquad C=4D+2.
\]
There are positive integers $p,q$ with
\[
 Y<q\leq Q,\qquad p\leq DQ+1,\qquad 0<|p^2-Dq^2|<C.
\]
The pair can be obtained by bounded search using integer arithmetic.
\end{lemma}
\begin{proof}
Put $A=\operatorname{isqrt}(DQ^4)$ and $\alpha=A/Q^2$. Then
\[
 1\leq\alpha\leq D,\qquad
 0\leq D-\alpha^2<\frac{2A+1}{Q^4}
 \leq\frac{2D}{Q^2}+\frac1{Q^4}.
\]
Indeed $A\geq Q^2$, since $D\geq2$, and
$A^2\leq DQ^4\leq D^2Q^4$.
Rational Dirichlet approximation gives $1\leq q\leq Q$ and $p$ with
$|p-q\alpha|<1/Q$. Since $Q\geq12$ and $\alpha\geq1$, $p>0$.
Also $p\leq Dq+1\leq DQ+1$. Factoring a difference of squares,
\[
 \begin{split}
 |p^2-Dq^2|
 &\leq |p-q\alpha|(p+q\alpha)+q^2(D-\alpha^2)\\
 &<\frac{2Dq+1}{Q}
       +q^2\left(\frac{2D}{Q^2}+\frac1{Q^4}\right)\\
 &\leq4D+\frac1Q+\frac1{Q^2}<C.
 \end{split}
\]
It is nonzero by Corollary~\ref{dio:powers}.
If $q\leq Y$, the same estimate gives
\[
 |p^2-Dq^2|
 <\frac{2DY+1}{4(D+1)(Y+1)}
  +\frac{(2D+1)Y^2}{16(D+1)^2(Y+1)^2}<1.
\]
The first summand is less than $1/2$, since
$2DY+1\leq2D(Y+1)$; the second is less than $1/2$, since
$2D+1<8(D+1)^2$. This contradicts the nonzero integral norm.
Thus $q>Y$. Clearing the positive denominators expresses the entire
proof as integer equalities and inequalities.
\end{proof}

\begin{theorem}[Uniform bounded Pell existence]\label{dio:pell-existence}
For nonsquare $D\geq2$, define
\[
 \begin{gathered}
 C=4D+2,\quad B=4(D+1),\quad K=2C^3,\\
 U=(B+1)^{K+1},\qquad H=4(D+1)^2(U+1)^2.
 \end{gathered}
\]
There are integers
\[
 2\leq x\leq H,\qquad1\leq y\leq H,\qquad x^2-Dy^2=1.
\]
The first such pair in the displayed finite box is therefore a primitive
recursive function of $D$ on nonsquare inputs.
\end{theorem}
\begin{proof}
Let $T_0=0,T_{i+1}=B(T_i+1)$. For each $i=0,\ldots,K$ use the
preceding lemma with $Y=T_i$ to obtain $(p_i,q_i)$ and
$m_i=p_i^2-Dq_i^2$, with
\[
 T_i<q_i\leq T_{i+1},\qquad 0<|m_i|\leq C,\qquad
 p_i\leq DT_{i+1}+1.
\]
The $q_i$ strictly increase. Induction gives
$T_i+1\leq(B+1)^i$, so $q_i<U$ and $p_i\leq DU+1$.
Classify a pair by
\[
 (m_i,\ p_i\bmod|m_i|,\ q_i\bmod|m_i|).
\]
There are at most $2C^3=K$ classes: use a sign-and-magnitude tag
less than $2C$, and two residue slots each less than $C$.
A collision among the $K+1$ pairs gives $i<j$ with common nonzero
norm $m$ and matching residues modulo $|m|$.
Define
\[
 x=\frac{p_ip_j-Dq_iq_j}{m},\qquad
 y=\frac{p_iq_j-p_jq_i}{m}.
\]
The first numerator is congruent to
$p_i^2-Dq_i^2=m$ modulo $|m|$, and the second to zero.
Thus both quotients are integers, also for $m<0$.
Expansion gives
\[
 (p_ip_j-Dq_iq_j)^2-D(p_iq_j-p_jq_i)^2=m^2,
\]
so their norm is $1$.

Positivity needs checking, especially for negative common norm.
Since $q_j>q_i>0$ and all $p_i$ are positive,
\[
 (p_iq_j)^2-(p_jq_i)^2=m(q_j^2-q_i^2).
\]
The second numerator has the same sign as $m$, hence $y>0$.
Also
\[
 (p_ip_j)^2-(Dq_iq_j)^2
 =m\bigl(D(q_i^2+q_j^2)+m\bigr).
\]
The bracket is positive: for $m<0$ use
$Dq_i^2+m=p_i^2>0$, and for $m>0$ it is immediate.
The first numerator therefore has the same sign as $m$, so $x>0$.
Now $D\geq2,y>0$, and the norm equation imply $x\geq2$.
Finally, because $|m|\geq1$,
\[
 x\leq(DU+1)^2+DU^2\leq H,\qquad
 y\leq2(DU+1)U\leq H.
\]
The deliberately loose estimates are sufficient.
There are exactly $K+1$ stages and each inner construction is bounded.
Alternatively the proved existence permits a direct scan of the
$H^2$ candidate pairs. Give the function a default value when
$D<2$ or $\operatorname{isqrt}(D)^2=D$ to make it total.
\end{proof}

\begin{theorem}[Pell composition and arbitrarily many solutions]
\label{dio:pell-powers}
Define
\[
 (x,y)\mathbin{\star}(u,v)=(xu+Dyv,xv+yu),\qquad
 \mathcal N_D(x,y)=x^2-Dy^2.
\]
The norm is multiplicative. If the norm is $1$, $(x,-y)$ is an inverse.
For $u>1,v>0,u^2-Dv^2=1$, define
\[
 (X_0,Y_0)=(1,0),\qquad
 X_{n+1}=uX_n+DvY_n,\quad Y_{n+1}=vX_n+uY_n.
\]
Every pair has norm $1$ and $Y_{n+1}>Y_n$. For every specified $n$,
this primitive recursive procedure gives $n$ distinct positive solutions.
\end{theorem}
\begin{proof}
Polynomial expansion proves norm multiplicativity, associativity,
and the conjugate inverse identity. Induction proves norm $1$ and
$X_n\geq1,Y_n\geq0$. The recurrence gives
$Y_{n+1}=vX_n+uY_n>Y_n$, since $vX_n\geq1,u\geq2$.
For a size bound put $L=u+Dv+1$. Induction gives
$X_n,Y_n\leq L^n$. The infinitude assertion is this finite witness
construction for each $n$, not an assertion that an infinite set exists.
\end{proof}

\begin{theorem}[All positive Pell solutions are fundamental powers]
\label{dio:pell-all}
Let $(u,v)$ be a norm-one solution with $u>1,v>0$ whose $v$ is
smallest among positive solutions. Such a pair is found in the box of
Theorem~\ref{dio:pell-existence}. Every solution with $x>0,y\geq0$
is uniquely $(u,v)^{\star n}$ for an integer $0\leq n\leq y$.
Its exponent is recovered by a primitive recursive procedure.
\end{theorem}
\begin{proof}
For $y=0$ the equation gives $x=1$. For $y>0$, minimality gives
$y\geq v$. Divide by $(u,v)$ using the integer pair
\[
 x'=ux-Dvy,\qquad y'=uy-vx.
\]
Its norm is $1$. Furthermore
\[
 (uy)^2-(vx)^2=y^2-v^2\geq0,
\]
so $y'\geq0$. Positivity of $x'$ follows from
\[
 (ux)^2-(Dvy)^2=1+D(v^2+y^2)>0.
\]
Finally
\[
 (vx)^2-\bigl((u-1)y\bigr)^2=v^2+2(u-1)y^2>0,
\]
so $vx>(u-1)y$, which says $y'<y$.
Repeat until reaching $(1,0)$, using at most the original $y$ steps.
The invariant states that the current pair has norm $1$, has
nonnegative coordinates with positive first coordinate, and multiplied
by the already recorded fundamental power equals the original pair.
The second coordinate supplies the decreasing fuel.
Reconstruction gives the power representation.
The preceding theorem makes second coordinates of distinct powers
strictly increasing, proving uniqueness.

For the fundamental solution, first find $(x_0,y_0)$ in the box of size
$H$. Scan $1\leq y\leq y_0,2\leq x\leq x_0$.
Any competitor with smaller $y$ has smaller $x$ by $x^2=1+Dy^2$,
so this scan finds a globally minimal positive $y$.
\end{proof}

\begin{corollary}[An explicit Pell family]\label{dio:pell-family}
For $t\geq2$, the nonsquare $D=t^2-1$ has fundamental solution
$(t,1)$. All its positive solutions are powers of this pair.
\end{corollary}
\begin{proof}
The inequalities $(t-1)^2<t^2-1<t^2$ prove nonsquareness.
The identity $t^2-(t^2-1)\cdot1^2=1$ gives a solution with smallest
possible positive second coordinate. Apply the preceding theorem.
\end{proof}

\subsection{Induction audit and Lean implementation route}

Each construction has a total primitive recursive version using a bounded
scan or fixed-fuel recurrence and a default on malformed inputs.
Rational claims are proved after clearing positive denominators.
Bounded universal conditions, such as preservation of previous entries,
are Boolean folds and can be part of a quantifier-free invariant.
The descents use the original positive size as fuel. Pell existence uses
$K+1$ stages and the explicit bound $H$; Pell classification uses the
input second coordinate.

Formalization has three distinct products. First define the algorithms
and prove their ordinary Lean specifications over naturals, integers, and
explicit rational pairs. Second give primitive recursive codes and
defining equations in the chosen iPRA presentation. Third construct
iPRA derivations of the invariants and specifications.
An ordinary Lean arithmetic proof alone does not complete the third step.

Useful module boundaries are \texttt{RationalCodes},
\texttt{PythagoreanParameters}, \texttt{BoundedDescent},
\texttt{Continuants}, \texttt{FiniteDirichlet},
\texttt{PellNormCollision}, and \texttt{PellClassification}.
The separate blueprint provides concrete definitions and proof obligations.
ProveIt modules named \texttt{Pythagoras}, \texttt{FermatQuartic},
\texttt{SquareArea}, \texttt{Dirichlet}, \texttt{SimultaneousGrid}, and
\texttt{PellGeneral} are integration points subject to the repository
audit; their names alone do not certify the explicit uniform bound above.

\clearpage
\section{Finite analytic and probabilistic number theory}\label{sec:analytic}

The root article links prime distribution, the zeta function, sieves, and probabilistic number theory. This section gives exact finite statements in those directions. Every sum and product is finite, every denominator is positive, and every assertion becomes a bounded arithmetic formula after clearing denominators. The only infinite objects mentioned are the external interpretations of explicitly specified rational names.

\subsection{Infinitude in two residue classes}
\begin{theorem}[Primes in the classes $1$ and $3$ modulo $4$]\label{an:progressions}
For every $N\ge2$ there is a prime $q>N$ with $q\equiv3\pmod4$ and $q\le4N!-1$. There is also a prime $p>N$ with $p\equiv1\pmod4$ and $p\le(N!)^2+1$. Both selections are primitive recursive and their correctness is provable in \ipra.
\end{theorem}
\begin{proof}
Factor $4N!-1$. It is odd and congruent to $3$ modulo $4$. If all of its prime factors were $1$ modulo $4$, their product would be $1$ modulo $4$. Thus a bounded scan of its factorization finds a prime factor $q\equiv3\pmod4$. No prime at most $N$ divides $4N!-1$, because it divides $N!$ and cannot also divide their difference $1$.

Next factor $(N!)^2+1$. The number is odd since $N!$ is even, and again every prime factor exceeds $N$. For a prime factor $p$, the nonzero residue $N!$ squares to $-1$ modulo $p$. The first supplement to quadratic reciprocity, proved in the quadratic chapter, says that this is possible only for $p\equiv1\pmod4$. The bounds are the numbers being factored. In both arguments the failure of a desired finite scan contradicts a decidable congruence computation; this constructs a witness. This theorem is not the full theorem on primes in every coprime arithmetic progression.
\end{proof}

\subsection{Exact sieve identities}
Write $\pi(X)=\#\{p\le X:p\text{ prime}\}$ and $P(z)=\prod_{p\le z}p$, with empty product $1$. For $X,z\in\N$ put
\[
 S(X,z)=\#\{1\le n\le X:\gcd(n,P(z))=1\}.
\]
The M\"obius function $\mu$ is defined and its divisor identity is proved in the elementary chapter.

\begin{theorem}[Finite inclusion--exclusion sieve]\label{an:sieve}
\[
 S(X,z)=\sum_{d\mid P(z)}\mu(d)\left\lfloor\frac Xd\right\rfloor,
 \qquad
 \left|S(X,z)-X\prod_{p\le z}\left(1-\frac1p\right)\right|
 \le 2^{\pi(z)}.
\]
\end{theorem}
\begin{proof}
For each fixed $n$, expand the finite product
\[
 \prod_{p\le z}(1-[p\mid n]).
\]
It is $1$ exactly when no listed prime divides $n$ and is $0$ otherwise. Choosing either $1$ or the negative indicator from each factor identifies its expansion with $\sum_{d\mid P(z),\ d\mid n}\mu(d)$. Sum over $1\le n\le X$ and interchange the two finite loops. There are exactly $\lfloor X/d\rfloor$ multiples of $d$ in that interval. The same subset expansion without the divisibility indicators gives $\sum_{d\mid P(z)}\mu(d)/d=\prod_{p\le z}(1-1/p)$. Subtract this identity multiplied by $X$. Each floor discrepancy has absolute value less than one. The squarefree number $P(z)$ has $2^{\pi(z)}$ divisors, so the triangle inequality proves the displayed bound. Subsets are bit vectors of length $\pi(z)$, with a scan bound $2^{\pi(z)}$.
\end{proof}

\begin{corollary}[Exact prime counting from a finite sieve]\label{an:sievepi}
For $X\ge2$ and $z=\lfloor\sqrt X\rfloor$,
\[
 \pi(X)=S(X,z)+\pi(z)-1.
\]
\end{corollary}
\begin{proof}
A composite $n\le X$ has a prime factor at most $\sqrt n\le z$: if its least prime factor $p$ exceeded $\sqrt n$, the complementary factor $n/p$, which is at least $p$, would make $n>n$. More explicitly $p\le n/p$ implies $p^2\le n$. Thus the survivors are exactly $1$ and the primes $z<p\le X$. All comparisons can use squares, without real square roots.
\end{proof}

\subsection{Elementary quantitative prime bounds}
\begin{lemma}[A central binomial coefficient and its prime powers]\label{an:binomial}
For $n\ge1$,
\[
 \frac{4^n}{2n+1}\le \binom{2n}{n}\le4^n,
 \qquad \binom{2n}{n}\le(2n)^{\pi(2n)}.
\]
Also $\prod_{n<p\le2n}p\le4^n$.
\end{lemma}
\begin{proof}
Pascal recursion defines the binomial coefficients as natural numbers. The sum of row $2n$ is $2^{2n}=4^n$, by induction using Pascal's identity. The central coefficient is the largest: the factorial identity
\[
 (j+1)\binom{2n}{j+1}=(2n-j)\binom{2n}{j}
\]
shows nondecrease for $j<n$, and symmetry handles the other half. There are $2n+1$ entries, proving the first two inequalities.

The factorial valuation formula gives, for each prime $p\le2n$,
\[
 v_p\binom{2n}{n}
 =\sum_{j\ge1}\left(\left\lfloor\frac{2n}{p^j}\right\rfloor
       -2\left\lfloor\frac n{p^j}\right\rfloor\right).
\]
Only $j\le2n$ need be included. Write $n=ap^j+r$ with $0\le r<p^j$; the corresponding summand is $\lfloor2r/p^j\rfloor$, either $0$ or $1$. It is zero when $p^j>2n$. If $h$ is the largest exponent with $p^h\le2n$, the valuation is at most $h$, and so the entire $p$-power contribution is at most $p^h\le2n$. Multiplying over the $\pi(2n)$ primes proves the third bound. Finally every prime $n<p\le2n$ occurs in $(2n)!$ once and in $n!$ zero times. Their product therefore divides the binomial coefficient, proving the last assertion.
\end{proof}

\begin{theorem}[Dyadic Chebyshev bounds without logarithms]\label{an:chebyshev}
For every integer $K\ge3$, \ipra\ proves
\[
 2^{K-1}\ \le\ K\pi(2^K)\ \le\ 6\,2^K.
\]
\end{theorem}
\begin{proof}
Set $n=2^{K-1}\ge4$. Induction from $n=3$ proves $2n+1\le2^n$: the next left side increases by $2$, while the right side doubles and is already at least $2$. Lemma~\ref{an:binomial} therefore gives
\[
 2^n\le\frac{4^n}{2n+1}\le \binom{2n}{n}
 \le (2n)^{\pi(2n)}=2^{K\pi(2^K)}.
\]
Strict monotonicity of powers of $2$ yields $n\le K\pi(2^K)$.

For the upper bound, multiply the interval prime-product bounds for $n=1,2,4,\ldots,2^{K-1}$:
\[
 P(2^K)\le 2^{2+4+\cdots+2^K}=2^{2^{K+1}-2}.
\]
Put $L=\lceil K/2\rceil$. Every prime in $(2^L,2^K]$ is at least $2^L$, whence
\[
 L\bigl(\pi(2^K)-\pi(2^L)\bigr)\le2^{K+1}.
\]
Since $\pi(2^L)\le2^L$ and $L2^L\le2^K$, this implies $L\pi(2^K)\le3\,2^K$. The auxiliary inequality follows by writing $K=2m$ or $K=2m+1$ and using $m\le2^m$ or $m+1\le2^m$, respectively; the latter holds for $m\ge0$ by induction. Finally $K\le2L$ gives the advertised upper bound. All exponent and counting operations have finite primitive recursive definitions. These coarse constants are chosen for transparent proofs, not sharpness.
\end{proof}

\subsection{Harmonic growth and the prime reciprocal sum}
Define $H_M=\sum_{n=1}^M1/n$. Rational inequalities below mean inequalities after cross multiplication.

\begin{lemma}[Dyadic harmonic growth]\label{an:harmonic}
For $m\ge0$, $H_{2^m}\ge1+m/2$.
\end{lemma}
\begin{proof}
The block $2^{j-1}<n\le2^j$ has $2^{j-1}$ terms, each at least $2^{-j}$, so its sum is at least $1/2$. Add the $m$ blocks and the initial term $1$.
\end{proof}

\begin{lemma}[Finite Euler domination at exponent one]\label{an:eulerdom}
For $M\ge1$,
\[
 H_M\le\prod_{p\le M}\left(1-\frac1p\right)^{-1}.
\]
\end{lemma}
\begin{proof}
Expand $\prod_{p\le M}\sum_{j=0}^{M}p^{-j}$. Unique prime factorization puts every $1/n$, $1\le n\le M$, among its terms, exactly once. No exponent of such an $n$ exceeds $M$. All other terms are positive. Each finite geometric sum is at most the rational number $(1-1/p)^{-1}$, by the finite geometric-sum identity. Thus no infinite product or convergence theorem is involved.
\end{proof}

\begin{lemma}[Bounding a finite product by a reciprocal sum]\label{an:blockproduct}
Let $0\le x_i\le1/2$ be a finite list of rationals and suppose $\sum_i x_i\le k$, where $k\in\N$. Then
\[
 \prod_i(1-x_i)^{-1}\le4^{4k+1}.
\]
\end{lemma}
\begin{proof}
Process the list in order, accumulating a block as long as its sum remains at most $3/4$; otherwise start a new block. A completed nonfinal block has sum greater than $1/4$, since the next summand is at most $1/2$ and would have taken the sum above $3/4$. If there are $r$ blocks, their total is greater than $(r-1)/4$ when $r>1$, so in all cases $r\le4k+1$. The empty list also satisfies the claim.

For nonnegative $u_i$ with sum at most one, induction proves $\prod_i(1-u_i)\ge1-\sum_i u_i$. In the induction step multiply the prior inequality by $1-u\ge0$ and use the extra nonnegative product $u\sum_i u_i$. Hence each block product is at least $1/4$ and its reciprocal at most $4$. Multiply the at most $4k+1$ block inequalities. The greedy partition and its count are finite recursive objects.
\end{proof}

\begin{theorem}[Divergence of prime reciprocals with an explicit bound]\label{an:primeharmonic}
Put
\[
 B(k)=2^{\,2\cdot4^{4k+1}}.
\]
For every $k\in\N$, \ipra\ proves
\[
 \sum_{p\le B(k)}\frac1p>k.
\]
\end{theorem}
\begin{proof}
The sum is a rational, so comparison with $k$ is decidable. In the branch where it is at most $k$, Lemma~\ref{an:blockproduct} applied to $x_i=1/p$ bounds its Euler product by $4^{4k+1}$. Lemmas~\ref{an:eulerdom} and \ref{an:harmonic}, with $m=2\cdot4^{4k+1}$, instead give
\[
 \prod_{p\le B(k)}(1-1/p)^{-1}\ge H_{B(k)}
 \ge1+4^{4k+1},
\]
a contradiction. The remaining comparison branch proves the assertion. This is a uniform first-order theorem with a primitive recursive bound, rather than just a statement that no finite list contains all primes.
\end{proof}

\subsection{Zeta values and a quantitatively controlled Euler product}
For integers $s\ge2$, define
\[
 Z_s(N)=\sum_{n=1}^N\frac1{n^s},\qquad
 E_s(N)=\prod_{p\le N}\frac1{1-p^{-s}}.
\]
Both are rational functions of natural-number inputs computed by finite loops.

\begin{lemma}[A rational tail estimate]\label{an:zetatail}
For $M\ge N\ge1$ and $s\ge2$,
\[
 0\le Z_s(M)-Z_s(N)\le N^{1-s}-M^{1-s}\le N^{1-s}.
\]
\end{lemma}
\begin{proof}
For each $n\ge2$, the difference-of-powers identity gives
\[
 n^{s-1}-(n-1)^{s-1}\ge(n-1)^{s-2}.
\]
Consequently
\[
 \frac1{(n-1)^{s-1}}-\frac1{n^{s-1}}
 \ge\frac1{(n-1)n^{s-1}}\ge\frac1{n^s}.
\]
Sum from $N+1$ to $M$ and telescope. This proves the estimate entirely in rational arithmetic, with no integration argument.
\end{proof}

\begin{theorem}[Finite Euler-product error]\label{an:eulererror}
For $N\ge1$ and $s\ge2$,
\[
 0\le E_s(N)-Z_s(N)\le N^{1-s}.
\]
\end{theorem}
\begin{proof}
The lower bound follows by the same finite expansion as Lemma~\ref{an:eulerdom}, replacing $1/p$ by $1/p^s$. For the upper bound let $r=\pi(N)$ and truncate every geometric series at exponent $J$:
\[
 T_J=\prod_{p\le N}\sum_{j=0}^Jp^{-sj}.
\]
Each term in its expansion is $n^{-s}$ for a distinct positive integer $n\le P(N)^J$. Thus $T_J\le Z_s(M)$ with $M=\max(N,P(N)^J)$, and Lemma~\ref{an:zetatail} gives $T_J\le Z_s(N)+N^{1-s}$.

Each full geometric factor is at most $2$, and the difference between it and its truncation is
\[
 \frac{p^{-s(J+1)}}{1-p^{-s}}\le2^{-J}.
\]
The finite telescoping identity for a difference of products therefore gives
\[
 0\le E_s(N)-T_J\le N2^{N-J}.
\]
This is a rational expression; a negative exponent means a reciprocal power. We have proved the desired upper inequality with an extra error $N2^{N-J}$ for every $J$.

To remove this error inside arithmetic, set $\delta=E_s(N)-Z_s(N)-N^{1-s}$, and compute a positive denominator $D$ for it, without requiring lowest terms. If $\delta>0$, then $\delta\ge1/D$. Take $J=2N+D+2$. The elementary inequalities $N\le2^N$ and $D<2^{D+1}$ give $N2^{N-J}<1/D$, contradicting the preceding error estimate. The rational sign test is decidable, so $\delta\le0$. The extra truncation index is a primitive recursive function of the original input denominators. This explicit step replaces a hidden appeal to a limiting process.
\end{proof}

\begin{corollary}[An iPRA name for $\zeta(s)$ and its Euler product]\label{an:zetaname}
Let $z_s(k)=Z_s(2^{k+2})$ and $e_s(k)=E_s(2^{k+2})$. For $l\ge k$,
\[
 0\le z_s(l)-z_s(k)\le2^{-k-2},\qquad
 0\le e_s(k)-z_s(k)\le2^{-k-2},
\]
and $|e_s(l)-e_s(k)|\le3\,2^{-k-2}\le2^{-k}$. In the external real interpretation both names have the same value, denoted $\zeta(s)$, and this proves the Euler-product identity for every integer $s\ge2$.
\end{corollary}
\begin{proof}
Use the tail and product estimates with $N=2^{k+2}$, noting $N^{1-s}\le N^{-1}$. For the last estimate insert $z_s(l)$ and $z_s(k)$ between the two Euler products and use the triangle inequality. Equality of the two names means precisely the displayed errors tending to zero with this supplied modulus. The theorem does not assert analytic continuation or a statement about arbitrary complex $s$.
\end{proof}

\subsection{Finite probabilistic number theory}
Let $\omega(n)$ be the number of distinct prime divisors of $n$, and $\Omega(n)$ the number counted with multiplicity, with both zero at $n=1$.

\begin{theorem}[Exact means of prime-divisor counts]\label{an:means}
For $X\ge1$,
\[
 \sum_{n=1}^{X}\omega(n)=\sum_{p\le X}\left\lfloor\frac Xp\right\rfloor,
 \quad
 \sum_{n=1}^{X}\Omega(n)=\sum_{p\le X}\sum_{1\le j\le X}
            \left\lfloor\frac X{p^j}\right\rfloor.
\]
In particular,
\[
 \left|\frac1X\sum_{n=1}^{X}\omega(n)-\sum_{p\le X}\frac1p\right|
 \le\frac{\pi(X)}X.
\]
\end{theorem}
\begin{proof}
Expand $\omega(n)=\sum_{p\le X}[p\mid n]$ and interchange the two finite sums. For $\Omega$, the valuation $v_p(n)$ is the number of positive $j$ with $p^j\mid n$; all such $j$ are at most $X$. Count multiples once again. For the estimate divide the first equality by $X$ and bound each floor discrepancy by $1/X$. This is a theorem about a uniform finite sample $1,\ldots,X$; it requires no probability measure on an infinite space.
\end{proof}

\begin{theorem}[Exact finite independence of distinct prime divisibility]\label{an:independence}
Let $p_1,\ldots,p_r$ be distinct primes and $M=\prod_i p_i$. In the uniform finite sample $0\le n<M$, the events $p_i\mid n$ are independent and have probabilities $1/p_i$. More explicitly, for disjoint index sets $A,B\subseteq\{1,\ldots,r\}$,
\[
 \frac{\#\{n<M:\ p_i\mid n\ (i\in A),\ p_j\nmid n\ (j\in B)\}}M
 =\prod_{i\in A}\frac1{p_i}\prod_{j\in B}\left(1-\frac1{p_j}\right).
\]
\end{theorem}
\begin{proof}
The Chinese remainder map sends $n<M$ bijectively to its tuple of residues modulo the $p_i$. A specified zero coordinate has one choice, a specified nonzero coordinate has $p_i-1$ choices, and an unspecified coordinate has $p_i$ choices. Multiply the choices and divide by $M$. This is an exact finite counting identity. It does not imply that all prime-divisibility indicators on every interval are independent.
\end{proof}

\begin{theorem}[Coprime pairs and M\"obius inversion]\label{an:coprimepairs}
For $X\ge1$, the number $C(X)$ of pairs $1\le a,b\le X$ with $\gcd(a,b)=1$ satisfies
\[
 C(X)=\sum_{d=1}^X\mu(d)\left\lfloor\frac Xd\right\rfloor^2.
\]
\end{theorem}
\begin{proof}
The divisor identity $\sum_{d\mid m}\mu(d)=[m=1]$ turns the indicator of coprimality into $\sum_{d\mid a,\ d\mid b}\mu(d)$. Interchange three finite sums. For each $d$ there are $\lfloor X/d\rfloor$ choices for each coordinate. Their product gives the formula. No passage to an asymptotic density is needed for this statement.
\end{proof}

\clearpage
\section{Algebraic, \texorpdfstring{$p$}{p}-adic, and combinatorial examples}\label{sec:algebraic}

\subsection{A concrete failure of unique factorization}
The algebraic-number-theory example in the root article can be formulated without constructing the complex numbers. Let $R$ be the set of integer pairs with operations
\[
 (a,b)+(c,d)=(a+c,b+d),\qquad
 (a,b)(c,d)=(ac-5bd,ad+bc).
\]
Write $a+b\theta$ for $(a,b)$, where $\theta^2=-5$. Equality means equality of both integer coordinates. All ring laws follow by expanding integer polynomials. Define conjugation by $\overline{a+b\theta}=a-b\theta$ and norm by $N(a+b\theta)=a^2+5b^2$.

\begin{theorem}[Norm and integrality in a quadratic ring]\label{ag:norm}
For $u,v\in R$, $N(uv)=N(u)N(v)$. The norm vanishes exactly at zero, and the units are $1$ and $-1$. Every $a+b\theta$ satisfies the monic integer polynomial
\[
 X^2-2aX+(a^2+5b^2).
\]
\end{theorem}
\begin{proof}
Expand $(ac-5bd)^2+5(ad+bc)^2$ and cancel the two mixed terms. The result is $(a^2+5b^2)(c^2+5d^2)$. Both norm summands are nonnegative, so norm zero forces $a=b=0$. If $uv=1$, the positive integers $N(u),N(v)$ multiply to one and therefore both equal one. The equation $a^2+5b^2=1$ forces $b=0$ and $a=\pm1$. Conversely these two elements have inverses. Direct substitution in the displayed monic polynomial gives zero using $\theta^2=-5$. This is the exact algebraic-integrality assertion in a finite presentation; no choice of a complex square root is required.
\end{proof}

\begin{theorem}[Two inequivalent irreducible factorizations]\label{ag:nonufd}
In $R$ the four elements $2,3,1+\theta,1-\theta$ are irreducible, and
\[
 6=2\cdot3=(1+\theta)(1-\theta)
\]
are factorizations with no matching associates. In particular $2$ is irreducible but is not prime.
\end{theorem}
\begin{proof}
There is no element of norm $2$ or $3$. Indeed, $b\ne0$ gives norm at least $5$, while $b=0$ would require an integer square equal to $2$ or $3$. A nontrivial factorization of $2$, of norm $4$, would require norms $2,2$, and a nontrivial factorization of $3$, of norm $9$, would require norms $3,3$. A nontrivial factorization of $1\pm\theta$, of norm $6$, would require norms $2,3$ in some order. All are impossible. The product identity follows by expansion. Associates differ by a unit, hence by sign; none of $2,3$ is a sign of $1+\theta$ or $1-\theta$, and the two latter elements are not associates of one another. Finally $2$ divides their product $6$ but divides neither factor, since an element divisible by $2$ has both coordinates even. Thus irreducible need not imply prime outside the ordinary integers.

To read irreducibility as a first-order formula, quantify over the two candidate factors. Norm multiplicativity immediately bounds each nonzero factor norm by the fixed norm of the element; this reduces every test just used to a finite bounded case analysis. The proof therefore fits \ipra\ directly.
\end{proof}

\subsection{Finite precision \texorpdfstring{$p$}{p}-adic arithmetic}
Fix a prime $p$. For an integer $a$ and a precision $k$, define
\[
 v_{p,k}(a)=\max\{0\le j\le k:p^j\mid a\}.
\]
In particular $v_{p,k}(0)=k$. This avoids introducing an infinite value. A residue of precision $k$ is an integer in $[0,p^k)$. Compatibility of residues $a_k,a_l$ for $k\le l$ means $a_l\equiv a_k\pmod{p^k}$.

\begin{theorem}[Truncated valuation laws]\label{ag:valuation}
For integers $a,b$,
\[
 v_{p,k}(ab)=\min(k,v_{p,k}(a)+v_{p,k}(b)),\qquad
 v_{p,k}(a+b)\ge\min(v_{p,k}(a),v_{p,k}(b)).
\]
If $v_{p,k}(a)\ne v_{p,k}(b)$, equality holds in the second formula.
\end{theorem}
\begin{proof}
If a factor is zero, or has valuation at least $k$, the first formula follows immediately from divisibility by $p^k$. Otherwise write $a=p^ru$, $b=p^sv$ with $p\nmid u,v$. Euclid's lemma gives $p\nmid uv$, so the product has exact valuation $r+s$, truncated at $k$. For the sum, set $r=v_{p,k}(a)$ and $s=v_{p,k}(b)$. Divisibility by the smaller of $p^r,p^s$ gives the inequality. If, say, $r<s$, then $r<k$; write $a=p^ru$ with $p\nmid u$ and $b=p^sv$. No indivisibility condition on $v$ is required when $s=k$. Now $a+b=p^r(u+p^{s-r}v)$ and the parenthesis is nonzero modulo $p$, giving exact valuation $r$. These arguments include negative integers because divisibility ignores signs. The maximum in the definition is found by at most $k+1$ tests.
\end{proof}

For a finite integer coefficient vector $c=(c_0,\ldots,c_d)$ put
\[
 f_c(X)=\sum_{j=0}^d c_jX^j,\qquad
 f'_c(X)=\sum_{j=1}^d j c_jX^{j-1}.
\]
Horner evaluation and these derivative coefficients are uniform primitive recursive operations on the vector code.

\begin{lemma}[Taylor congruence at a finite precision]\label{ag:taylor}
If $k\ge1$, then for integers $a,t$,
\[
 f_c(a+t p^k)\equiv f_c(a)+t p^k f'_c(a)\pmod{p^{k+1}}.
\]
\end{lemma}
\begin{proof}
Expand each $(a+tp^k)^j$ by the finite binomial theorem. The terms with zero and one occurrences of $tp^k$ are precisely $a^j$ and $j a^{j-1}tp^k$. Every remaining term contains $p^{2k}$, divisible by $p^{k+1}$ since $k\ge1$. Sum the finitely many identities with the coefficients $c_j$. This is a polynomial identity with an explicitly computed remainder divisible by $p^{k+1}$.
\end{proof}

\begin{theorem}[Simple-root Hensel lifting with a unique digit]\label{ag:henselstep}
Suppose $k\ge1$, $0\le a<p^k$, $p^k\mid f_c(a)$, and $p\nmid f'_c(a)$. There is a unique $t\in[0,p)$ such that $b=a+t p^k$ satisfies $p^{k+1}\mid f_c(b)$. Consequently $b$ is the unique root modulo $p^{k+1}$ in $[0,p^{k+1})$ reducing to $a$ modulo $p^k$.
\end{theorem}
\begin{proof}
Compute the signed exact quotient $q=f_c(a)/p^k$. The derivative has an inverse $u$ modulo $p$, by B\'ezout. Choose $t\equiv-qu\pmod p$, $0\le t<p$. Lemma~\ref{ag:taylor} shows that $f_c(b)/p^k\equiv q+t f'_c(a)\equiv0\pmod p$. Thus $b$ is a root at the next precision. Its size follows from $a<p^k$ and $t<p$.

Conversely any candidate $b$ in the indicated residue class has a unique digit expression $a+t p^k$ with $0\le t<p$. Its being a root forces $q+t f'_c(a)\equiv0\pmod p$, and cancellation by the invertible derivative forces the chosen value of $t$. The inverse and digit can be found by scans of length $p$; this proves a bounded witness theorem.
\end{proof}

\begin{theorem}[A primitive recursive compatible Hensel name]\label{ag:henselname}
Given $0\le a_1<p$ with $f_c(a_1)\equiv0\pmod p$ and $f'_c(a_1)\not\equiv0\pmod p$, there is a primitive recursive sequence $a_k$, $k\ge1$, such that
\[
 0\le a_k<p^k,\qquad f_c(a_k)\equiv0\pmod{p^k},\qquad
 a_l\equiv a_k\pmod{p^k}\quad(l\ge k).
\]
It is unique among compatible root sequences with the prescribed first residue.
\end{theorem}
\begin{proof}
Iterate the previous theorem exactly $k-1$ times. At each stage the derivative remains congruent to $f'_c(a_1)$ modulo $p$, because it is an integer polynomial and all residues remain equal modulo $p$. Thus the next lift always applies. Store the bounds, root divisibility, and the finite conjunction of previous compatibility conditions as the loop invariant. Uniqueness follows by induction from uniqueness of the successive digits. If another fixed primitive recursive sequence is supplied, the same argument is an iPRA proof that the sequences agree at every index; arbitrary second-order quantification over sequences is not being asserted. Externally this compatible sequence is a $p$-adic integer root. For example, lifting $X^2-2$ from $3$ modulo $7$ begins with $3,10,108$ at moduli $7,49,343$.
\end{proof}

\subsection{Integer partitions as finite combinatorics}
A partition of $n$ is a multiplicity vector $(c_1,\ldots,c_n)$ with $0\le c_j\le n$ and $\sum_{j=1}^n j c_j=n$. The empty vector partitions zero. Counting these vectors is bounded by $(n+1)^n$ for $n\ge1$ and is primitive recursive. Denote the count by $p(n)$ and the count with parts at most $m$ by $p_m(n)$.

\begin{theorem}[Partition recursion and finite generating product]\label{ag:partitionrec}
For $m\ge1$,
\[
 p_m(n)=p_{m-1}(n)+[n\ge m]p_m(n-m),
 \quad p_0(0)=1,\quad p_0(n)=0\ (n>0).
\]
For every fixed truncation $N$, the coefficient of $X^n$, $0\le n\le N$, in
\[
 \prod_{j=1}^N\left(\sum_{e=0}^{\lfloor N/j\rfloor}X^{je}\right)
\]
is $p(n)$.
\end{theorem}
\begin{proof}
Split partitions into those having no part equal to $m$ and those having at least one. Removing one such part is a bijection from the second class to partitions of $n-m$ with parts at most $m$. The second class is empty if $n<m$. For the generating product, an expansion term chooses one multiplicity $e$ for each part size $j$ and has exponent $\sum je$. Therefore the coefficient counts exactly the partitions of that exponent. Terms above degree $N$ can be discarded throughout; polynomial multiplication is a finite coefficient convolution. The recursion is implemented by a table, first increasing $m$, then $n$; the within-row argument $n-m$ is strictly smaller than $n$ whenever it is used.
\end{proof}

\begin{theorem}[Euler's odd-parts/distinct-parts identity]\label{ag:oddpartitions}
For each $n$, the number of partitions of $n$ into odd parts equals the number into distinct parts.
\end{theorem}
\begin{proof}
For every odd part size $u$, write its multiplicity in binary, $c_u=\sum_j\epsilon_{u,j}2^j$ with digits $\epsilon_{u,j}\in\{0,1\}$. Replace these copies of $u$ by the distinct part $u2^j$ whenever $\epsilon_{u,j}=1$. Every positive integer has a unique odd part times a power of two, so outputs coming from different pairs $(u,j)$ cannot coincide. The sum is preserved because $u c_u=\sum_j\epsilon_{u,j}u2^j$.

For the inverse, factor each distinct part as $u2^j$ with $u$ odd and add $2^j$ to the multiplicity of $u$. Binary uniqueness proves that the two transformations are inverse. All part sizes and multiplicities are at most $n$, and exponents need only be scanned up to $n$. The explicit inverse bijection establishes equality of the two finite counts in \ipra.
\end{proof}

\begin{theorem}[Conjugation of partitions]\label{ag:conjugate}
Partitions of $n$ with exactly $k$ parts are equinumerous with partitions whose largest part is $k$, where the largest part of the empty partition is defined to be zero.
\end{theorem}
\begin{proof}
List the parts as $\lambda_1\ge\cdots\ge\lambda_k>0$. For $1\le j\le\lambda_1$ define $\lambda'_j=\#\{i:\lambda_i\ge j\}$. The new list is nonincreasing, has largest part $k$, and has the same sum: both sums count the pairs $(i,j)$ with $j\le\lambda_i$. Applying the construction again recovers $\lambda$, since at least $i$ rows have length at least $j$ exactly when $j\le\lambda_i$. Hence conjugation is an involutive bijection. The empty partition handles $n=k=0$; if $k=0<n$ or $k>n$, both classes are empty. Counting cells in an $n$ by $n$ box proves the identities by bounded sums.
\end{proof}

\clearpage
\section{An explicit transcendental rational name}\label{sec:transcendence}

Transcendence admits a substantial constructive example using only finite
polynomials and rapidly convergent rational approximations. The classical
constant is Liouville's $\sum_{j\ge1}10^{-j!}$; see, for example, the
construction in \cite{miller-liouville}. Here the internal theorem will
state that every nonzero integer polynomial is separated from zero on
all sufficiently late rational approximations, with a primitive recursive
index and a positive rational separation bound. The proof does not
quantify over arbitrary real numbers or decide equality of real names.

For an integer coefficient vector $c=(c_0,\ldots,c_d)$ with $c_d\ne0$,
put
\[
 f_c(X)=\sum_{j=0}^d c_jX^j,
 \qquad A(c)=\sum_{j=1}^d j|c_j|.
\]
The degree, coefficient sums, rational Horner evaluation, and every
finite search below are uniform primitive recursive functions of the
coefficient code. Trailing zeros can first be removed by a bounded scan.

\begin{theorem}[Finite polynomial estimates and a rational Liouville bound]
\label{tr:polynomial}
Let $f=f_c$ be a nonzero integer polynomial of degree $d$.
\begin{enumerate}
\item Among any $d+1$ distinct rational numbers, at least one has
nonzero value under $f$, and a bounded scan finds such an argument.
\item If $y=a/q$, where $a\in\Z$ and $q\ge1$, and $f(y)\ne0$, then
$|f(y)|\ge q^{-d}$.
\item If $x,y\in[0,1]\cap\Q$, then
$|f(x)-f(y)|\le A(c)|x-y|$. Consequently, if $A(c)>0$,
$y=a/q$, $f(y)\ne0$, and $|f(x)|\le1/(2q^d)$, then
\[
 |x-y|\ge\frac1{2A(c)q^d}.
\]
If $f(x)=0$, the stronger bound $|x-y|\ge1/(A(c)q^d)$ holds.
\end{enumerate}
\end{theorem}
\begin{proof}
For the first assertion, suppose the finite scan finds no nonzero
value at the distinct arguments $x_0,\ldots,x_d$. Synthetic division
at $x_0$ writes $f(X)=(X-x_0)g(X)$, with $g$ of degree $d-1$
and the same nonzero leading coefficient. Explicitly, for coefficients
$a_0,\ldots,a_e$ and a root $u$, start with $b_{e-1}=a_e$ and
recurse downwards by $b_{i-1}=a_i+ub_i$. Coefficient comparison gives
$f(X)=(X-u)\sum_{i<e}b_iX^i+(a_0+ub_0)$, and the remainder
is $f(u)=0$. At every other listed root,
$x_i-x_0\ne0$, so cancellation of nonzero rational numbers shows
$g(x_i)=0$. Repeat at $x_1$, then $x_2$, and so on. After exactly
$d$ divisions the remaining polynomial is the original nonzero leading
coefficient, but still vanishes at $x_d$, a contradiction. For $d=0$
the contradiction is immediate. Thus the decidable failure branch of
the finite scan is impossible and a witness is returned.

This is a clocked algorithm on rational coefficient arrays. Its invariant
after $i\le d$ divisions is: the current degree is $d-i$, its leading
coefficient is the original leading coefficient, and its values at the
remaining listed arguments vanish. All argument indices are bounded by
$d$; rational identities are cleared of denominators. Hence the proof
uses bounded, and therefore quantifier-free-compilable, induction.

For the second assertion, clearing the positive denominator gives
\[
 q^d f(a/q)=\sum_{j=0}^d c_j a^j q^{d-j}\in\Z.
\]
This integer is nonzero, so its absolute value is at least one.
Dividing by $q^d$ proves the inequality; the fraction $a/q$ need not
be in lowest terms.

For the third assertion use, for $j\ge1$,
\[
 x^j-y^j=(x-y)\sum_{i=0}^{j-1}x^{j-1-i}y^i.
\]
Each summand in the finite sum belongs to $[0,1]$, so
$|x^j-y^j|\le j|x-y|$. Multiplication by $|c_j|$, summation,
and the triangle inequality give the claimed Lipschitz bound. Finally
\[
 q^{-d}\le |f(y)|\le |f(x)|+A(c)|x-y|.
\]
Insert the stated upper bound for $|f(x)|$ and rearrange. When
$f(x)=0$, omit that term instead. These are exact rational inequalities;
no mean-value theorem or differentiability is being used.
\end{proof}

For $n\ge1$, define the rational truncations and their common denominators
by
\[
 r_n=\sum_{j=1}^n10^{-j!}=\frac{P_n}{q_n},\qquad
 q_n=10^{n!},\qquad P_n=\sum_{j=1}^n10^{n!-j!}.
\]
All factorials, powers, sums, and integer quotients here have fixed
primitive recursive definitions.

\begin{theorem}[The Liouville name and its approximation bounds]
\label{tr:name}
The rational sequence $r_n$ is strictly increasing and lies in $(0,1)$.
For $m\ge n\ge1$,
\[
 0\le r_m-r_n\le2\,10^{-(n+1)!}\le2^{-n}.
\]
Thus $u(k)=r_{k+1}$ is a primitive recursive rational name with a
specified Cauchy modulus. Moreover, for every $h\ge1$, every $n\ge h$,
and every $m>n$,
\[
 0<r_m-r_n<q_n^{-h}.
\]
The truncation $r_n=P_n/q_n$ is in lowest terms.
\end{theorem}
\begin{proof}
Every added term is positive. Since $j!\ge j$ for $j\ge1$, the
partial sum is bounded by the finite geometric sum
$\sum_{j=1}^n10^{-j}<1/9<1$. For $j\ge n+1$, successive factorials
increase by at least one, so
\[
 j!\ge(n+1)!+(j-n-1).
\]
Consequently, using only a finite geometric sum,
\[
 r_m-r_n\le10^{-(n+1)!}
       \sum_{i=0}^{m-n-1}10^{-i}
       \le\frac{10}{9}\,10^{-(n+1)!}
       <2\,10^{-(n+1)!}
\]
when $m>n$; for $m=n$ the difference is zero. Since
$(n+1)!\ge n+1$ and $10^{n+1}\ge2^{n+1}$, the latter bound is
at most $2^{-n}$. This proves the asserted modulus directly.

If $n\ge h$, then $q_n\ge10$ and
\[
 2\,10^{-(n+1)!}=2q_n^{-(n+1)}<q_n^{-h},
\]
because $n+1-h\ge1$ and therefore $q_n^{n+1-h}\ge10>2$.
Finally, the last summand in $P_n$ is one and each earlier summand
is divisible by ten. Thus $P_n\equiv1\pmod{10}$, so it is
coprime to $q_n=2^{n!}5^{n!}$. This gives the lowest-terms assertion.
The external interpretation of the name is the real number commonly
written $L=0.110001000000000000000001\ldots$; the internal content
is precisely the rational estimates just proved.
\end{proof}

\begin{theorem}[Uniform polynomial separation and transcendence]
\label{tr:transcendence}
There are primitive recursive functions which, from any nonzero integer
polynomial $f_c$, return an index $n\ge1$ and a positive integer $D$
such that \ipra\ proves
\[
 \forall m\ge n\qquad |f_c(r_m)|\ge\frac1D.
\]
For a nonconstant polynomial, explicit choices are
\[
 A=A(c),\qquad n_0=d+A+1,\qquad
 Q=10^{(n_0+d)!},\qquad D=2Q^d,
\]
with $n$ selected from $n_0,\ldots,n_0+d$. The external real number
represented by $r_n$ is therefore transcendental.
\end{theorem}
\begin{proof}
If $d=0$, every value is the nonzero integer $c_0$, so choose
$n=D=1$. Suppose now $d\ge1$. Then $A\ge1$ because the leading
coefficient is nonzero. The $d+1$ rational truncations
$r_{n_0},\ldots,r_{n_0+d}$ are distinct. The polynomial-root bound
therefore makes the search for the first index $n$ with $f_c(r_n)\ne0$
succeed. Writing $q=q_n$, the rational-value bound gives
$|f_c(r_n)|\ge q^{-d}$.

For $m\ge n$, the Lipschitz estimate and the tail bound give
\[
 |f_c(r_m)-f_c(r_n)|\le 2A\,10^{-(n+1)!}.
\]
We check the required numerical domination explicitly. Since
$n\ge d+A+1$, one has $n+1-d\ge A+2$. Thus
\[
 (n+1-d)n!\ge A+2,
 \qquad 10^{(n+1-d)n!}\ge10^{A+1}>4A.
\]
For the last inequality, elementary induction gives $10^A\ge A+1$,
so $10^{A+1}\ge10(A+1)>4A$. Using
$(n+1)!-dn!=(n+1-d)n!$, it follows that
\[
 2A\,10^{-(n+1)!}\le\frac12\,10^{-dn!}
       =\frac12q^{-d}.
\]
The triangle inequality now yields
\[
 |f_c(r_m)|\ge |f_c(r_n)|-|f_c(r_m)-f_c(r_n)|
              \ge\frac12q^{-d}\ge\frac1{2Q^d},
\]
since $n\le n_0+d$ implies $q\le Q$. This is the announced
uniform margin, independent of the later index $m$.

The witness algorithm has exactly $d+1$ candidate evaluations and uses
only factorial, exponentiation, rational polynomial evaluation, and a
bounded nonzero test. Its fallback value on the zero polynomial can
be fixed arbitrarily; the theorem assumes the decidable nonzero
coefficient condition. The proof for any later $m$ uses finite sums
with that $m$ as a numerical parameter. No universal evaluator for
arbitrary programs is introduced: only coefficient-coded polynomials
are evaluated.

For the external interpretation, polynomial continuity follows already
from the displayed Lipschitz estimate on $[0,1]$. Hence the same positive
margin separates $f_c(L)$ from zero for every nonzero integer
polynomial. This is exactly transcendence. Internally the stronger,
explicitly witnessed eventual separation above is the theorem; no
second-order object or unrestricted real-zero test is needed.
\end{proof}


The concrete Lean encoding and denominator-cleared endpoint are described in Section~\ref{sec:lean}.

\clearpage
\section{Lean formalization: definitions, targets, and completion route}
\label{sec:lean}

\subsection{Three products that must not be conflated}
There are three separate artifacts in a complete formalization. A \emph{semantic specification} defines the algorithms on Lean's \texttt{Nat}, \texttt{Int}, rational pairs, and finite lists, and states their mathematical correctness. A \emph{primitive recursive representation} gives the fixed function symbols and their defining equations. An \emph{object derivation} is a proof term in the explicit first-order calculus, with only the axioms of \ipra. An ambient Lean proof of a theorem about natural numbers does not automatically produce the third artifact. Even avoiding Lean's \texttt{Classical} namespace would not remove this distinction: unrestricted inductive propositions and higher-order constructions are available in the ambient metatheory.

The delivered files provide concrete definitions and a detailed dependency plan. Their status is \textbf{uncompiled blueprint}. The runtime used to prepare this report had Python and \LaTeX, but no \texttt{lean}, \texttt{lake}, or \texttt{elan} executable. Therefore this package does not claim any new Lean kernel check. The existing repository sources were read at a pinned commit; their import closure was not rebuilt here. A record called \texttt{AnalyticGoals} or \texttt{RequiredDerivations} lists obligations and has no constructed inhabitant: such a record is not an axiom declaring the obligations true.

\subsection{Concrete semantic interfaces}
\texttt{ElementaryDefinitions.lean} defines fixed-clock iteration, bounded divisibility and primality bits, division states, an extended Euclidean state with two coefficient rows, factor lists, valuation, congruence, totient and divisor sums, M\"obius values, convolution, a sieve, finite CRT checking, bounded order search, digit extraction, and the Lucas product. The definitions have explicit defaults outside their mathematical domains. For example, the mathematical specification for \texttt{divide a b} assumes $b>0$; the definition still returns a total value for $b=0$.

\texttt{AnalyticDefinitions.lean} defines prime lists, $\pi$, primorials, harmonic and zeta partial sums, rational Euler products and their truncations, the explicit prime-reciprocal bound, prime-divisor counts, finite sieve survivors, and the M\"obius formulas. The exact-rational semantics can initially use \texttt{Rat}; the object translation replaces a rational by signed numerator and positive denominator, with all comparisons cleared of denominators.

\texttt{AlgebraicDefinitions.lean} defines the quadratic integer-pair multiplication and norm, truncated valuations, polynomial evaluation and formal derivative, a bounded Hensel digit search, compatible lifted residues, and finite multiplicity-vector partition predicates. Derivative coefficients are computed from the input coefficient vector. This is a restricted uniform polynomial interpreter, whose primitive recursiveness follows from its explicit finite degree loop.

The Diophantine blueprint specifies rational normalization, continuants, approximation witnesses, the bounded norm-collision construction, and Pell classification. \texttt{QuadraticDefinitions.lean} adds residue searches, the Legendre symbol, primitive-root predicates, and bounded two-square and four-square witness searches; the quadratic chapter specifies the finite polynomial operations and both descent steps. \texttt{TranscendenceDefinitions.lean} supplies factorial truncations, rational polynomial evaluation, the coefficient bound, and the explicit separation search and denominator. The report itself defines every mathematical predicate and construction used in the goals; the machine-readable theorem inventory records the label and source file of each result.

For a concrete example, a four-square witness can be defined without any complicated recursive control flow: enumerate $0\le a,b,c,d\le n$, encoded by four base-$(n+1)$ digits; choose the first tuple whose squared norm is $n$; return the zero tuple on failure. The scan has $(n+1)^4$ candidates. The hard proof obligation is that it succeeds, which is supplied mathematically by the prime-seed and descent proofs. Giving a bounded definition alone would not prove its success.

\subsection{The exact ProveIt object language}
At the pinned revision, \texttt{ConciseIPRA.PR} defines the type of function symbols by the following constructors:
\begin{lstlisting}[language={}]
inductive PR : Nat -> Type
  | zero : PR 0
  | succ : PR 1
  | proj {k : Nat} (i : Fin k) : PR k
  | comp {m k : Nat} (f : PR m) (gs : Fin m -> PR k) : PR k
  | prec {k : Nat} (f : PR k) (g : PR (k + 2)) : PR (k + 1)
\end{lstlisting}
Terms \texttt{PTm n} have $n$ free variables and are built from variables and applications of these symbols. The symbols are syntactic descriptions, not object-level arguments to a universal evaluator. The primitive recursion constructor takes its recursion variable first. The first-order theory \texttt{PRAAx} contains the projection, composition, recursion, successor, and quantifier-free-induction axioms described in Section~\ref{sec:foundations}.

The inspected \texttt{Foundations.BoundedSyntax} module gives an arity-indexed bounded formula type \texttt{BFm n}. It has equality, conjunction, disjunction, implication, negation, bounded universal, and bounded existential constructors. The module defines \texttt{formula} for the actual first-order interpretation and \texttt{bitF} for a compiled primitive recursive Boolean symbol. Its source contains the endpoint
\begin{lstlisting}[language={}]
boolean_reflection_closed (A : BFm n) :
  ([] : List (Fm PTm 0)) |- [PRAAx]
    univ n (iff (app (bitF A) var = one) (formula A))
\end{lstlisting}
where the listing uses plain ASCII for the source's object equality and derivability notation. This is a derivation of Boolean reflection, not an extra reflection axiom. Its compilation is metatheoretic recursion on a fixed finite formula. This is exactly the infrastructure required for the bounded invariants in this report.

\texttt{Core.Search.search\_found} finds a \emph{zero} of its supplied term. By contrast the characteristic functions in this report use $1$ for truth. A port must therefore supply the complement $1-[A]$, or explicitly use the corresponding zero-specification. Reversing these conventions silently would invalidate a witness routine. The exact names and bit conventions were checked in the source, rather than inferred from a filename.

\texttt{ProveItTargets.lean} defines actual formulas in the object language, parameterized by the still-to-be-constructed primitive recursive witness codes. For example its four-square body has one free variable and is the literal object equality
\[
 w_0(n)^2+w_1(n)^2+w_2(n)^2+w_3(n)^2=n.
\]
Its universal closure is a closed \texttt{Fm PTm 0}. The Pell body uses a Boolean nonsquare guard and the equality $x(D)^2=1+Dy(D)^2$, together with $y(D)>0$. The file also reuses the inspected bounded-reflection theorem in a small wrapper; that wrapper remains uncompiled here. The missing theorem-specific derivations are explicitly fields of a requirement record.

\subsection{Source audit and reusable interfaces}
The source paths below are relative to
\path{Logic/PrimitiveRecursiveArithmetic/Derivations/Lean/IPRADerivations/}.
They were read at commit \texttt{bbb85238dac533fc4b7b84a750cca864fe35fed5}.
\begin{longtable}{>{\raggedright\arraybackslash}p{0.43\linewidth}>{\raggedright\arraybackslash}p{0.50\linewidth}}
\toprule
Inspected source & Useful interface or precise role\\
\midrule\endhead
\path{Core/Definitions.lean} & \texttt{PTm.toPR}, \texttt{recDef}, \texttt{sumF}, \texttt{prodF}; converts fixed arithmetic terms into function symbols.\\
\path{Core/Division.lean} & \texttt{div\_mod}, \texttt{euclid\_unique}, \texttt{div\_mod\_unique}; quotient and remainder equations.\\
\path{Core/Bezout.lean} & \texttt{Bezout.invariant\_open} and \texttt{Bezout.bezout}; a closed B\'ezout derivation.\\
\path{Core/Search.lean} & \texttt{Search.search\_spec}, \texttt{search\_minimal}, \texttt{search\_found}; bounded zero search.\\
\path{Core/BoundedSums.lean} & \texttt{sum\_zero\_elim}, \texttt{sum\_zero\_of\_bounded}; finite sum reasoning.\\
\path{Core/Pigeonhole.lean} & \texttt{Pigeonhole.pigeonhole}; explicit collision selectors and their object-level specification.\\
\path{Core/Descent.lean} & \texttt{Descent.descent}; a reusable finite descent interface with hypotheses that must be discharged.\\
\path{Foundations/BoundedSyntax.lean} & \texttt{BFm}, \texttt{bitF}, \texttt{boolean\_reflection\_closed}; the bounded-formula compiler and derivation.\\
\path{DiophantineNumberTheory/}\newline\path{QuadraticReciprocity.lean} & \texttt{quadratic\_reciprocity}; a closed proof of the separately defined reciprocity sentence.\\
\path{DiophantineNumberTheory/}\newline\path{PellGeneral.lean} & Positive- and negative-norm reduction from a \emph{supplied} norm-one unit; does not by its filename prove the explicit bound in Theorem~\ref{dio:pell-existence}.\\
\path{AnalyticNumberTheory/}\newline\path{FinalPNT.lean} & \texttt{arith\_pnt}; a source endpoint for a stronger arithmetic prime-number estimate, recorded below without claiming a fresh rebuild.\\
\bottomrule
\end{longtable}

The actual toolchain file specifies Lean \texttt{v4.32.0}; the lockfile pins Mathlib to
\begin{center}\small\ttfamily 81a5d257c8e410db227a6665ed08f64fea08e997.\end{center}
These exact files, rather than a potentially older prose dependency summary, determine the proposed replay environment. The source audit is selective: it checks interfaces and visible definitions, not every transitive imported proof. No new theorem in this report depends on treating an unverified research packet as an axiom.

\subsection{A precise completion schedule}
The following work order minimizes repeated encoding and proof work.
\begin{enumerate}
\item \textbf{Freeze the syntax and encodings.} Use the pinned \texttt{PR}, \texttt{PTm}, and \texttt{BFm} frontend. Represent integers by the repository's existing signed arithmetic layer, and rational values by integer numerator/positive denominator pairs. Prove agreement with the semantic definitions for every operation actually used. Select one list representation and prove its bounded access and update equations.
\item \textbf{Port the shared loops.} Instantiate the proved Boolean compiler for primality, divisibility, squarehood, tuple validity, and the loop invariants. Reuse bounded zero search with the correct complemented truth bit. Implement frozen iteration, list factorization, bounded maxima, and finite sum/product templates as fixed primitive recursive syntax. Derive their initial and successor equations.
\item \textbf{Discharge arithmetic invariants.} Reuse the division and B\'ezout derivations. Establish prime factorization with its remainder/product invariant, finite list uniqueness, and valuation formulas. Prove finite permutations and finite sum reindexings via coded inverse maps. These finish the elementary chapter and supply nearly all arithmetic dependencies for the later chapters.
\item \textbf{Build the two main descent packets.} Compile the two-square and four-square norm identities into integer polynomial certificates. Prove the exact divisibility of every numerator before defining its quotient. Treat the even four-square multiplier separately from the odd multiplier. Use the fixed-clock invariant for the multiplier and then prove bounded witness-search success.
\item \textbf{Build the Pell packet.} Establish the rational box collision, the norm estimate with $B=4(D+1)$ and $C=4D+2$, and the separation $q>Y$. Compile the $K+1$-stage sequence, where $K=2C^3$. Prove that equal class codes yield integral quotients, verify both signs by the two squared-difference identities, and discharge the explicit $H$ bound. This proves the new bound independently of any continued-fraction-period theorem.
\item \textbf{Build the analytic packet.} Formalize the finite binomial valuation bound, the dyadic product recurrence, and the block partition recurrence for reciprocal products. For the Euler product, define its rational denominator and the primitive recursive error-removal index $J=2N+D+2$. The decisive obligation is the discrete fact that a positive rational with denominator $D$ is at least $1/D$.
\item \textbf{Build the remaining finite packets.} Prove the Hensel Taylor remainder certificate, single-digit uniqueness, and the compatible residue loop. Verify the partition bijections by finite inverse transformations. The explicit transcendence construction, when read as separation of a fixed rational name from integer polynomial zeros, uses the same finite-polynomial and error-budget infrastructure.
\item \textbf{Close and audit endpoints.} Assemble a derivation of each closed translated sentence, print its exact type and its Lean axiom dependencies, check that no \texttt{sorryAx} appears, and check the entire selected import closure. A clean \texttt{\#print axioms} for an ambient semantic theorem is not a substitute for inspecting that the conclusion really has the type \texttt{Prf PRAAx []} for the intended sentence.
\end{enumerate}

The packet-level manifests and the theorem inventory make this schedule concrete: every numbered result has a source label, a proof location, and the mathematical construction to port. The provided definitions may be optimized later without changing the mathematical obligations; the deliberately simple bounded searches are preferable first targets because they expose all quantifiers and default cases.

\paragraph{A concrete transcendence endpoint.}
Use integer coefficient vectors and rational pairs with positive
denominators. Formalize Horner evaluation and synthetic division, proving
the clocked remaining-roots invariant of Theorem~\ref{tr:polynomial}.
Establish the finite power-difference identity and its rational inequality
by induction on the exponent. Define $P_n,q_n$ by the displayed finite
sums and prove the factorial-gap and finite-geometric inequalities.
Finally define the $d+1$-entry search and its margin $D$ as numerical
primitive recursive terms. The final object-theory formula has the
nonzero-vector test as antecedent and, after denominator clearing,
\[
 m\ge n(c)\ \longrightarrow\ D(c)|\operatorname{num}(f_c(r_m))|
                    \ge\operatorname{den}(f_c(r_m)).
\]
The positivity of the denominator is part of the rational-code invariant.
These are concrete semantic and derivation obligations, not claims of
completed Lean certification.

\subsection{A stronger repository endpoint, kept separate}
The inspected \texttt{FinalPNT.lean} source states a derivation of the arithmetic estimate
\[
 k\,\bigl|\pi(n)A(n)-nF(n)\bigr|\le nF(n)
 \quad\text{for }k\ge1,\ n\ge N(k),
\]
where $F(n)=(n^2)!$, $A(n)=\sum_{j=n+1}^{n^2}F(n)/j$, and $N$ is a concrete primitive recursive term assembled from its imported stage functions~\cite{proveit-pnt}. In particular $A(n)/F(n)=H_{n^2}-H_n$ is an exact rational harmonic proxy for $\log n$. The source supplies a proof term using other library modules, rather than merely naming a goal. Reusing that endpoint with its complete checked dependency closure is a sensible next integration step.

This report does not reproduce that long proof or define all its imported stage functions, so it does not count the prime number theorem among the self-contained results proved here. Likewise the existence of a source file named after a major theorem would not be enough to count its full classical formulation. This separation prevents the source inventory from enlarging the mathematical claims beyond the arguments actually supplied.

\clearpage
\appendix
\section{Result inventory}
The following table is generated from the actual theorem environments and labels in the source. It includes lemmas and corollaries, and keeps the five infrastructure lemmas separate by their \texttt{fd:} prefix. The JSON inventory also records statements and proof locations; the topic manifests supply additional witnesses, bounds, and dependency notes.
\begin{longtable}{p{0.16\linewidth}p{0.75\linewidth}}
\toprule Result & Statement title\\
\midrule\endhead
Lemma~\ref{fd:reflection} & Bounded Boolean reflection\\
Lemma~\ref{fd:search} & Bounded search, finite sums, and finite products\\
Lemma~\ref{fd:loops} & Clocked loops and simultaneous invariants\\
Lemma~\ref{fd:descent} & Finite descent\\
Lemma~\ref{fd:php} & Finite pigeonhole principle with witnesses\\
Theorem~\ref{el:division} & Division with a bounded algorithm\\
Theorem~\ref{el:bezout} & Euclidean algorithm and B\'ezout identity\\
Theorem~\ref{el:linear} & Linear Diophantine equations\\
Lemma~\ref{el:leastprime} & A least divisor is prime\\
Theorem~\ref{el:euclidlemma} & Euclid's lemma\\
Theorem~\ref{el:fta} & Fundamental theorem of arithmetic\\
Theorem~\ref{el:valuation} & Valuations, gcd, and lcm\\
Theorem~\ref{el:euclid} & Euclid's theorem with a bound\\
Theorem~\ref{el:sieve} & Correctness of the sieve of Eratosthenes\\
Theorem~\ref{el:linearcong} & Solving and counting linear congruences\\
Theorem~\ref{el:crt2} & Generalized Chinese remainder theorem for two moduli\\
Theorem~\ref{el:crt} & Chinese remainder theorem for a finite system\\
Theorem~\ref{el:euler} & Euler's theorem\\
Corollary~\ref{el:fermat} & Fermat's little theorem\\
Theorem~\ref{el:order} & Order of a residue\\
Theorem~\ref{el:wilson} & Wilson's theorem and converse\\
Theorem~\ref{el:rsa} & RSA correctness, including nonunits\\
Theorem~\ref{el:digits} & Base-expansion divisibility rules\\
Theorem~\ref{el:binomial} & Prime binomial divisibility and the Frobenius identity\\
Theorem~\ref{el:legendre} & Legendre's factorial valuation formula\\
Theorem~\ref{el:lucas} & Lucas's binomial congruence\\
Theorem~\ref{el:kummer} & Kummer's carry formula\\
Theorem~\ref{el:divisorfunc} & Divisor formulas and multiplicativity\\
Theorem~\ref{el:squareparity} & Squarefree decomposition and parity of divisor functions\\
Theorem~\ref{el:totient} & Totient formula and multiplicativity\\
Theorem~\ref{el:phisum} & The divisor sum of the totient\\
Theorem~\ref{el:mobiuscancel} & M\"obius cancellation\\
Theorem~\ref{el:liouville} & Liouville's function and square detection\\
Theorem~\ref{el:convolution} & Finite Dirichlet convolution algebra\\
Theorem~\ref{el:mobius} & M\"obius inversion\\
Theorem~\ref{el:summatory} & Divisor summatory and hyperbola identities\\
Theorem~\ref{el:perfect} & Euclid--Euler characterization of even perfect numbers\\
Theorem~\ref{qr:roots} & Polynomial root bound\\
Lemma~\ref{qr:orders} & Orders and coprime products\\
Theorem~\ref{qr:primitive} & Primitive roots modulo a prime\\
Corollary~\ref{qr:primitive-count} & Orders of powers and the number of primitive roots\\
Lemma~\ref{qr:squares-half} & The square residues\\
Theorem~\ref{qr:euler-criterion} & Euler's criterion\\
Corollary~\ref{qr:character} & Multiplicativity\\
Theorem~\ref{qr:gauss} & Gauss's lemma\\
Lemma~\ref{qr:eisenstein} & The floor-sum form\\
Theorem~\ref{qr:reciprocity} & Quadratic reciprocity\\
Theorem~\ref{qr:supplements} & The two supplementary laws\\
Lemma~\ref{qr:two-product} & The two-square product identity\\
Lemma~\ref{qr:two-descent} & A strict, bounded two-square descent\\
Theorem~\ref{qr:two-prime} & Fermat's two-square theorem for primes\\
Lemma~\ref{qr:two-obstruction} & Primes congruent to three modulo four\\
Theorem~\ref{qr:two-classification} & The sum-of-two-squares classification\\
Lemma~\ref{qr:four-product} & Euler's four-square identity\\
Lemma~\ref{qr:four-seed} & A small represented multiple of an odd prime\\
Lemma~\ref{qr:four-descent} & A strict four-square descent, including the even case\\
Theorem~\ref{qr:four-squares} & Lagrange's four-square theorem\\
Corollary~\ref{qr:rational-four} & Four squares over the nonnegative rationals\\
Proposition~\ref{qr:three-obstruction} & The three-square obstruction and optimality\\
Theorem~\ref{dio:normal} & Canonical rational representation\\
Theorem~\ref{dio:rational-root} & Rational root test\\
Corollary~\ref{dio:powers} & Rational powers and irrational roots\\
Theorem~\ref{dio:pythagoras} & Primitive Pythagorean triples\\
Corollary~\ref{dio:all-triples} & All integer right triangles\\
Theorem~\ref{dio:square-area} & No integer right triangle has positive square area\\
Theorem~\ref{dio:fermat-four} & Fermat's quartic descent\\
Theorem~\ref{dio:cf-rational} & Finite continued fractions for rationals\\
Theorem~\ref{dio:continuants} & Continuant identity, coprimality, and growth\\
Theorem~\ref{dio:cf-error} & Exact tail and error identities\\
Corollary~\ref{dio:cf-real} & A primitive recursive continued-fraction real\\
Theorem~\ref{dio:farey} & Farey separation and the mediant\\
Theorem~\ref{dio:dirichlet} & Rational Dirichlet approximation\\
Theorem~\ref{dio:simultaneous} & Simultaneous rational Dirichlet approximation\\
Theorem~\ref{dio:real-transfer} & Certified transfer to represented reals\\
Lemma~\ref{dio:pell-approx} & Bounded norm with a prescribed lower denominator\\
Theorem~\ref{dio:pell-existence} & Uniform bounded Pell existence\\
Theorem~\ref{dio:pell-powers} & Pell composition and arbitrarily many solutions\\
Theorem~\ref{dio:pell-all} & All positive Pell solutions are fundamental powers\\
Corollary~\ref{dio:pell-family} & An explicit Pell family\\
Theorem~\ref{an:progressions} & Primes in the classes $1$ and $3$ modulo $4$\\
Theorem~\ref{an:sieve} & Finite inclusion--exclusion sieve\\
Corollary~\ref{an:sievepi} & Exact prime counting from a finite sieve\\
Lemma~\ref{an:binomial} & A central binomial coefficient and its prime powers\\
Theorem~\ref{an:chebyshev} & Dyadic Chebyshev bounds without logarithms\\
Lemma~\ref{an:harmonic} & Dyadic harmonic growth\\
Lemma~\ref{an:eulerdom} & Finite Euler domination at exponent one\\
Lemma~\ref{an:blockproduct} & Bounding a finite product by a reciprocal sum\\
Theorem~\ref{an:primeharmonic} & Divergence of prime reciprocals with an explicit bound\\
Lemma~\ref{an:zetatail} & A rational tail estimate\\
Theorem~\ref{an:eulererror} & Finite Euler-product error\\
Corollary~\ref{an:zetaname} & An iPRA name for $\zeta(s)$ and its Euler product\\
Theorem~\ref{an:means} & Exact means of prime-divisor counts\\
Theorem~\ref{an:independence} & Exact finite independence of distinct prime divisibility\\
Theorem~\ref{an:coprimepairs} & Coprime pairs and M\"obius inversion\\
Theorem~\ref{ag:norm} & Norm and integrality in a quadratic ring\\
Theorem~\ref{ag:nonufd} & Two inequivalent irreducible factorizations\\
Theorem~\ref{ag:valuation} & Truncated valuation laws\\
Lemma~\ref{ag:taylor} & Taylor congruence at a finite precision\\
Theorem~\ref{ag:henselstep} & Simple-root Hensel lifting with a unique digit\\
Theorem~\ref{ag:henselname} & A primitive recursive compatible Hensel name\\
Theorem~\ref{ag:partitionrec} & Partition recursion and finite generating product\\
Theorem~\ref{ag:oddpartitions} & Euler's odd-parts/distinct-parts identity\\
Theorem~\ref{ag:conjugate} & Conjugation of partitions\\
Theorem~\ref{tr:polynomial} & Finite polynomial estimates and a rational Liouville bound\\
Theorem~\ref{tr:name} & The Liouville name and its approximation bounds\\
Theorem~\ref{tr:transcendence} & Uniform polynomial separation and transcendence\\
\bottomrule
\end{longtable}

\clearpage
\section{Validation, review, and reproducibility}
\label{sec:validation}

\subsection{What the evidence establishes}
The report's mathematical evidence is the full proof of every numbered statement, including the explicit induction invariants and bounds that reduce the argument to \ipra. Independent mathematical cross-review was used on the quadratic, Diophantine, foundational, analytic, algebraic, and transcendence arguments. The review found and corrected a circular presentation of bounded Boolean reflection and search: the final proof first verifies the Boolean scan directly by quantifier-free induction, then derives bounded existential reflection. Smaller clarifications addressed zero in prime-divisor sums and capped valuations.

The computational evidence is a separate layer. The Python programs use exact integers and \texttt{fractions.Fraction}; they compare algorithms against independent finite specifications and check identities and invariants on stated ranges. Passing these checks does not prove a universally quantified mathematical theorem, establish a statement about all primitive recursive programs, or provide an object-language \ipra\ proof certificate. Very large symbolic bounds, such as the Pell worst-case box, the prime-reciprocal threshold, and the Liouville separation denominator, were not fully enumerated in the checks.

\subsection{Recorded finite checks}
\begin{longtable}{>{\raggedright\arraybackslash}p{0.24\linewidth}>{\raggedright\arraybackslash}p{0.67\linewidth}}
\toprule
Check family & Actual scope and purpose\\
\midrule\endhead
Elementary algorithms & Division on 1,377 cases; 2,116 extended Euclidean inputs with coefficient and descent invariants; 6,405 generalized CRT systems compared with brute force; 20,301 sieve entry comparisons.\\
Binomial and divisor identities & 19,880 Lucas cases, 17,892 Kummer cases, arithmetic identities for positive inputs through 400, 1,200 M\"obius inversion cases, and every even integer through 10,000 for the perfect-number characterization.\\
Quadratic and square representations & 1,980 ordered reciprocity prime pairs; primitive-root counts below 100; prime and composite square-representation algorithms through the stated bounds; 93,728 signed assignments for the even and odd four-square descent branches.\\
Diophantine algorithms & 16,545 cases across continued fractions, rational and simultaneous approximation, bounded norm approximants, Pell collisions and power descent, and Pythagorean parameter recovery.\\
Pell norm collision & Every nonsquare $D$ from 2 through 100, including collisions with both positive and negative common norm. For $D=61$ the recorded construction returns $(1766319049,226153980)$ using common norm $-9$.\\
Finite analytic estimates & Exact rational sieve identities, 12 dyadic Chebyshev cases, 150 Euler-product errors, 450 zeta tails, 750 truncation estimates, 1,500 denominator-removal inequalities, and complete-period CRT event patterns.\\
Algebraic and $p$-adic examples & Norm multiplication, capped valuation edge cases, Hensel simple-root digit uniqueness and compatible residues, and explicit partition bijections. Precise cases are in the companion JSON record.\\
Transcendence & Finite factorial-tail and polynomial-value bounds, root avoidance including a polynomial vanishing at the first two truncations, and the domination inequalities for the deliberately huge separation denominator.\\
\bottomrule
\end{longtable}

The JSON files beside the programs or in \texttt{verification/} contain the detailed ranges, counts, examples, and declared limitations. The recorded Pell example is an illustration of the explicit collision mechanism, not evidence that the worst-case bound is close to optimal. The programs are small reference implementations, not claims about best computational complexity.

\subsection{Reproducing the package}
From the package directory, run \texttt{make check} to execute the finite checks and regenerate the inventory, and \texttt{make pdf} to compile the \LaTeX\ report. The checks require Python 3 and only its standard library. The PDF build requires \texttt{latexmk}, pdf\LaTeX, and the standard packages listed in the preamble. Intermediate \LaTeX\ files are excluded from the delivered archive.

The standalone Lean directory pins the inspected Lean and Mathlib revisions. After installing that toolchain and obtaining dependencies, its proposed first replay is \texttt{lake build}. This replay is an instruction for future validation, not a command reported as successful here. \texttt{ProveItTargets.lean} is a separate integration file requiring the pinned ProveIt checkout and its selected source paths; it is excluded from the standalone semantic-definition target. The accompanying blueprints explain which proofs still need to be implemented.

The PDF was compiled and visually reviewed after the final text changes. Its result inventory is generated from the same theorem labels as the source, and the package includes a SHA-256 manifest for the delivered files. The archive is self-contained for reading the proofs, recompiling the PDF, and rerunning the Python checks. Completing the Lean object derivations requires the additional formalization work explicitly described in Section~\ref{sec:lean}.

\clearpage
\begin{thebibliography}{99}
\addcontentsline{toc}{section}{References and inspected sources}
\bibitem{wiki-number} Wikipedia contributors, \emph{Number theory}, consulted 22 November 2026. Used as the requested scope index. \url{https://en.wikipedia.org/wiki/Number_theory}.
\bibitem{wiki-prime} Wikipedia contributors, \emph{Prime number}. Scope reference. \url{https://en.wikipedia.org/wiki/Prime_number}.
\bibitem{wiki-dio} Wikipedia contributors, \emph{Diophantine approximation}. Scope reference. \url{https://en.wikipedia.org/wiki/Diophantine_approximation}.
\bibitem{wiki-padic} Wikipedia contributors, \emph{$p$-adic number}. Scope reference. \url{https://en.wikipedia.org/wiki/P-adic_number}.
\bibitem{mit781} Abhinav Kumar and Joseph Lee, \emph{Theory of Numbers}, MIT 18.781, Spring 2012, lecture-note collection. Primary teaching reference for the conventional topics. \url{https://ocw.mit.edu/courses/18-781-theory-of-numbers-spring-2012/pages/lecture-notes/}.
\bibitem{kumar-qr} Abhinav Kumar and Joseph Lee, \emph{Quadratic Residues, Quadratic Reciprocity}, MIT 18.781, Lecture 9 notes, 2012. \url{https://ocw.mit.edu/courses/18-781-theory-of-numbers-spring-2012/resources/mit18_781s12_lec9/}.
\bibitem{lynn-qr} Ben Lynn, \emph{Number Theory: Quadratic Reciprocity}, Stanford-hosted notes. \url{https://crypto.stanford.edu/pbc/notes/numbertheory/quadrecip.html}.
\bibitem{conrad-descent} Keith Conrad, \emph{Proofs by Descent}, University of Connecticut-hosted mathematical notes. \url{https://kconrad.math.uconn.edu/blurbs/ugradnumthy/descent.pdf}.
\bibitem{miller-liouville} Steven J. Miller and Ramin Takloo-Bighash, \emph{An Invitation to Modern Number Theory: Countable, Uncountable, Algebraic and Transcendentals}, draft dated 23 October 2003, Section 2.3.1. \url{https://web.williams.edu/Mathematics/sjmiller/public_html/BrownClasses/1/CountableAlgTran.pdf}.
\bibitem{proveit-readme} Vladimir Reshetnikov, \emph{ProveIt: Primitive recursive arithmetic}, repository overview at commit \texttt{bbb85238dac533fc4b7b84a750cca864fe35fed5}. \url{https://github.com/VladimirReshetnikov/ProveIt/blob/bbb85238dac533fc4b7b84a750cca864fe35fed5/Logic/PrimitiveRecursiveArithmetic/README.md}.
\bibitem{proveit-pr} \emph{ProveIt: ConciseIPRA.PR}, definition of \texttt{PR}, \texttt{PTm}, \texttt{PRAAx}, and quantifier-free induction at the same commit. \url{https://github.com/VladimirReshetnikov/ProveIt/blob/bbb85238dac533fc4b7b84a750cca864fe35fed5/Logic/PrimitiveRecursiveArithmetic/ConciseIPRA/Lean/ConciseIPRA/PR.lean}.
\bibitem{proveit-bounded} \emph{ProveIt: Foundations.BoundedSyntax}, bounded syntax, compiler, and reflection derivations at the same commit. \url{https://github.com/VladimirReshetnikov/ProveIt/blob/bbb85238dac533fc4b7b84a750cca864fe35fed5/Logic/PrimitiveRecursiveArithmetic/Derivations/Lean/IPRADerivations/Foundations/BoundedSyntax.lean}.
\bibitem{proveit-pnt} \emph{ProveIt: AnalyticNumberTheory.FinalPNT}, arithmetic prime number theorem source endpoint, at the same commit. The complete transitive proof closure was not rebuilt for this report. \url{https://github.com/VladimirReshetnikov/ProveIt/blob/bbb85238dac533fc4b7b84a750cca864fe35fed5/Logic/PrimitiveRecursiveArithmetic/Derivations/Lean/IPRADerivations/AnalyticNumberTheory/FinalPNT.lean}.
\end{thebibliography}
\end{document}
